\documentclass[12pt]{book}

% ─── Fonts ────────────────────────────────────────────────────────────────────
\usepackage[sc]{mathpazo}          % Palatino for text and maths (elegant serif)
\usepackage[T1]{fontenc}
\linespread{1.06}                  % Palatino benefits from slightly open leading

% ─── Core packages ────────────────────────────────────────────────────────────
\usepackage{microtype}             % Microtypography: better justification
\usepackage{amsmath, amssymb}
\usepackage{graphicx}
\usepackage{float}
\usepackage{wrapfig}
\usepackage{listings}
\lstset{basicstyle=\ttfamily\footnotesize, breaklines=true}
\usepackage{amsthm}
\usepackage{tikz}
\usetikzlibrary{arrows.meta, positioning, shapes.geometric, calc}
\usepackage{emptypage}             % No headers/footers on intentionally blank pages
\usepackage{etoolbox}

% ─── Colour definitions ───────────────────────────────────────────────────────
\usepackage[dvipsnames]{xcolor}
\definecolor{accent}{RGB}{0, 62, 125}       % Deep UNB navy
\definecolor{boxbg}{RGB}{237, 243, 252}     % Soft blue tint for boxes

% ─── Page geometry ────────────────────────────────────────────────────────────
\usepackage[
  top    = 1.1in,
  bottom = 1.1in,
  inner  = 1.15in,
  outer  = 0.9in,
  headheight = 15pt
]{geometry}
\setlength{\parindent}{0pt}
\setlength{\parskip}{5pt}

% ─── Hyperlinks ───────────────────────────────────────────────────────────────
\usepackage{hyperref}
\hypersetup{
  colorlinks = true,
  linkcolor  = accent,
  urlcolor   = accent,
  citecolor  = accent,
  pdftitle   = {Random Variables and Stochastic Processes},
  pdfauthor  = {Kevin Englehart, PhD, PEng},
}

% ─── Chapter & section headings ───────────────────────────────────────────────
\usepackage{titlesec}

% Chapter: "CHAPTER  1" label above, large title, thick colour rule beneath
\titleformat{\chapter}[display]
  {\normalfont\bfseries}
  {\normalfont\small\scshape\color{accent}\MakeLowercase{\chaptertitlename}\ \thechapter}
  {4pt}
  {\Huge}
  [\vspace{6pt}{\color{accent}\rule{\linewidth}{2pt}}]
\titlespacing*{\chapter}{0pt}{-10pt}{32pt}

% Starred chapter (Preface, etc.)
\titleformat{name=\chapter,numberless}[display]
  {\normalfont\bfseries}
  {}
  {0pt}
  {\Huge}
  [\vspace{6pt}{\color{accent}\rule{\linewidth}{2pt}}]

\titleformat{\section}
  {\normalfont\Large\bfseries}
  {\color{accent}\thesection}
  {0.8em}{}[\vspace{-2pt}]

\titleformat{\subsection}
  {\normalfont\large\bfseries}
  {\color{accent}\thesubsection}
  {0.8em}{}

\titleformat{\subsubsection}
  {\normalfont\normalsize\bfseries\itshape}
  {\color{accent}\thesubsubsection}
  {0.8em}{}

% ─── Running headers & footers ────────────────────────────────────────────────
\usepackage{fancyhdr}
\pagestyle{fancy}
\fancyhf{}
\fancyhead[LE]{\small\color{accent!70!black}\nouppercase{\leftmark}}
\fancyhead[RO]{\small\color{accent!70!black}\nouppercase{\rightmark}}
\fancyfoot[LE,RO]{\small\color{accent}\thepage}
\renewcommand{\headrulewidth}{0.4pt}
\renewcommand{\footrulewidth}{0pt}

% Plain pages (chapter openers, TOC)
\fancypagestyle{plain}{%
  \fancyhf{}%
  \fancyfoot[C]{\small\color{accent}\thepage}%
  \renewcommand{\headrulewidth}{0pt}%
}

% ─── Table of contents ────────────────────────────────────────────────────────
\usepackage{tocloft}
\renewcommand{\cftchapfont}{\bfseries\color{accent}}
\renewcommand{\cftchappagefont}{\bfseries\color{accent}}
\renewcommand{\cftchapleader}{\color{accent!60}\cftdotfill{\cftdotsep}}
\renewcommand{\cftsecleader}{\color{accent!40}\cftdotfill{\cftdotsep}}
\setlength{\cftbeforechapskip}{8pt}
\setlength{\cftsecindent}{1.6em}

% ─── Boxes & callouts ─────────────────────────────────────────────────────────
\usepackage[breakable, skins]{tcolorbox}

% Default style applied to every plain \begin{tcolorbox}
\tcbset{
  every tcolorbox/.style={
    enhanced,
    breakable,
    colback   = boxbg,
    colframe  = accent,
    boxrule   = 0.8pt,
    arc       = 4pt,
    left      = 8pt,
    right     = 8pt,
    top       = 6pt,
    bottom    = 6pt,
    fontupper = \small,
  }
}

% Aside / Proof box (named environment, distinct lighter border)
\newtcolorbox{asidebox}[1]{
  enhanced,
  breakable,
  colback     = boxbg,
  colframe    = accent!45,
  boxrule     = 0.7pt,
  arc         = 3pt,
  title       = {#1},
  fonttitle   = \bfseries\small,
  coltitle    = accent,
  attach boxed title to top left = {yshift=-2mm, xshift=5mm},
  boxed title style = {
    colback   = white,
    colframe  = accent!45,
    boxrule   = 0.7pt,
    arc       = 2pt,
  },
  left        = 8pt,
  right       = 8pt,
  top         = 10pt,
  bottom      = 6pt,
}

% ─── Theorem / example style ──────────────────────────────────────────────────
\theoremstyle{definition}
\newtheorem{exmp}{Example}[chapter]

% ─── Misc commands ────────────────────────────────────────────────────────────
\newcommand{\F}{\mathcal{F}}
\newcommand*\rfrac[2]{{}^{#1}\!/_{#2}}
\graphicspath{{figures/}}

% ─── Custom cover page ────────────────────────────────────────────────────────
\newcommand{\makecover}{%
  \begin{titlepage}
  \thispagestyle{empty}
  \begin{tikzpicture}[remember picture, overlay]
    % Full-page dark navy banner at top (tall enough for all title text)
    \fill[accent] (current page.north west) rectangle
      ([yshift=-10.2cm] current page.north east);
    % Thin accent rule below banner
    \fill[accent!50!white] ([yshift=-10.2cm] current page.north west)
      rectangle ([yshift=-10.45cm] current page.north east);
    % Bottom band
    \fill[accent!12!white] (current page.south west) rectangle
      ([yshift=3.5cm] current page.south east);
    \fill[accent] ([yshift=3.5cm] current page.south west)
      rectangle ([yshift=3.3cm] current page.south east);
  \end{tikzpicture}
  \vspace*{1.5cm}
  {\centering
    % Title block — all white against the navy band
    \color{white}
    {\fontsize{34}{42}\selectfont\bfseries Random Variables\par}
    \vspace{0.5cm}
    {\fontsize{34}{42}\selectfont\bfseries and Stochastic\par}
    \vspace{0.5cm}
    {\fontsize{34}{42}\selectfont\bfseries Processes\par}
    \vspace{0.7cm}
    {\normalsize\color{white!75!accent}\itshape Course Notes for EE~6513\par}
    \vspace{4.8cm}
    \color{accent}
    {\small\scshape authored by\par}
    \vspace{0.35cm}
    {\Large\bfseries Kevin Englehart, PhD, PEng\par}
    \vspace{0.25cm}
    {\small University of New Brunswick \textbullet\ Fredericton, NB\par}
    \vspace{0.4cm}
    {\small\color{accent!70!black} March 2026\par}
  }
  \end{titlepage}
}

% Colophon / title verso page
\newcommand{\makecolophon}{%
  \thispagestyle{empty}
  \vspace*{\fill}
  \begin{center}
    \small
    \textit{Random Variables and Stochastic Processes}\\[4pt]
    Kevin Englehart, PhD, PEng\\
    University of New Brunswick\\
    Faculty of Engineering\\
    Fredericton, NB, Canada\\[12pt]
    Course Notes for EE~6513\\[12pt]
    \textcopyright\ 2026 Kevin Englehart, PhD, PEng. \\ This work is licensed under a Creative Commons Attribution-NonCommercial-NoDerivatives 4.0 International License.\\[4pt]
  \end{center}
  \vspace*{\fill}
  \newpage
}

\title{Random Variables and Stochastic Processes}
\author{Kevin Englehart, PhD, PEng}
\date{March 2026}

\begin{document}

\makecover
\makecolophon

\frontmatter
\tableofcontents
\clearpage

\mainmatter

%\chapter{First Chapter}
%\section{Section name}
%\lipsum[1]
%\setcounter{chapter}{1}
%\setcounter{section}{1}
%\section{Section name}
%\lipsum[1-2]

\chapter*{Preface}

Welcome to \textit{An Introduction to Probability, Random Variables, and Stochastic Processes}, a comprehensive guide designed to provide a thorough understanding of the fundamental concepts and methods used in the analysis of random phenomena. This textbook is meticulously structured to cater to students, educators, and professionals who seek a solid foundation in probability theory, random variables, and stochastic processes, which are pivotal in various fields such as engineering, finance, and data science.

\section*{Scope and Objectives}

The primary objective of this textbook is to build a robust conceptual framework that allows readers to confidently analyze and interpret random processes and events. We aim to bridge the gap between theory and practical application, enabling you to apply probabilistic methods to solve real-world problems. This book covers a wide range of topics, beginning with the basics of probability theory and advancing to more complex concepts like stochastic processes and optimal filtering.

\section*{Structure of the Book}

The content is organized into seven detailed chapters, each focusing on a specific area of study:

\begin{description}
    \item[1. Probability] This chapter lays the groundwork with an introduction to probability theory, exploring its axioms, conditional probability, and the concept of repeated trials. Topics such as Bayes' Rule and Bernoulli Trials are examined to provide a comprehensive understanding of probability fundamentals.
    
    \item[2. Random Variables] Delving into the heart of probabilistic analysis, this chapter discusses probability distributions and density functions. It covers properties of discrete and continuous random variables, the Central Limit Theorem, and functions of random variables, alongside expected values and moments.
    
    \item[3. Stochastic (Random) Processes] Expanding on the concept of random variables, this chapter introduces stochastic processes, discussing their first and second-order properties, stationary processes, and the power density spectrum. Key concepts such as ergodicity and autocovariance functions are explained in detail.
    
    \item[4. Estimation] This chapter addresses the methods and performance of different estimators. It covers Bayesian estimation, maximum likelihood estimation, and minimum mean square estimation, providing essential techniques for parameter estimation in statistical models.
    
    \item[5. Detection] Focusing on binary detection, this chapter explores various detection strategies, including the Bayesian binary detector and binary detection in Gaussian white noise. Performance measures of detectors are also discussed, offering insight into their practical applications.
    
    \item[6. Optimal Filtering] A critical aspect of signal processing, this chapter examines linear time-invariant systems with stochastic inputs. Topics such as Wiener filters and the Kalman filter are covered extensively, providing the tools needed to implement optimal filtering techniques.
    
    \item[7. Analysis of Time-Series Data] The final chapter equips readers with the skills to analyze time-series data. It includes tests for self-stationarity, criteria for ergodicity, and estimation of power density spectra, enabling a comprehensive understanding of time-series analysis.
\end{description}

\section*{Learning Approach}

Each chapter begins with an introduction to the fundamental concepts, followed by detailed explanations and mathematical formulations. Examples and exercises are provided to reinforce the material and facilitate hands-on learning. By progressing through the chapters, readers will develop a deep and practical understanding of probabilistic and stochastic methods.

\section*{Who Should Read This Book}

This textbook is ideal for undergraduate and graduate students studying probability, statistics, and engineering. It is also a valuable resource for researchers and professionals who need to apply probabilistic models and stochastic processes in their work.



\chapter{Probability}

\section{Introduction}

Understanding the nature of uncertainty and randomness is fundamental to comprehending the behavior of random processes, which appear in a myriad of scientific and engineering disciplines. The study of random processes is a domain characterized by its intrinsic uncertainty, where we endeavor to model and analyze systems whose outcomes are subject to unpredictable variability. This chapter lays the foundation for our exploration of random processes by introducing the central concept of probability.\\

Probability serves as the cornerstone of the modern scientific framework that navigates us through uncertain and stochastic phenomena. It provides a framework for quantifying our level of belief about different outcomes, thereby enabling us to make informed decisions and perform meaningful analyses. \\

This chapter will navigate you through the fundamental principles of probability, drawing from a probabilistic paradigm. We will explore the rules of probability, Bayesian reasoning and probability distributions. By weaving together these complementary approaches, you will gain a comprehensive understanding of how probability theory underpins the study of random processes.

\section{Axioms of Probability Theory}

\begin{enumerate}
    \item \textbf{First Axiom:} To each event \(A\) of a set of all possible events \(S\), there is assigned a non-negative real number \(P(A)\) called the probability of that event:

\[
\quad P(A) \geq 0, \quad P(A) \in \mathbb{R}
\]

This might be conceptualized as the area of event $P({A})$ relative to the area of ${S}$ as illustrated in Figure \ref{fig:MutExcl} panel (a).

\begin{figure}[h]
\centering
\includegraphics{mutually_exclusive.png}
\caption{Probability space describing (a) probability of event \(A\) is the relative area of it's locus to the area of total probability, \(S\); (b) mutually exclusive events, and (c) events that are not mutually exclusive, \(P(A\cap B) \ne 0\).}
\label{fig:MutExcl}
\end{figure}

\item \textbf{Second Axiom:} The probability of a certain event $S$ (entire the entire sample space) is 
    \[
P(S) = 1
\]

\item \textbf{Third Axiom:} Two events, \(A\) and \(B\) are defined as \emph{mutually exclusive} when, if one of the events happens, the other cannot. Mathematically, two events are mutually exclusive if:


\[
P(A \cap B) = 0
\]

This is illustrated in Figure \ref{fig:MutExcl}, panel (b).


\begin{asidebox}{Aside: Terminology}
Some notes pertaining to terminology.
\begin{itemize}
    \item The term \(A \cap B\) reads `the intersection of \(A\) and \(B\).' This is equivalent to the term \(A \cdot B\).
    \item For mutually exclusive events \(A\) and \(B\), \(A \cap B = A \cdot B = \varnothing\), where \(\varnothing\) is the null (empty) set.
\end{itemize}
\end{asidebox}

\item \textbf{Fourth Axiom:} If \(A_1, A_2, A_3, \ldots\) are mutually exclusive events, then the probability of the union of these events is equal to the sum of their individual probabilities:

\[
P(A_1 \cup A_2 \cup A_3 \cup \ldots) = P(A_1) + P(A_2) + P(A_3) + \ldots
\]

or

\[
P\left(\bigcup_{i=1}^{n} A_i\right) = \sum_{i=1}^{n} P(A_i).
\]

The general result for the sum of probabilities of events \(A_i\) that are not mutually exclusive is given by the principle of inclusion-exclusion. It allows us to calculate the probability of the union of these events while accounting for the overlap between them.

For \(n\) events \(A_1, A_2, \ldots, A_n\), the probability of their union is calculated as:

\begin{multline*}
P(A_1 \cup A_2 \cup \cdots \cup A_n) = 
    \sum_{i=1}^n P(A_i) 
    - \sum_{i<j} P(A_i \cap A_j) \\
    + \sum_{i<j<k} P(A_i \cap A_j \cap A_k) 
    - \cdots 
    + (-1)^{n-1}P(A_1 \cap A_2 \cap \cdots \cap A_n)
\end{multline*}

In this formula, \(P(A_i)\) represents the probability of event \(A_i\), and \(P(A_i \cap A_j)\) represents the probability of the intersection of events \(A_i\) and \(A_j\). The principle of inclusion-exclusion accounts for the double-counting of overlapping events and subtracts it to find the correct probability of the union.

\end{enumerate}

\begin{tcolorbox}
\begin{exmp}{Mutually Exclusive Events: Throwing a Die.}


Consider an experiment of rolling a six-sided fair die. Let's define two events:

\begin{itemize}
    \item Event \(A\): Getting an even number (2, 4, or 6).
    \item Event \(B\): Getting an odd number (1, 3, or 5).
\end{itemize}


These events are defined in such a way that no outcome can belong to both \(A\) and \(B\). In other words, if event \(A\) occurs (e.g., rolling a 4), then event \(B\) cannot occur simultaneously (e.g., rolling a 3).\\

Mathematically, we can express this mutual exclusivity as:

\[
P(A \cap B) = 0
\]

Calculate the probabilities of each event.

\end{exmp}
\end{tcolorbox}


\[
P(A) = \frac{\text{Number of favorable outcomes for } A}{\text{Total number of outcomes}} = \frac{3}{6} = \frac{1}{2}
\]

\[
P(B) = \frac{\text{Number of favorable outcomes for } B}{\text{Total number of outcomes}} = \frac{3}{6} = \frac{1}{2}
\]\\

Notice that the probabilities of events \(A\) and \(B\) sum to 1, which is consistent with the principle together they comprise the complete set \(S\) and that they are mutually exclusive.\\






\begin{tcolorbox}
    \begin{exmp}{A Problem Involving Non-Mutually Exclusive Events.}\\
        In a classroom of 30 students, 15 students play soccer (Event A), 12 students play basketball (Event B), and 8 students play both soccer and basketball (Events A and B). Calculate the following probabilities:

\begin{enumerate}
    \item The probability that a randomly selected student plays soccer (Event A).
    \item The probability that a randomly selected student plays basketball (Event B).
    \item The probability that a randomly selected student plays either soccer or basketball.
\end{enumerate}

    \end{exmp}
\end{tcolorbox}


\textbf{Solution:}

Let's solve each part of the problem:

\begin{enumerate}
    \item The probability that a randomly selected student plays soccer (Event A) is given by the ratio of students playing soccer to the total number of students:

   \[P(A) = \frac{\text{Number of students playing soccer}}{\text{Total number of students}} = \frac{15}{30} = 0.5\]

   \item The probability that a randomly selected student plays basketball (Event B) is given by:

   \[P(B) = \frac{\text{Number of students playing basketball}}{\text{Total number of students}} = \frac{12}{30} = 0.4\]

    \item To find the probability that a student plays either soccer or basketball (not mutually exclusive), we can use the principle of inclusion and exclusion:

   \[P(A \cup B) = P(A) + P(B) - P(A \cap B)\]

   \[P(A \cup B) = 0.5 + 0.4 - P(A \cap B) \]
   
   We are given that 8 students play both soccer and basketball (Event A and Event B), so 
   \[P(A \cap B) = \frac{8}{30} = 0.267\]
   \[P(A \cup B) = 0.9 - 0.267 = 0.633\]
   
\end{enumerate}

Therefore, the probability that a randomly selected student plays either soccer or basketball is 0.633.


\begin{tcolorbox}
\begin{exmp}{Toss a Coin Three Times.}  

For a single coin toss the space $S$ consists of the outcomes heads and tails ($h$ and $t$):
\[ S = \{h,t\} \]

Now, for three consecutive coin tosses, the sample space $S$ includes all possible outcomes:
\[ S = \{hhh, hht, hth, thh, htt, tth, tht, ttt\} \]

Let's define event \(A\) as obtaining at least two heads in the outcomes, and event \(B\) as having the first two coin tosses result in heads. Calculate the probability of \(P(A)\), \(P(B)\) and \(P(A \cap B)\).

\end{exmp}
\end{tcolorbox}

Each of these outcomes is mutually exclusive of the others, e.g:

\[
P(hhh \cap hht) = 0
\]

Consequently,

\[
P(hhh \cup hht) = P(hhh) + P(hht) 
\]

\textbf{Probability Calculation:}
\begin{align*}
P(A) &= P(\text{at least two heads in the outcomes}) \\
& = P(hhh) + P(hht) + P(hth) + P(thh) \\
& = 4 \times \frac{1}{8} \\
& = \frac{1}{2}
\end{align*}


\textbf{Probability Calculation:}
\begin{align*}
P(B) &= P(\text{first two coin tosses result in heads}) \\
&= P(hhh) + P(hht) \quad \text{these events are mutually exclusive}\\
&= \frac{1}{8} + \frac{1}{8} \\
&= \frac{1}{4}
\end{align*}

Consider now \(P(A \cap B)\):
\[
P(A \cap B) = P(hhh) + P(hht) = P(B)
\]

Clearly, \(P(A \cap B) \ne 0\) and so events \(A\) and \(B\) are not mutually exclusive. \\

The union of \(A\) and \(B\) is:

\[
\begin{aligned}
    P(A \cup B) & = P(A) + P(B) - P(A \cap B)\\
        & = P(A) + P(B) - P(B) \\
        & = P(A) \\
        & = \frac{1}{2}
\end{aligned}
\]







\section{Conditional Probability}
The probability of an event \(A\) given that event \(B\) has occurred is defined as:

\[
P(A | B) = \frac{P(A \cap B)}{P(B)}
\]

\textbf{Cases:}
A depiction of these cases is provided in Figure \ref{fig:MECases}.
\begin{enumerate}
\item If \(A\) and \(B\) are mutually exclusive (disjoint):
\[ P(A | B) = 0 \]

\item If \(B \in A\):
\[ P(A | B) = \frac{P(B)}{P(B)} = 1 \]

\item If \(A \in B\):
\[ P(A | B) = \frac{P(A)}{P(B)} \]


\item \(P(A \cup B | C) = P(A | C) + P(B | C)\) if \(A\) and \(B\) are mutually exclusive.\\

\item Given the mutually exclusive events \(A_i\), \(i = 1, \ldots, n\), and \(S = \bigcup_{i=1}^{n} A_i\):
\[ P(B) = \sum_{i=1}^{n} P(B | A_i) \cdot P(A_i) \text{ for } B \in S \]

\end{enumerate}

\begin{figure}[!h]
\centering
\includegraphics[scale=0.6]{ME_Cases.png}
\caption{Conditional probability (a) the general case, (b) $A$ and $B$ mutually exclusive, (c) \(B \in A\), (d) \(A \in B\), (e) three events with $A$ and $B$ mutually exclusive and, (f) the general case for mutually exclusive events.}
\label{fig:MECases}
\end{figure}

\bigbreak

\begin{asidebox}{Proof of Case 5}
Let's start with the set \(S\) and express \(B\) in terms of the mutually exclusive events:
\[
B = B \cap S = (B \cap A_1) + (B \cap A_2) + \ldots + (B \cap A_n)
\]
Because the events \(A_i\), \(i = 1, \ldots, n\), are mutually exclusive, we can use the definition of conditional probability:
\[
P(B) = P(B \cap A_1) + P(B \cap A_2) + \ldots + P(B \cap A_n)
\]
Furthermore, we can express each term using conditional probability:
\[
P(B) = P(B | A_1) \cdot P(A_1) + P(B | A_2) \cdot P(A_2) + \ldots + P(B | A_n) \cdot P(A_n)
\]
This concludes the proof.
\end{asidebox}






\subsection{Bayes' Rule}

Bayes' Rule is an extension of conditional probability that allows us to update our beliefs based on new evidence. It is expressed as:

\begin{equation}
    P(A | B) = \frac{P(B | A) \cdot P(A)}{P(B)}
\end{equation}

Here's a breakdown of the components in Bayes' Rule:

\begin{itemize}
    \item $P(A | B)$: Posterior probability of $A$ given $B$.
    \item $P(B | A)$: Likelihood of $B$ given $A$.
    \item $P(A)$: Prior probability of $A$.
    \item $P(B)$: Total probability of $B$.
\end{itemize}

Bayes' Rule allows us to update our initial beliefs (prior probability) about the occurrence of event $A$ based on the new information provided by event $B$. It is a powerful tool in statistical inference, machine learning, and various scientific disciplines.


\begin{tcolorbox}
\begin{exmp}{Bayes' Rule Example: Medical Testing}

Consider a rare disease, Disease X, affecting 1 in 1000 people in a population. A diagnostic test for Disease X has a sensitivity (true positive rate) of 95\% and a false positive rate of 10\%.

Let's define the events:
\begin{itemize}
    \item $A$: The individual has Disease X.
    \item $B$: The test result is positive.
\end{itemize}

We want to find $P(A | B)$, the probability that the individual has the disease given a positive test result, using Bayes' Rule.

\end{exmp}
\end{tcolorbox}

The true positive rate is the likelihood of $B$ given $A$: 
\[
P(B|A) = 95\%
\]

The true negative rate is the likelihood of $B$ given $\bar{A}$: 
\[
P(B|\bar{A}) = 10\%
\]

The probability of the individual having disease X given a positive test result is given by Bayes' Rule:
\[
    P(A | B) = \frac{P(B | A) \cdot P(A)}{P(B)}
\]

The law of total probability is used to calculate $P(B)$:

\[
P(B) = P(B | A) \cdot P(A) + P(B | \bar A) \cdot P(\bar A)
\]

Substitute the given values:
\[
P(A | B) = \frac{0.95 \cdot \frac{1}{1000}}{0.95 \cdot \frac{1}{1000} + 0.10 \cdot \frac{999}{1000}}
\]

After computation, $P(A | B) \approx 0.0094$, which is approximately 0.94\%. Despite the positive test result, the probability of actually having the disease is relatively low.\\


\subsection{Multivariate Bayes' Rule}

Bayes's rule can be extended to include more than one possible prior event. 
Consider the mutually exclusive events $A_1, \ldots, A_n$, such that 
\[
S = \sum_{i = 1}^n A_i
\]
and another event \(B\) in this same space,  
as illustrated in Figure \ref{fig:Bayes}. \\


From the definition of conditional probability:

\[ 
\begin{aligned}
    P(A_i | B) & = \frac{P(A_i \cap B)}{P(B)} \\
    & = \frac{P(B \cup A_i)}{P(B)} \\
    & = \frac{P(B | A_i) P(A_i)}{P(B)} 
\end{aligned} \]


Bayes' theorem can be stated as follows:
\[ 
\boxed{P(A_i | B) = \frac{P(B | A_i)  P(A_i)}{\sum_{i=1}^nP(B|A_i)P(A_i)}.
}
\]


This theorem provides the probability that event $A_i$ might occur, given that event $B$ occurred. 

\begin{figure}[h]
\centering
\includegraphics[scale=0.40]{figures/Bayes.png}
\caption{For the mutually exclusive events $A_1, \ldots, A_n$, Bayes' theorem allows us to determine the probability that a given \(A_i\) may occur, given that event \(B\) has occurred. }
\label{fig:Bayes}
\end{figure}

\begin{tcolorbox}
\begin{exmp}{Multivariate Application of Bayes' Theorem.}

Suppose there are two factories, Factory A and Factory B, that produce a certain type of light bulb. Factory A produces 60\% of the bulbs, and Factory B produces the remaining 40\%. The bulbs from Factory A have a defect rate of 5\%, while the bulbs from Factory B have a defect rate of 3\%.\\

Determine the probability that the bulb came from Factory A given that it is defective.
\end{exmp}
\end{tcolorbox}

Let's define our events:
\begin{itemize}
    \item \(A\): The bulb comes from Factory A.
    \item \(B\): The bulb comes from Factory B.
    \item \(D\): The bulb is defective.
\end{itemize}

\bigskip

We have the following \emph{a priori} probabilities:
\begin{itemize}
    \item \(P(A) = 0.60\) (probability of a randomly selected bulb coming from Factory A).
    \item \(P(B) = 0.40\) (probability of a randomly selected bulb coming from Factory B).
\end{itemize}

\bigskip 

We also have the following conditional probabilities:
\begin{itemize}
    \item \(P(D | A) = 0.05\) (probability of a bulb being defective given that it came from Factory A).
    \item \(P(D | B) = 0.03\) (probability of a bulb being defective given that it came from Factory B).
\end{itemize}

Now, we want to find \(P(A | D)\), the probability that the bulb came from Factory A given that it is defective, using Bayes' theorem:

\[
P(A | D) = \frac{P(D | A) \cdot P(A)}{P(D | A) \cdot P(A) + \cdot P(D | B) \cdot P(B)}
\]

Substituting the values:

\[
P(A | D) = \frac{0.05 \cdot 0.60}{0.05 \cdot 0.06 + 0.03 \cdot 0.40} \approx 0.7143
\]
So, the probability that the bulb came from Factory A given that it is defective is approximately 0.7143 or 71.43\%.

\bigskip

Further treatments of Baye's rule are provided in sections on Bayesian Estimation (\ref{sec:Bayesian_Estimation}) and Binary Detection (\ref{sec:Binary_Detection}).

\subsection{Mutually Exclusive and Independent Events}

In the context of stochastic processes, it's essential to understand the concepts of independence and mutual exclusivity when dealing with events. These concepts play a significant role in probability theory.


\subsubsection{Mutually Exclusive Events}


Two events, \(A\) and \(B\), are considered mutually exclusive if they cannot occur simultaneously. See Figure \ref{fig:MutExcl} panel (b). In other words, if one of the events happens, the other must not. Mathematically, two events are mutually exclusive if:

\[
P(A \cap B) = 0
\]

Additionally, the conditional probability of \(A\) given \(B\) and the conditional probability of \(B\) given \(A\) are both zero:

\[
\begin{aligned}
    P(A | B) & = P(A \cap B) = 0 \\
    P(B | A) & = P(A \cap B) = 0    
\end{aligned}
\]


\subsubsection{Independent Events}

Two events, \(A\) and \(B\), are considered independent if the occurrence of one event does not affect the probability of the other event occurring. Mathematically, two events are independent if:

\[
P(A \cap B) = P(A) \cdot P(B)
\]

Additionally, the conditional probability of \(A\) given \(B\) and the conditional probability of \(B\) given \(A\) are equal to the individual probabilities:

\[
\begin{aligned}
    P(A | B) & = P(A) \\
    P(B | A) & = P(B)
\end{aligned}
\]




\subsubsection{Relationship between Mutually Exclusive and Independent Events}

It's important to note that if two events, \(A\) and \(B\), are mutually exclusive, then they are not independent. This relationship can be expressed as:

\[
P(A \cap B) = 0 \quad \text{and} \quad P(A \cap B) \neq P(A) \cdot P(B)
\]

Conversely, if two events are independent, they are not mutually exclusive. Is important that this distinction is recognized, as these two terms are often misused.



\section{Repeated Trials}
Consider events \(A\) in experiment \(\zeta_1\) and event \(B\) in experiment \(\zeta_2\). The experiments are independent if and only if:

\[
P(A \cap B) = P(A) \cdot P(B)
\]

for all \(A \in \zeta_1\) and \(B \in \zeta_2\).

\subsection{Bernoulli Trials}
Bernoulli trials are a sequence of random experiments or trials in which each trial has only two possible outcomes: success and failure. These trials are characterized by the following key properties:

\begin{itemize}
    \item \textbf{Two Outcomes:} In a Bernoulli trial, there are exactly two possible outcomes, which are often labeled as 'success` (often denoted as 1) and 'failure` (often denoted as 0). These outcomes are mutually exclusive, meaning that if one occurs, the other cannot.
    \item \textbf{Independent Trials:} Each trial is independent of the others. The outcome of one trial does not affect the outcome of any other trial. This property ensures that the trials are conducted under the same conditions and that past results do not influence future outcomes.

    \item \textbf{Constant Probability:} The probability of success (usually denoted as \(p\)) remains the same for each trial. Likewise, the probability of failure is \(1 - p\). This means that the likelihood of success or failure is consistent across all trials.

    \item \textbf{Discrete Random Variable:} The outcome of a Bernoulli trial can be described as a discrete random variable. It takes on the values 0 (failure) or 1 (success) with associated probabilities.

    \item \textbf{Fixed Number of Trials:} Bernoulli trials are often described within the context of a fixed number of trials, such as \(n\) trials. In this case, the random variable representing the number of successes in \(n\) trials follows a \emph{binomial} distribution.
\end{itemize}

The binomial distribution is characterized by having exactly $k$ outcomes of event $A$ in $n$ independent trials. Its probability mass function is:
\[ P_n(k) = \binom{n}{k}  p^k (1 - p)^{n - k}. \]


Common examples of Bernoulli trials include:
\begin{itemize}
    \item Flipping a fair coin: The outcome of each coin flip is a Bernoulli trial. "Heads" may be considered a success (1), and "tails" a failure (0).
    \item Testing the reliability of a product: Each product is either reliable (success) or unreliable (failure) when tested.
    \item Conducting a medical test: Each individual either tests positive (success) or negative (failure) for a specific condition.
\end{itemize}


Bernoulli trials serve as the building blocks for more complex probability models, such as the binomial distribution (which models the number of successes in a fixed number of trials) and other related probability distributions. They are fundamental in probability theory and statistics, and many real-world applications involve situations that can be modeled as Bernoulli trials.

\newpage


\begin{tcolorbox}
\begin{exmp}{Bernoulli Trials: Roll of Two Dice.}

Two dice are rolled 10 times. We want to find the probability that a sum of 7 is obtained at least once.

\end{exmp}
\end{tcolorbox}


\begin{itemize}
    \item Let \(A\) be the event that a sum of 7 is obtained on a single roll of the two dice. 
    \item To determine \(P(A)\), we can calculate the probability of obtaining a sum of 7. There are a total of 36 possible outcomes when rolling two dice (each die has 6 faces). Out of these, there are 6 outcomes that result in a sum of 7: (1, 6), (2, 5), (3, 4), (4, 3), (5, 2), and (6, 1).
\end{itemize}

So, \(P(A)\), the probability of obtaining a sum of 7 on a single roll, is:

\[
P(A) = \frac{6}{36} = \frac{1}{6}
\]

The probability of not obtaining a sum of 7 on a single roll is:

\[
P(\bar A) = 1 - P(A) = 1 - \frac{1}{6} = \frac{5}{6}
\]


Now, we want to find the probability of getting zero sums of 7 in 10 rolls of two dice. Using the binomial formula, this probability is:

\[
P_{10}(0) = \binom{10}{0} \left(\frac{1}{6}\right)^0 \left(\frac{5}{6}\right)^{10} = \left(\frac{5}{6}\right)^{10}
\approx 0.1615\]

Now, we can find the probability that a sum of 7 is obtained at least once by taking the complement of the event that it is not obtained on any roll:
\[
P_{10}(k\ge1) = 1 - P_{10}(0) = 1 - \left(\frac{5}{6}\right) ^{10}
 \approx 0.8385\]



\begin{tcolorbox}
\begin{exmp}{Events in a time series.}


Suppose we have a long line segment (such as time) of length \(T\), and we want to randomly place \(n\) points on this interval. What is the probability that exactly \(k\) of these points fall within the subinterval \((0, \tau)\), where \(\tau\) is a specific value within the interval \((0, T)\)?\\

\end{exmp}
\end{tcolorbox}

This problem can be approached as a problem in repeated trials, where the experiment is the placement of a single point within the interval \((0, T)\). Let's define the following events:

\begin{itemize}
    \item \(A\): The point is placed within the subinterval \((0, \tau)\).
    \item \(A'\): The point is placed outside the subinterval \((0, \tau)\).
\end{itemize}


We want to calculate \(P(k)\), the probability that exactly \(k\) out of the \(n\) points fall within \((0, \tau)\). In each trial, the probability of success \(P(A)\) is 

\[P(A) = p = \frac{\tau}{T}\]

and the probability of failure \(P(A')\) is 

\[P(A') = 1 - p = \frac{T - \tau}{T}\]

Now, we can use the binomial probability formula to find \(P(k)\):

\[
P(k) = \binom{n}{k} p^k (1 - p)^{n - k} = \binom{n}{k} \left(\frac{\tau}{T}\right)^k \left(\frac{T - \tau}{T}\right)^{n - k}
\]

This formula gives us the probability that exactly \(k\) points out of the \(n\) placed within the interval \((0, T)\) fall within the subinterval \((0, \tau)\).


\begin{tcolorbox}
\begin{exmp}{Example of Events in a Series: Train Arrival Times.}

Imagine a railway station where trains arrive at regular intervals throughout the day. The station operates from 6:00 AM to 10:00 PM, a total of 16 hours (\( T = 16 \) hours). Throughout these 16 hours, a total of 40 trains (\( n = 40 \)) are scheduled to arrive. Assume that the likelihood of arrival is equally spread throughout this time.\\

Now, let's say we're interested in a specific time window, for instance, from 8:00 AM to 10:00 AM (\( \tau = 2 \) hours). We want to know the probability that exactly 5 out of the 40 trains fall within this specific time window.

\end{exmp}
\end{tcolorbox}



Using the same approach as before:

\[
P(k) = \binom{n}{k} p^k (1 - p)^{n - k} = \binom{n}{k} \left(\frac{\tau}{T}\right)^k \left(\frac{T - \tau}{T}\right)^{n - k}
\]

\[ 
\begin{aligned}
    P(k=5) & = \binom{40}{5} \left(\frac{2}{16}\right)^5 \left(\frac{16 - 2}{16}\right)^{35}\\
    &  = \left( \frac{40!}{5!35!} \right ) \left(\frac{1}{8}\right)^{5} \left(\frac{14}{16}\right)^{35} \\
    & \approx 0.1875 
\end{aligned}
    \]

So, the probability that exactly 5 out of the 40 trains scheduled to arrive at the railway station fall within the time window from 8:00 AM to 10:00 AM is approximately 0.1875 or 18.75\%. \\

It is instructive to visualize this result with respect to the corresponding binomial distribution. For \(n=40\) trains arriving in \(T=16\) hours, and a 2-hour duration of interest, the probability of a `successful' event is  \(p = \frac{2}{16}=\frac{1}{8}\). The corresponding binomial distribution is illustrated in Figure \ref{fig:binom_trains}.

\begin{figure}[h]
\centering
\includegraphics[scale=0.60]{figures/binom_trains.png}
\caption{The binomial distribution with \(n=40, p=1/8\). }
\label{fig:binom_trains}
\end{figure}

Since the average rate of arrival is 40/16 = 2.5 trains per hour, it is intuitive that, in the 2-hour time window of interest, the case of \(k=5\) trains arriving would have the greatest probability. 



\begin{tcolorbox}
\begin{exmp}{Component failure.}


Given:
\begin{itemize}
    \item Number of components, $n = 1000$
    \item Probability of failure for each component, $p = 0.1$
    \item Number of failures we want, $k = 12$
\end{itemize}


We want to find the probability that exactly 12 out of 1000 components will fail between (0, $\tau$).

\end{exmp}
\end{tcolorbox}

Using the binomial probability formula:

\[
P_n(k) = \binom{n}{k} p^k (1 - p)^{n - k}
\]

Substituting the values:

\[
P(12) = \binom{1000}{12} \cdot (0.1)^{12} \cdot (1 - 0.1)^{1000 - 12} \approx 1.2092^{-30}
\]

Note, this can be difficult to compute, as numerical precision may be insufficient.



\begin{tcolorbox}
\begin{exmp}{Probability of Getting Heads 5000 Times.}


Given:

\begin{itemize}
    \item Number of coin tosses, $n = 10000$
    \item Probability of getting heads (event A), $P(A) = \frac{1}{2}$
    \item Number of times we want to get heads, $k = 5000$
\end{itemize}


We want to find the probability that heads will show up 5000 times when the coin is tossed 10,000 times. 

\end{exmp}
\end{tcolorbox}

Using the binomial probability formula:

\[
P(k) = \binom{n}{k} p^k (1 - p)^{n - k}
\]

Substituting the values:

\[
P_{10000}(5000) = \binom{10000}{5000} \left(\frac{1}{2}\right)^{5000} \left(1 - \frac{1}{2}\right)^{5000}
\]

\[
P_{10000}(5000) = \frac{10000!}{2 \cdot 5000!} = \left(\frac{1}{2}\right)^{10000} \approx 0.008
\]

This result can be evaluated numerically built-in scripts in a computing platform such as Matlab or Python. For large \(n\) and \(p\), this is typically done using approximation formulas to avoid numerical precision errors. \\


One such method is a normal approximation to the binomial distribution which is valid for a mean $np \geq 10$ and variance $n(1-p) \geq 10$. Here, $\mu=np=5000$ and variance $\sigma^2=np(1-p)=2500$. Evaluating the corresponding normal distribution pdf:

\[
f(x | \mu, \sigma^2) = \frac{1}{\sqrt{2 \pi \sigma^2}} \exp\left(-\frac{(x - \mu)^2}{2 \sigma^2}\right)
\]

for \(x=k=5000\), yields 0.008. 



\chapter{Random Variables}

\section{Introduction}
A random variable is a real-valued function that maps the outcomes of a random experiment to numerical values. This mathematical construct serves as a bridge between the abstract world of probability and the concrete world of data analysis. Random variables come in two principal forms: \textbf{discrete} and \textbf{continuous}.\\

Discrete random variables are associated with countable sets of outcomes of an experiment. These outcomes, often denoted as \(A_i\), are mapped by a random variable \(X\) to a finite set of real numbers \(x_i\):

\[
X(A_i) = x_i
\]

Consider a coin toss in which the outcomes are: \(A_1 = heads\) and \(A_2 = tails\). We may define a random variable \(X\) that operates on the outcomes, assigning values as 

\[
    \begin{aligned}
         X(A_1) & = 0 \\ 
        X(A_2) & = 1
    \end{aligned}
\]
Note that the index of the events \(Ai\) in general does not equal the value assigned by the random variable \(X\).\\

For a continuous random variable, the outcomes \(A\) of an experiment are mapped to elements of the real line, \(\mathbb{R}\), by a random variable \(X\):

\[
X(A) \rightarrow \mathbb{R}
\]


\section{Probability Distribution and Density Functions}

The properties of random variables are characterized by their Probability Density Functions (PDFs), Probability Mass Functions (PMFs), and Cumulative Density Functions (CDFs).

\begin{itemize}
    \item PDFs describe the likelihood of continuous random variables taking specific values or falling within ranges.
    \item PMFs fulfill a similar role for discrete random variables, assigning probabilities to individual outcomes.
    \item CDFs offer a cumulative perspective on probability distributions, summarizing the probability that a random variable is below or equal to a given value.
\end{itemize}

The concepts of the PDF, PMF and CDF are illustrated in \ref{fig:pdf_pmf_cdf}.

\begin{figure}[h!]
\centering
\includegraphics[scale=0.5]{figures/pdf_pmf_cdf.png}
\caption{Illustrative example of the PDF, PMF and CDF. For a continuous random variable $X$ (e.g. normal distribution), the probability is distributed over a continuous range (top left). The corresponding CDF (top right) shows the probability of $X$ being less than or equal to a given value. The PMF of a discrete random variable is shown (bottom left) for a binomial distribution This depicts the probability that $X$ assigns a specific value, $k$.  The corresponding corresponding CDF is shown (bottom right), exhibiting incremental cumulative probability at each value of $k$.}
\label{fig:pdf_pmf_cdf}
\end{figure}

\subsection{Properties of PDFs, PMFs and CDFs}
Let's first consider continuous random variables. The CDF provides insights into the cumulative probability associated with a random variable \(X\). In simple terms, it answers the question, "What is the probability that \(X\) takes on a value less than or equal to \(x\)?"

\[F_\mathrm{x}(x) = P(X \leq x)\]

Some general properties of cumulative distribution functions are:
\begin{enumerate}
    \item \(F_\mathrm{x}(\infty) = 1\) \\
        - all values assigned by \(X\) must be less than \(\infty\), so the probability is 1. 
    \item \(F_\mathrm{x}(-\infty) = 0\) \\
        - no values of \(X\) are less than \(-\infty\), so the probability is 0. 
    \item \(F_\mathrm{x}(x_1) \leq F_\mathrm{x}(x_2)\), for \(x_1 \leq x_2\) \\
        - since \(F_\mathrm{x}(x)\) is cumulative. 
\end{enumerate}

The PDF, denoted as \(f_\mathrm{x}(x)\), focuses on quantifying the probability that the random variable \(X\) assumes a particular value \(x\), or how the probability changes as \(x\) varies. The PDF can be computed as the derivative of the CDF

\[f_\mathrm{x}(x) = \frac{dF_\mathrm{x}(x)}{dx}\]

or equivalently, the CDF can be viewed as the integral of the PDF:

\[F_\mathrm{x}(x) = P(X \leq x)\ = \int_{-\infty}^{x} f_\mathrm{x}(x) \, dx \]

The PDF and the CDF can also be used to express the likelihood of \(X\) falling within an interval \((x_1,x_2)\).

\[P(x_1 \leq X \leq x_2) = \int_\mathrm{x_1}^{x_2} f_\mathrm{x}(x) \, dx = F_\mathrm{x}(x_2) - F_\mathrm{x}(x_1)\]

This is illustrated in Figure \ref{fig:prob_calc}.

\begin{figure}[h!]
\centering
\includegraphics[scale=0.5]{figures/prob_calc.png}
\caption{Illustrative example of how the PDF and the CDF can also be used compute the probability of \(X\) falling within an interval $(x_1, x_2)$. This is the area under the PDF between $x_1$ and $x_2$ and the difference in height of the CDF between $x_1$ and $x_2$ .}
\label{fig:prob_calc}
\end{figure}


\subsection{Discrete Random Variables}

Some properties of discrete random variables differ from continuous random variables.
\begin{itemize}
    \item A finite number of values are assigned by the RV.

    \item The CDF \(F_\mathrm{x}(x)\) becomes a staircase function. The height of the "stair" at value \(x_i\) equals \(P(X=x_i)\). 

    \item The PMF \(f_\mathrm{x}(x)\) becomes a series of impulses: 

\[
f_\mathrm{x}(x) = \sum_{i=1}^{n} p_i  \delta(x - x_i)
\]

where:
\begin{itemize}
    \item \(n\) is the number of values.
    \item \(p_i\) is the probability at \(P(X=x_i)\).
    \item \(\delta(x - x_i)\) is the Dirac delta function centered at \(x_i\).
\end{itemize}


\end{itemize}

This is illustrated in Figure \ref{fig:pmp_discrete}.


\begin{figure}[h!]
\centering
\includegraphics[scale=0.5]{figures/pmp_discrete.png}
\caption{Illustrative example of how the PMF and the CDF can also be used compute the probability of \(X\) falling within an interval $(x_1, x_2)$. This is the area under the PMF (equally, the sum of the height of the impulses) or the difference in height of the CDF. Here, for values of \(X\) between 2 and 4 (inclusive), the probability \(P(X=\{2,3,4\}) = 0.2+0.4+0.2 = 0.8\)}
\label{fig:pmp_discrete}
\end{figure}



\begin{tcolorbox}
\begin{exmp}{Experiment - Die throw.}


\begin{itemize}
    \item Events: \(A_i\), where \(i = 1, 2, \ldots, 6\).
    \item Random variable \(X(A_i) = i\).
    \item Probability of each event: \(P(A_i) = \frac{1}{6}\) (assuming a fair die). 
\end{itemize}

Show the probability distribution function and the probability density function of the random variable \(X\).

\end{exmp}
\end{tcolorbox}

The PMF and CDF of the fair die throw are shown in Figure \ref{fig:Fx_die}.

\begin{figure}[h!]
\centering
\includegraphics[scale=0.3]{figures/f_x_die_throw.png}
\includegraphics[scale=0.3]{figures/Fx_die_throw.png}
\caption{The PMF (left) and CDF (right) of a fair die throw.}
\label{fig:Fx_die}
\end{figure}


\subsection{Example Distributions}

\bigskip
\subsubsection{Poisson (discrete)}
The Poisson probability distribution is a statistical distribution used to model the number of events occurring in a fixed interval of time or space when the events happen independently at a constant rate. It is commonly used in various fields such as insurance, telecommunications, biology, and economics to model the occurrence of events like arrivals of customers in a queue, accidents in a given time frame, radioactive decay events, etc. Here's a breakdown of its key characteristics:

\begin{itemize}
    \item \textbf{Constant rate:} The events occur at a constant average rate, denoted by the symbol \(\lambda\), which represents the average number of events in the given interval.
    \item \textbf{Discrete outcomes:} The possible number of events in a given interval is discrete, meaning it can only take on integer values (0, 1, 2, 3, ...)
    \item \textbf{Events occur independently:} Each event is independent of the others, meaning that the occurrence of one event does not affect the probability of another event occurring.
\end{itemize}


The probability mass function \(P(X = k)\) gives the probability that the random variable \(X\) takes on a specific value \(k\). It describes the likelihood of observing exactly \(k\) events in the specified interval.
\[
P(X = k) = \frac{e^{-\lambda} \lambda^k}{k!}, \quad k = 0, 1, 2, \ldots
\]

where \(k\) represents the number of events, \(\lambda\) is the average rate of events, and \(e\) is the base of the natural logarithm.

\[
f_\mathrm{x}(x) = \begin{cases} 
\frac{e^{-\lambda} \lambda^x}{x!} & \text{for } x = 0, 1, 2, \ldots \\
0 & \text{otherwise}
\end{cases}
\]

or 

\[
f_\mathrm{x}(x) = \sum_{i=0}^{\infty}\frac{e^{-\lambda} \lambda^x}{x!} \delta(x-i)
\]

This is illustrated in \ref{fig:Possion} for \(\lambda=2\) and \(\lambda=4\).

\begin{figure}[h!]
\centering
\includegraphics[scale=0.5]{figures/Possion.png}
\caption{PMFs of Poisson RVs with \(\lambda=2\) and \(\lambda=4\).}
\label{fig:Possion}
\end{figure}

Some important properties are as follows.

\begin{enumerate}

  \item \textbf{Mean and Variance:} The mean (\(\mu\)) and variance (\(\sigma^2\)) of the Poisson distribution are both equal to \(\lambda\): \(\mu = \sigma^2 = \lambda\).

  \item \textbf{Shape:} The distribution is right-skewed for small \(\lambda\) and becomes more symmetric as \(\lambda\) increases.

\end{enumerate}


\subsubsection{Binomial Distribution (discrete)}

Consider an experiment with two outcomes, \(A\) and \(\bar A\), each with probabilities \(p\) and \(q\), respectively. If we repeat this experiment \(n\) times, let \(k\) be the number of outcomes A. The random variable \(X\) is a discrete random variable with the following probability mass function (PMF):

\[
P(X = k) = \binom{n}{k} p^k q^{n-k}.
\]


The probability density function of \(X\) is given by:

\[
f_\mathrm{x}(x) = \sum_{k=0}^n\binom{n}{k} p^k q^{n-k} \delta(x-k)
\]


Some important properties include: 
\begin{enumerate}

    \item \textbf{Mean and Variance:} The mean and variance of the binomial distribution are:
    \[
    \mu = np \quad \text{and} \quad \sigma^2 = np(1-p)
    \]

    \item \textbf{Shape:} The shape of the binomial distribution depends on \(n\) and \(p\). It can be skewed or bell-shaped.

\end{enumerate}

\begin{tcolorbox}
\begin{exmp}{Coin Toss.}

Let's consider a coin toss experiment with \(p = q = \frac{1}{2}\) (i.e., a fair coin) and \(n\) total tosses. The PMF of \(X\) for this experiment is:

\[
f_\mathrm{x}(x) = \begin{cases}
\binom{n}{x} \left(\frac{1}{2}\right)^x \left(\frac{1}{2}\right)^{n-x}, & \text{for } x = 0, 1, 2, \ldots, n \\
0, & \text{otherwise}
\end{cases}
\]

If \(p\) is defined as an outcome of \emph{heads}, determine the PMF of \(X=k\) heads for \(n=5\) coin tosses.
\end{exmp}
\end{tcolorbox}

For example, if \(n = 5\), the PDF is:

\[
f_\mathrm{x}(x) = \begin{cases}
\frac{1}{32}, & \text{for } x = 0 \\
\frac{5}{32}, & \text{for } x = 1 \\
\frac{10}{32}, & \text{for } x = 2 \\
\frac{10}{32}, & \text{for } x = 3 \\
\frac{5}{32}, & \text{for } x = 4 \\
\frac{1}{32}, & \text{for } x = 5 \\
0, & \text{otherwise}
\end{cases}
\]

This PMF is illustrated in Figure \ref{fig:binomial}.

\begin{figure}[h!]
\centering
\includegraphics[scale=0.5]{figures/binomial.png}
\caption{PMF of binomial a RV with \(n=5\).}
\label{fig:binomial}
\end{figure}


\subsubsection{Gaussian (Normal) Distribution (continuous)}
The Gaussian distribution, denoted as \(\mathcal{N}(\mu, \sigma^2)\), is described by the probability density function:
    \[
    f_\mathrm{x}(x) = \frac{1}{\sigma \sqrt{2\pi}} e^{-\frac{(x - \mu)^2}{2\sigma^2}}
    \]

where  \(\mu\) is the mean and \(\sigma\) is the standard deviation. The variance is the square of the standard deviation, \(\sigma^2\). This distribution is denoted as \(\mathcal{N}(\mu, \sigma)\). This is illustrated in \ref{fig:Gaussian} for \(\mathcal{N}(1, 1)\) and \(\mathcal{N}(0, 3)\).\\


\begin{figure}[h!]
\centering
\includegraphics[scale=0.4]{figures/Gaussian.png}
\caption{Probability density function of a Gaussian distribution with different values of \(\mu\) and \(\sigma\).}

\label{fig:Gaussian}
\end{figure}

\begin{itemize}   
    
    \item \textbf{Standard Normal Distribution:} The standard normal distribution has \(\mu = 0\) and \(\sigma = 1\). It is denoted as \(\mathcal{N}(0, 1)\).



    \item \textbf{Properties:}
    \begin{itemize}
        \item Symmetric and bell-shaped.
        \item Mean, median, and mode are equal (\(\mu\)).
        \item Empirical Rule: About 68\% of data falls within one standard deviation, 95\% within two, and 99.7\% within three.
    \end{itemize}
\end{itemize}

The normal distribution is appropriate for many applications due to several factors:

\begin{enumerate}
    \item \textbf{Natural Phenomena:} Many natural processes and phenomena exhibit behaviors that approximate a normal distribution. For instance, human characteristics like height, weight, and IQ tend to follow a normal distribution due to the combined effect of numerous genetic and environmental factors. Additionally, measurement errors and random fluctuations often conform to a normal distribution due to the central limit theorem.
    \item \textbf{The Central Limit Theorem:} The CLT states that the sum (or average) of a large number of independent and identically distributed random variables tends to have an approximately normal distribution, regardless of the original distribution of the variables. This property makes the normal distribution a suitable model for the aggregate behavior of many random processes.
    \item \textbf{Mathematical Simplicity:} The normal distribution has simple and well-defined mathematical properties, making it tractable for analytical and computational purposes. Its probability density function and cumulative distribution function have closed-form expressions, facilitating calculations and statistical analysis.
    \item \textbf{Statistical Methods and Inference:} Many statistical methods and inference procedures are based on assumptions of normality. For instance, hypothesis testing, confidence interval estimation, and linear regression often assume that the data are normally distributed. When these assumptions are met, these methods tend to be robust and efficient.
    
\end{enumerate}





\subsection{Jointly Distributed Random Variables}

Consider two random variables, \(X\) and \(Y\). The joint distribution function is defined as:

\[
F_\mathrm{xy}(x, y) = P(X \leq x, Y \leq y)
\]

And the joint density function is given by:

\[
f_\mathrm{xy}(x, y) = \frac{\partial^2 F_\mathrm{xy}(x, y)}{\partial x \partial y}
\]

To find probabilities involving jointly distributed random variables \(X\) and \(Y\), you can use the following integral expressions:

\[
P(a \leq X \leq c, b \leq Y \leq d) = \int_a^c \int_b^d f_\mathrm{xy}(x, y) \,  dx \, dy
\]

\[
P(X \le x) = \int_{-\infty}^x \int_{-\infty}^{\infty} f_\mathrm{xy}(x, y) \, dx  \,dy  = F_\mathrm{x}(X)
\]

\bigskip

\textbf{Note:} The integral of the joint density functions are the \emph{marginal} density functions. 
\[
f_\mathrm{x}(X) = \int_{-\infty}^\infty f_\mathrm{xy}(x, y)  \, dy
\]

\[
f_\mathrm{y}(Y) = \int_{-\infty}^\infty f_\mathrm{xy}(x, y) \, dx
\]

\subsection{Jointly Gaussian Random Variables}


Assuming that $X$ and $Y$ are jointly Gaussian, their joint PDF is given by a bivariate normal distribution:

\[
f_\mathrm{xy}(x, y) = \frac{1}{2\pi\sigma_\mathrm{x}\sigma_\mathrm{y}\sqrt{1-\rho^2}} \exp\left(-\frac{1}{2(1-\rho^2)}\left[\frac{(x-\mu_\mathrm{x})^2}{\sigma_\mathrm{x}^2} - 2\rho\frac{(x-\mu_\mathrm{x})(y-\mu_\mathrm{y})}{\sigma_\mathrm{x}\sigma_\mathrm{y}} + \frac{(y-\mu_\mathrm{y})^2}{\sigma_\mathrm{y}^2}\right]\right)
\]

Where:
\begin{itemize}
    \item $\mu_\mathrm{x}$ and $\mu_\mathrm{y}$ are the means (expected values) of $X$ and $Y$, respectively.
    \item $\sigma_\mathrm{x}^2$ and $\sigma_\mathrm{y}^2$ are the variances of $X$ and $Y$, respectively.
    \item $\sigma_\mathrm{xy}$ is the covariance between $X$ and $Y$.
    \item $\rho$ is the correlation coefficient between $X$ and $Y$, computed as 
    \[
    \rho = \frac{\sigma_\mathrm{xy}}{\sigma_\mathrm{x} \sigma_\mathrm{y}}
    \]
\end{itemize}

If $\rho = 0$, $X$ and $Y$ are uncorrelated, and the PDF simplifies to the product of two independent univariate normal distributions. If $\rho \neq 0$, the PDF shows how $X$ and $Y$ are correlated.\\


\begin{tcolorbox}
\begin{exmp}{Jointly Gaussian RV.}\\
Let \(X\) and \(Y\) be jointly Gaussian random variables with means $\mu_\mathrm{x} = 3$ and $\mu_\mathrm{y} = 2$, variances $\sigma_\mathrm{x}^2 = 4$ and $\sigma_\mathrm{y}^2 = 1$, and a covariance of $\sigma_\mathrm{xy} = 1.5$. \\

Calculate the probability that \(X\) takes a value between 1 and 5 (inclusive) and \(Y\) takes a value between 1 and 3 (inclusive).\\
\end{exmp}

\end{tcolorbox}


First:
   \[
    \rho = \frac{\sigma_\mathrm{xy}}{\sigma_\mathrm{x} \sigma_\mathrm{y}} = \frac{1.5}{(2)(1)} = 0.75
    \]


Substituting the values of mean, variance and correlation:

\begin{multline*}
f_\mathrm{xy}(x, y) = \frac{1}{2\pi(2)(1)\sqrt{1-(0.75)^2}} \\
    \times \exp\left(-\frac{1}{2(1-(0.75)^2)}
    \left[\frac{(x-3)^2}{4} 
    - \frac{2(0.75)(x-3)(y-2)}{(2)(1)} 
    + \frac{(y-2)^2}{1}\right]\right)
\end{multline*}

Now, the probability that \(X\) takes a value between 1 and 5 (inclusive) and \(Y\) takes a value between 1 and 3 (inclusive) is: 
\begin{equation}
P(1 \leq X \leq 5, 1 \leq Y \leq 3) = \int_{1}^{5} \int_{1}^{3} f_\mathrm{xy}(x, y) \, dy \, dx \approx 0.683.
\end{equation}

This is evaluated in numerical form in the code \href{https://www.ece.unb.ca/kengleha/Textbook/Code/Joint_Gauss.py}{\texttt {Joint\_Gauss.py}}, which also draws the corresponding PDF (Figure \ref{fig:joint-Gaussian}).

\begin{figure}[h!]
\centering
\includegraphics[scale=0.8]{figures/joint-Gauss.png}
\caption{The area under the jointly Gaussian PDF that yields \( P(1 \leq X \leq 5, 1 \leq Y \leq 3) \) .}

\label{fig:joint-Gaussian}
\end{figure}





\section{Central Limit Theorem (CLT)}

The Central Limit Theorem states that the distribution of the sum (or average) of a large number of independent, identically distributed random variables approaches a Gaussian distribution as the sample size (\(N\)) becomes large, regardless of the original distribution.\\
    
\textbf{Mathematical Representation:} If \(X_1, X_2, \ldots, X_N\) are independent and identically distributed random variables with mean \(\mu\) and variance \(\sigma^2\), then the distribution of the random variable \(Z = \sum_{i=1}^{N} X_i \) approaches a  normal distribution as \(N\) becomes large. A proof of the CLT will be explored in Section \ref{sec:moments}\\


As an example of the manifestation of the CLT in physical signals, consider electromagnetic interference (EMI) that might be encountered in a measurement. This interference can arise from various sources, both natural and man-made, and exhibit different probability distribution functions. For example:

\begin{itemize}
    \item \emph{Mobile Phones (Man-Made Source)}: Mobile phones emit electromagnetic radiation due to their communication signals. The distribution of signal strength (in dBm) from mobile phones can be modeled as a normal distribution. Each phone call or data transmission contributes to the overall EMI in the environment.
    \item \emph{Lightning Strikes (Natural Source)}: Lightning generates intense electromagnetic fields during a strike. The occurrence of lightning strikes follows a Poisson distribution. The more frequent the strikes, the higher the EMI levels. These sudden discharges can affect nearby electronic systems.
    \item \emph{Solar Flares (Natural Source)}: Solar flares release bursts of high-energy particles and electromagnetic radiation. The intensity of solar flares follows a power-law distribution. During strong solar events, EMI can disrupt satellite communications and power grids.
    \item \emph{Microwave Ovens (Man-Made Source)}: Microwave ovens emit microwave radiation during cooking. The distribution of microwave power levels can be approximated by a gamma distribution. Leakage from faulty ovens can contribute to EMI in the vicinity.
    \item \emph{Ignition Systems (Man-Made Source)}: Ignition systems in vehicles generate sparks that emit electromagnetic radiation. The timing of ignition events follows a uniform distribution. These sparks can interfere with nearby radio receivers.
\end{itemize}


Typical PDFs of these EMI sources are depicted in Figure \ref{fig:EMI}. Assuming that the total EMI measured is the sum of contributions from each source, the PDF of the total EMI is depicted on the bottom row, right.

\begin{figure}[h!]
\centering
\includegraphics[scale=0.4]{figures/CLT.png}
\caption{The CLT illustrated by a mixture of independent electromagnetic signals.}

\label{fig:EMI}
\end{figure}


This is a simplified illustration, but it demonstrates how different EMI sources combine to create a composite interference signal that is roughly normal. This is observed in many physical systems. 


\subsection{Conditional Distributions}

\[
F_\mathrm{x}(x|A) = P(X \le x | A) = \frac{P(X \le x, A)}{P(A)}
\]

\[
f_\mathrm{x}(x|A) = \frac{d F_\mathrm{x}(x|A)}{dx}
\]


\textbf{Note:}
\begin{align*}
F_\mathrm{x}(\infty|A) &= 1 \\
F_\mathrm{x}(-\infty|A) &= 0
\end{align*}



\begin{tcolorbox}
    \begin{exmp}
        Let the event \(A\) be defined by the random variable \(Y\), i.e., \(A = \{Y \leq y\}\). The conditional probability function \(F_\mathrm{x}(x|A)\) is given by 
   
       
\[
 \begin{aligned}
        F_\mathrm{x}(x|A) & = F_\mathrm{x}(x|Y \leq y)\\
        & = P(X \le x| Y \le y) \\
        & = \frac{P(X \le x , Y \le y)}{P(Y \le y)}
    \end{aligned}
\]
    \end{exmp}
\end{tcolorbox}
    


\begin{tcolorbox}
    \begin{exmp}
        {Example}
    Let the event \(A\) be defined by the random variable \(X\), i.e., \(A = \{X \leq a\}\). The conditional probability function \(F_\mathrm{x}(x|A)\) is given by 
   
       
\[
 \begin{aligned}
        F_\mathrm{x}(x|A) & = P(X \le x| X \le a) \\
        & = \frac{P(X \le x , X \le a)}{P(X \le a)}
    \end{aligned}
\]

Two cases:
\begin{enumerate}
    \item If \(x \ge a\), then 
    \[F_\mathrm{x}(x|A) = \frac{P(X \le a)}{P(X \le a)} = 1\]
     \item If \(x < a\), then 
    \[F_\mathrm{x}(x|A) = \frac{P(X \le x)}{P(X \le a)}\]
\end{enumerate}
    \end{exmp}
\end{tcolorbox}



\section{Functions of a Random Variable}

Given a random variable \(X\), we can form a new random variable \(Y = g(X)\) by applying a function \(g\) to \(X\). It's important to understand how to determine the probability distribution of \(Y\). There are two approaches to achieve this, which we'll label the \emph{direct} approach and the \emph{derivative} approach.

\subsection{Direct Approach}

In the direct approach, we substitute the function \(g\) into the probability distribution of \(X\). For example, if \(Y = g(X) = aX + b\), where \(a\) and \(b\) are constants, we can find the probability distribution of \(Y\) by direct substitution into \(F_\mathrm{y}(y)\).\\


First, recall the CDF of \(Y\), which is \(F_\mathrm{y}(y) = P(Y \leq y)\). Using the transformation \(Y = aX + b\), we substitute for \(Y\):

\[
F_\mathrm{y}(y) = P(aX + b \leq y)
\]

Assuming that we know \(F_\mathrm{x}(x)\), to find \(F_\mathrm{y}(y)\) the first step is determine \(F_\mathrm{x}(y)\). To do this, we must express \(x\) in terms of \(y\) using the inverse transformation:

\[
x = \frac{y - b}{a}
\]

Substitute this expression for \(x\) back into the CDF of \(Y\):

\[
F_\mathrm{y}(y) = P\left(Y \leq y\right) = P\left(X \leq x\right) = P\left (X \leq \frac{y - b}{a} \right)
\]

So that:

\[
F_\mathrm{y}(y) = F_\mathrm{x}\left(\frac{y - b}{a}\right)
\]

Now, 

\[
f_\mathrm{y}(y) = \frac{d}{dy}F_\mathrm{y}(y)
\]

Now, differentiate \(F_\mathrm{y}(y)\) with respect to \(y\):

\[
f_\mathrm{y}(y) = \frac{d}{dy}\left(F_\mathrm{x}\left(\frac{y - b}{a}\right)\right)
\]

Using the chain rule:

\[
f_\mathrm{y}(y) = F_\mathrm{x}\left(\frac{y - b}{a}\right) \cdot \frac{d}{dy}\left(\frac{y - b}{a}\right)
\]

Since \(X\) is a random variable with a probability density function \(f_\mathrm{x}(x)\), we can write:


\[
f_\mathrm{y}(y) = \frac{1}{a} \cdot f_\mathrm{x}\left(\frac{y - b}{a}\right)
\]

\begin{tcolorbox}
    \begin{exmp}{Find \(f_\mathrm{y}(y)\) for \(Y=x^2\).}
    
         Given \(F_\mathrm{x}(x)\), the CDF of \(X\), we want to find the CDF and PDF of \(Y\).
    \end{exmp}
\end{tcolorbox}


\begin{figure}[h!]
\centering
\includegraphics[scale=0.4]{figures/y_x_squared.png}
\caption{Direct mapping of \(y\) onto \(x\) axis, for the function \(Y = X^2\)..}
\label{fig:y=x}
\end{figure}

First, we start with the definition of \(F_\mathrm{y}(y)\):

\[
F_\mathrm{y}(y) = P(Y \leq y)
\]

Since \(Y = X^2\), we can rewrite this as:

\[
F_\mathrm{y}(y) = P(X^2 \leq y)
\]

Now, we express this in terms of the original random variable \(X\):

\[
F_\mathrm{y}(y) = P(-\sqrt{y} \leq X \leq \sqrt{y})
\]

Using the CDF of \(X\), \(F_\mathrm{x}(x)\), we can write:

\[
F_\mathrm{y}(y) = F_\mathrm{x}(\sqrt{y}) - F_\mathrm{x}(-\sqrt{y})
\]

Now,

\[
f_\mathrm{y}(y) = \frac{d}{dy}F_\mathrm{y}(y)
\]

Using the chain rule, we can write this as:

\[    \begin{split}
f_\mathrm{y}(y)  & = \frac{d}{dy}\left(F_\mathrm{x}(\sqrt{y}) - F_\mathrm{x}(-\sqrt{y})\right) \\
        & = f_\mathrm{x}(\sqrt{y}) \frac{d\sqrt{y}}{dy} - f_\mathrm{x}(-\sqrt{y}) \frac{d (\sqrt{y})}{dy} \\
        & = \frac{1}{2 \sqrt{y}} x(\sqrt{y}) - \frac{-1}{2 \sqrt{y}} x(- \sqrt{y})
    \end{split}
\]
\[    
f_\mathrm{y}(y) = \frac{1}{2 \sqrt{y}} \left (f_\mathrm{x}(\sqrt{y}) + f_\mathrm{x}(-\sqrt{y}) \right ) u(y) 
\]\\

Here, \(u(y) \) denotes that this expression for \(f_\mathrm{y}(y) \) is defined only for \(y \ge 0\). \\

Now, consider resolving \(f_\mathrm{y}(y)\) for a specific \(f_\mathrm{x}(x)\). If we let 
    \[ 
    f_\mathrm{x}(x) = \frac{1}{\sigma_\mathrm{x} \sqrt{2 \pi}} \exp \left ( {- \frac{x^2}{2 \sigma^2_\mathrm{x}}} \right )
    \]

Substituting this in the expression for \(f_\mathrm{y}(y)\) (above) yields:
\[ f_\mathrm{y}(y) = \frac{1}{\sigma_\mathrm{x} \sqrt{2\pi y}} \exp \left ( - \frac{y}{2 \sigma^2_\mathrm{x}} \right) u(y) \]

This is illustrated in Figure \ref{fig:y=gx}. \\

\begin{figure}[h!]
\centering
\includegraphics[scale=0.5]{figures/y_exp_x.png}
\caption{The density function of \(f_\mathrm{y}(y)\) with \(Y = X^2\) and \(X = \mathcal{N}(0,1) \).}
\label{fig:y=gx}
\end{figure}

\textbf{Note:} that for complex \(Y=g(X)\), this method may be inconvenient to use. That is, the  expression of $x$
in terms of $y$ using an inverse transformation may be difficult or impossible.

\subsection{Derivative Approach}

The derivative approach requires determining the roots of the function \(Y = g(X)\), as illustrated in in Figure \ref{fig:roots}. In this diagram, \(x_1, x_2, x_3\) are the roots of y.\\


\begin{figure}[h!]
\centering
\includegraphics[scale=0.5]{figures/y_g_x.png}
\caption{Determining roots of the function \(Y = g(X)\).}
\label{fig:roots}
\end{figure}

Now, 
\[ 
F_\mathrm{y}(y) = P(Y \le y)
\]

and 

\[    \begin{split}
        P(y \le Y \le y + dy)   & = \int_\mathrm{y}^{y + dy} f_\mathrm{y}(y) \approx f_\mathrm{y}(y) |dy| \\
                                & \approx f_\mathrm{y}(y) |dy|    \quad \text{(rectangle under function)}
\end{split}
\]

Also,


\[    \begin{split}
        P(y \le Y \le y + dy)   & = \sum_{i=1}^n P(x_i \le X \le x_i + dx_i) \\
                                & \approx \sum_{i=1}^n f_\mathrm{x}(x_i) |dx_i|
\end{split}
\]

Therefore,

\[
f_\mathrm{y}(y) |dy| \approx \sum_{i=1}^n f_\mathrm{x}(x_i) |dx_i|
\]

and 

\[
f_\mathrm{y}(y) = \sum_{i=1}^n \frac{ f_\mathrm{x}(x_i)}{|dy/dx_i|}
\]

or 

\[
f_\mathrm{y}(y) = \sum_{i=1}^n \frac{ f_\mathrm{x}(x_i)}{|g{'}(x_i)|}
\]

where 

\[
g{'}(x) = \frac{dg(x)}{dx}
\]

\begin{tcolorbox}
    \begin{exmp}
    For \( Y = X^2\) and
    \[ 
    f_\mathrm{x}(x) = \frac{1}{\sigma_\mathrm{x} \sqrt{2 \pi}} \exp \left ( {- \frac{x^2}{2 \sigma^2_\mathrm{x}}} \right )
    \]
 solve for \(f_\mathrm{y}(y)\).
\end{exmp}
\end{tcolorbox}


First,

\[
f_\mathrm{y}(y) = \sum_{i=1}^2 \frac{ f_\mathrm{x}(x_i)}{|g{'}(x_i)|}
\]

Now,

\[    \begin{split}
        f_\mathrm{y}(y)  & = \left \{ \frac{f_\mathrm{x}(\sqrt{y})}{|2 \sqrt{y}| }  + \frac{f_\mathrm{x}(-\sqrt{y})}{|- 2 \sqrt{y}| }  \right \} u(y)\\
                & \quad \text{where } u(y) \text{ is the unit step function)} \\
                & = \frac{1}{\sigma_\mathrm{x} \sqrt{2 \pi}} \left \{ \frac{e^{- \frac{x_1^2}{2 \sigma_\mathrm{x}}}}{2 \sqrt{y}}  + \frac{e^{- \frac{x_2^2}{2 \sigma_\mathrm{x}}}}{2\sqrt(y)} \right \}  u(y) \\
                & = \frac{1}{\sigma_\mathrm{x} \sqrt{2 \pi}} \left \{ \frac{e^{- \frac{y}{2 \sigma_\mathrm{x}}}}{2 \sqrt{y}}  + \frac{e^{- \frac{y}{2 \sigma_\mathrm{x}}}}{2\sqrt{y}} \right \}  u(y) \\
                & = \frac{1}{\sigma_\mathrm{x} \sqrt{2\pi y}} \exp \left ( - \frac{y}{2 \sigma^2_\mathrm{x}} \right) u(y)
\end{split}
\]

This is the same result obtained using the direct approach in the previous section.

\begin{tcolorbox}
    \begin{exmp}{Rectification in an Electrical Circuit.}

Consider an  electrical circuit that has a full-wave rectifier (unity gain). The input voltage to the rectifier \(X\) has a normal distribution with mean \(\mu = 0\) volts and standard deviation \(\sigma = 2\) volts.  Find the PDF of the voltage \(Y\) at the rectifier output.

    \end{exmp}
\end{tcolorbox}

Given that \(X\) follows a normal distribution, the PDF of \(X\) is given by:
\[ f_\mathrm{x}(x) = \frac{1}{\sqrt{2\pi }\sigma} e^{-\frac{(x-\mu)^2}{2\sigma^2}} \]

Substituting the given values \(\mu = 0\) and \(\sigma = 2\), we have:
\[ f_\mathrm{x}(x) = \frac{1}{2\sqrt{2\pi }} e^{-\frac{x^2}{8}} \]

The PDF of the output of the rectifier circuit is:
\[
f_\mathrm{y}(y) = \sum_{i=1}^n \frac{ f_\mathrm{x}(x_i)}{|g^{'}(x_i)|}
\]

For a rectifier circuit, \(Y = |X|\), there are two roots: \(x_1 = -y\) and \(x_2 = y\), and  
\[
|g'(x_i)| = 1    \quad \forall x_i
\]

Therefore

\[
\begin{aligned}
f_\mathrm{y}(y) & = \frac{f_\mathrm{x}(x_1)}{1} + \frac{f_\mathrm{x}(x_2)}{1}\\
            &  = \left \{ f_\mathrm{x}(-y) + f_\mathrm{x}(y) \right \} u(y) \\
            & = \left \{ \frac{1}{\sqrt{2\pi }\sigma} e^{-\frac{(-y-\mu)^2}{2\sigma^2}} + \frac{1}{\sqrt{2\pi }\sigma} e^{-\frac{(y-\mu)^2}{2\sigma^2}} \right \} u(y)
\end{aligned}
\]

This simplifies to:
\[ f_\mathrm{y}(y) = = \left \{ \frac{1}{\sqrt{2\pi }\sigma} e^{-\frac{(y+\mu)^2}{2\sigma^2}} + \frac{1}{\sqrt{2\pi }\sigma} e^{-\frac{(y-\mu)^2}{2\sigma^2}} \right \} u(y)\]

Substituting the given values \(\mu = 0\) and \(\sigma = 2\), we have:

\[
\begin{aligned}
    f_\mathrm{y}(y) & = \left \{ \frac{1}{2\sqrt{2\pi }} e^{-\frac{(y)^2}{8}} + \frac{1}{2\sqrt{2\pi }} e^{-\frac{(y)^2}{8}} \right \} u(y) \\
    & = \frac{1}{\sqrt{2\pi}} e^{-\frac{y^2}{8}}  u(y)
\end{aligned}
\]



\[ f_\mathrm{y}(y) = \frac{1}{ 2\sqrt{2\pi}} e^{-\frac{y^2}{4}} u(y)\]

The pdf of \(X\) and \(Y\) are plotted in Figure \ref{fig:pdf_x_y}. The effect whereas the output of the full-wave rectifier remains normally distributed, the negative values of the input are superimposed on the positive values, resulting in a pdf with twice the amplitude, but only existing for positive values of \(Y\).
\begin{figure}[h!]
\centering
\includegraphics[scale=0.5]{figures/pdf_x_y.png}
\caption{The pdf of the input \(X\) and output \(Y\) of the full-wave rectifier.}
\label{fig:pdf_x_y}
\end{figure}

Sometimes, if the form of \(f_\mathrm{x}(x)\) or \(Y=g(X)\) is complex, it may be difficult to evaluate \(f_\mathrm{y}(y)\). In this case, a numerical evaluation may be required. The pdf \(f_\mathrm{x}(x)\) is used to generate a random vector of samples of \(X\), and the transformation \(Y=g(X)\) is applied. For the above example, using 1000 random samples, these histograms of \(X\) and \(Y\) are generated, as illustrated in Figure \ref{fig:hist_x_y}.

\begin{figure}[h!]
\centering
\includegraphics[scale=0.5]{figures/hist_x_y.png}
\caption{The pdf of the input \(X\) and output \(Y\) of the full-wave rectifier.}
\label{fig:hist_x_y}
\end{figure}


\section{Expected Values \& Moments}

\subsection{The Mean of a Random Variable}

Consider first a discrete random variable $X$ with mean $\mu_\mathrm{x} = E[X]$ Here, $E[X$] is the \emph{expectation} operator whcih is defined as the average value of a random variable $X$ over all possible outcomes.

One can estimate $\mu_\mathrm{x}$ by taking $N$ independent samples of $X$:


\[
\hat{\mu_\mathrm{x}} = \frac{n_1 x_1}{N} + \frac{n_2 x_2}{N} \ldots \frac{n_m x_m}{N}
\]
where the $x_i$ are \(m\) unique values assigned by $X$, and $n_i$ is the number of occurrences of $x_i$.  and $P(X_i)$ is the probability of occurrence of $X_i$.\\

Now, let \( N \rightarrow \infty\), 

\[
\mu_\mathrm{x} = P(x_1)x_1 + P(x_2)x_2 \ldots P(x_m)x_m
\]

and 
\[
\mu_\mathrm{x} = \frac{1}{N} \sum_{i=1}^{N} P(x_i) x_i 
\]

This is now the true value of \( \mu_\mathrm{x}\) and not an estimate \( \hat{\mu_\mathrm{x}}\).\\

For a continuous RV 

\[
\mu_\mathrm{x}  = E[X] = \int_{- \infty}^{\infty} xf_\mathrm{x}(x) \, dx
\]

For a function of a RV \(Y = g(X)\), the expected value of \(Y\) is

\[
\mu_\mathrm{y}  = E[Y] = \int_{- \infty}^{\infty} yf_\mathrm{y}(y) \, dx = \int_{- \infty}^{\infty} g(x) x(x) \, dx
\]

\bigskip
\textbf{Comment:} The expected value of a random variable obeys superposition:
\[
E\left[\sum_{i=1}^n X_i\right] = \sum_{i=1}^n E\left[X_i\right]
\]


\subsection{The Variance of a Random Variable}


\[
\sigma_\mathrm{x}^2 \doteq E[(X - \mu_\mathrm{x})^2]
\]

The standard deviation is
\[
\sigma_\mathrm{x} = \sqrt{E[(X - \mu_\mathrm{x})^2]}
\]

Consider that 
\[    \begin{split}
        \sigma_\mathrm{x}^2  & = E[(X - \mu_\mathrm{x})^2]\\
                    & = E[X^2 - 2X\mu_\mathrm{x} + \mu_\mathrm{x}^2]\\
                    & = E[X^2] - 2\mu_\mathrm{x} E[X] +\mu_\mathrm{x}^2\\
                    & = E[X^2] - \mu_\mathrm{x}^2
\end{split}
\]

or 
\[
E[X^2] = \sigma_\mathrm{x}^2 + \mu_\mathrm{x}^2
\]
In this sense, \(E[X^2]\) is the mean squared value (MSV) and may be considered \emph{total power}, while \(\sigma_\mathrm{x}^2\) is the \emph{AC} power and \(\mu_\mathrm{x}^2\) is the \emph{DC} power. \\

\subsection{$n^{th}$-Order Moments}

Consider the $n^{th}$-order moment of the RV $X$
\[
m_n = E[X^n]
\]

such that 

\[    
\begin{split}
        m_0     & = E[X^0] = 1\\
        m_1     & = E[X^1] = \mu_\mathrm{x}\\
        m_2     & = E[X^2] = MSV
\end{split}
\]

The \emph{central} moments are

\[
\eta_n = E[(X-\mu_\mathrm{x})^n]
\]

such that 

\[    
\begin{split}
        \eta_0     & = E[(X-\mu_\mathrm{x})^0] = 1\\
        \eta_1     & = E[(X-\mu_\mathrm{x})] = 0\\
        \eta_2    & = E[(X-\mu_\mathrm{x})^2] = \sigma_\mathrm{x}^2
\end{split}
\]



\subsection{The Characteristic Function of a Random Variable}\label{sec:moments}

The characteristic function is defined as 
\[
\begin{split}
    \Phi_\mathrm{x}(\omega) & = E\left[ e^{j\omega X} \right] \\
                    & = \int_{-\infty}^{\infty}f_\mathrm{x}(x) e^{j\omega X} \, dx
\end{split}
\]\\

Upon inspection, this is the Fourier transform of the density function $f_\mathrm{x}(x)$. The characteristic function encodes all the statistical properties of the random variable $X$ and is particularly useful for:

\begin{itemize}
    \item \textbf{Finding Moments:} The moments of a random variable can be obtained from its characteristic function. 
    \item \textbf{Sum of Independent Random Variables}: The characteristic function of the sum of independent random variables is equal to the product of their individual characteristic functions. This property makes it especially useful in finding the distribution of the sum of independent random variables.
    \item \textbf{Convolution of Distributions:} The characteristic function is also helpful in finding the distribution resulting from the convolution of two probability distributions. The convolution theorem relates the characteristic functions of the convolved distributions to the characteristic functions of the original distributions.
\end{itemize}


\textbf{The Moment Theorem} The characteristic function is sometimes referred to as the Moment Generating Function (MGF). 
\bigskip

If 
\[
m_n = E[X^n]
\]

then 
\[
j^n m_n = \left. {\frac{d^n}{dt^n} \phi_\mathrm{x}(\omega)} \right|_{\omega=0}
\]


\textbf{Proof:}
\begin{enumerate}
    \item Start with the definition of the MGF:
    \[
    \Phi_\mathrm{x}(\omega) = E[e^{jX\omega}]
    \]
    
    \item Expand the exponential function as its Maclaurin series:
    \[
    e^{jX\omega} = 1 + (jX\omega) - \frac{1}{2!}(jX\omega)^2 + \frac{1}{3!}(jX\omega)^3 - \ldots
    \]
    
    \item Take the expected value of both sides:
    \[
    \Phi_\mathrm{x}(\omega) = E[1] + j\omega E[X] - \frac{1}{2!}(j\omega)^2 E[X^2] + \frac{1}{3!}(j\omega)^3 E[X^3] - \ldots
    \]
    
    \item Recognize that the coefficients of the Maclaurin series are actually moments:
    \[
    \Phi_\mathrm{x}(\omega) = 1 + j\omega E[X] - \frac{1}{2!}(j\omega)^2 E[X^2] + \frac{1}{3!}(j\omega)^3 E[X^3] - \ldots
    \]
    
    \item Differentiate both sides with respect to $\omega$:
    \[
    \frac{d}{d\omega}\Phi_\mathrm{x}(\omega) = 0 + jE[X] - (j\omega)E[X^2] + \frac{(j\omega)^2}{2!}E[X^3] - \ldots
    \]
    
    \item Evaluate the derivative at $\omega=0$ to find the first moment:
    \[
    \left. \frac{d}{d\omega}\Phi_\mathrm{x}(\omega) \right|_{\omega=0} = jE[X]
    \]
    
    \item To find higher moments, differentiate more times and evaluate at $\omega=0$:
    \[
    \left. \frac{d^n}{d\omega^n}\Phi_\mathrm{x}(\omega) \right|_{\omega=0} = (j)^n E[X^n]
    \]
    
\end{enumerate}

This proves the theorem.

\bigskip
\textbf{Example: Convolution Using Characteristic Functions}

Let's consider two independent random variables, \(X\) and \(Y\), with probability density functions \(f_\mathrm{x}(x)\) and \(f_\mathrm{y}(y)\) and respective characteristic functions \(\Phi_\mathrm{x}(t)\) and \(\Phi_\mathrm{y}(t)\). We want to find the probability density function \(f_\mathrm{z}(z)\) for the random variable \(Z = X + Y\).\\

The characteristic function of \(Z\) is given by:
\[
\Phi_\mathrm{z}(\omega) = E[e^{itZ}] = E[e^{it(X+Y)}] = E[e^{itX}e^{itY}] = \Phi_\mathrm{x}(t) \cdot \Phi_\mathrm{y}(t)
\]


\[
\begin{split}
f_\mathrm{z}(z)  & = \mathcal{F}^{-1} \left \{ \Phi_\mathrm{z}(\omega) \right \}  \\
        & = \mathcal{F}^{-1} \left \{ \Phi_\mathrm{x}(\omega) \cdot \Phi_\mathrm{y}(\omega) \right \} \\
        & = \mathcal{F}^{-1} \left \{ \Phi_\mathrm{x}(\omega) \right \} * \mathcal{F}^{-1} \left \{ \Phi_\mathrm{y}(\omega) \right \}
\end{split}
\]

Therefore,
\[
f_\mathrm{z}(z) = f_\mathrm{x}(x) * f_\mathrm{y}(y)
\]

This demonstrates how characteristic functions can be used to find the probability density function of the sum of independent random variables.\\

In the limit, for \(Z = X_1 + X_2 + X_2 + \ldots\) independent RVs,
\[
f_\mathrm{z}(z) = f_\mathrm{x_1}(x_1) * f_\mathrm{x_2}(x_2) * f_\mathrm{x_3}(x_3) *\cdots
\]

This implies how the central limit theorem works; these independent densities will smooth each other toward Gaussian with many convolutions.    

\section{Statistical Independence, Correlation, and Orthogonality}

Consider two random variables, $X$ and $Y$. Let $E[XY]$ denote the joint expectation of $X$ and $Y$, and $E[X]$ and $E[Y]$ represent their individual expectations.

\begin{itemize}

   
   \item $X$ and $Y$ are \emph{uncorrelated} if their covariance is zero:
     \[
     \text{Cov}(X, Y) = E[XY] - E[X]E[Y] = 0
     \]

     or
     \[
     E[XY] = E[X]E[Y]
     \]
     \begin{itemize}
         \item Uncorrelation indicates that linear relationships between $X$ and $Y$ are absent. However, uncorrelated random variables can still be dependent in a nonlinear sense, and therefore not independent.

     \end{itemize}
       \item $X$ and $Y$ are \emph{statistically independent} if and only if their joint probability density function $f_\mathrm{xy}(x, y)$ can be factorized into the product of their marginal density functions:
     \[
     f_\mathrm{xy}(x, y) = f_\mathrm{x}(x) \cdot f_\mathrm{y}(y)
     \]
     \begin{itemize}
         \item This means that knowing the value of one variable provides no information about the other. In this case, conditional probabilities factor as well: $P(X = x | Y = y) = P(X = x)$ and vice versa.
             \item Two random variables being uncorrelated suggests only the lack of a linear predictive relationship, while independent indicates NO relationship at all, including all joint moments up to and including the entire pdf. For this reasons, proving (or even asserting) independence is usually difficult.
     \end{itemize}
     
   \item \emph{Orthogonality} is a stronger condition than uncorrelation. Two random variables $X$ and $Y$ are orthogonal if and only if their expected product is zero:
     \[
     E[XY] = 0
     \]
     \begin{itemize}
         \item Orthogonal random variables are uncorrelated, but they go beyond that. They are independent in the sense that their joint behavior is completely decoupled.
     \end{itemize}
\end{itemize}

In summary, statistical independence implies uncorrelation, but uncorrelation does not necessarily imply statistical independence. Orthogonality is a stronger condition than uncorrelation and implies statistical independence.


\chapter{Stochastic (Random) Processes}

\section{Introduction}
Stochastic processes are  mathematical models used to describe the evolution of random phenomena over time. A stochastic process is essentially a collection of random variables indexed by time or another parameter (such as position). These random variables represent the uncertain nature of events that evolve with time. By studying stochastic processes, we gain insights into probabilistic behavior, make predictions about future events, and model complex systems influenced by randomness.\\

A stochastic process can be modeled as an ensemble of sample functions, which themselves are a function of time, as depicted in Figure \ref{fig:stochastic_ensemble}. Each sample function represents a possible time series that might be selected as an outcome of an experiment, or equivalently, the assignment by a random variable.\\

\begin{figure}[h!]
\centering
\includegraphics[scale=0.3]{figures/Stochastic-process.png}
\caption{A graphical depiction of a stochastic process and an ensemble of sample functions.}
\label{fig:stochastic_ensemble}
\end{figure}

Here:
 \begin{itemize}
     \item The outcome of an experiment \(A_i =  \left \{ A_1, A_2, A_3, A_4 \right \}\) allows the random process to assign a sample function \(X(t,A_i) = x_i(t)\) where \(x_i(t)\) is drawn from   \(\left \{ x_1(t), x_2(t), x_3(t), x_4(t) \right \}\).
     \item At some specific time \(t_1\), the random process is reduced to a random variable \(X(t_1,A_i) = X(t_1)\).
 \end{itemize}

Thus, a random process is a collection of sample functions and also a collection of random variables. Here, we define the full notation for random processes, sample functions, and random variables using the outcome of the experiment \(A_i\) (left side), and the shorthand notation (right) where the \(A_i\) is implicit. 
\begin{itemize}
    \item \(X(t,A_i) = X(t) \rightarrow\) random process
    \item \(X(t,A_1) = x_1(t) \rightarrow\) sample function
    \item \(X(t_1,A_i) = X(t_1) \rightarrow\) random variable
    \item \(X(t_1,A_1) = x_1(t_1) \rightarrow\) a number assigned by the RV \(X(t_1)\)
\end{itemize}

We will use the shorthand notation from this point onward.

\section{Distribution Functions and Density Functions}

The distribution and density functions of a stochastic process are those describing the random variables at all possible time instants, \(t_i\).\\

For a random variable \(X(t_i)\) at time \(t_i\): \\

Distribution function:
$F_\mathrm{x_i} (x) \rightarrow F_\mathrm{x}(x, t_i)$.\\

Density function:
$f_\mathrm{x_i}(x) \rightarrow f_\mathrm{x}(x, t_i)$.\\

Joint density function:
$f_\mathrm{x_i x_j}(x_i, x_j) \rightarrow$ Joint density for two random variables at different times, \(t_i, t_j\).\\

To completely define the process \(X(t)\), we would require a knowledge of the \(n^{th}\)-order joint probability density function as \(n \rightarrow\infty\) since there are an infinite number of random variables on the time axis. It is not usual to have this knowledge. 

\[
\lim_{{n \to \infty}} f_\mathrm{x_1 X_2 \ldots, X_n} (x_1, x_2, \ldots, x_n)
\]\\

Fortunately however, for most applications, only \(2^{nd}\)-order statistics are required.\\


\section{\(1^{st}\) Order Properties of a Stochastic Process}

In the context of a stochastic process, we often encounter first-order properties that provide essential insights into the process's behavior. 

\begin{itemize}
  \item \textbf{Mean of the process \(X(t)\)}. Denoted as $\mu_\mathrm{x}(t)$, this represents the expected value or average of the process at time $t$. It is defined as:
  \[
  \mu_\mathrm{x}(t) = E[X(t)] = \int_{-\infty}^{\infty} xf_\mathrm{x}(x, t) \, dx
  \]\\

  Note that in general, \(\mu_\mathrm{x}(t)\) and \(f_\mathrm{x}(x, t)\) may vary with time. \\

  At \(t=t_1\):
 \[
  \mu_\mathrm{x}(t_1) = E[X(t_1)] = \int_{-\infty}^{\infty} xf_\mathrm{x}(x, t_1) \, dx
  \]
  or equivalently,
  
  \[
  \mu_\mathrm{x}(t_1) = E[X_1] = \int_{-\infty}^{\infty} x f_\mathrm{x_1}(x) \, dx
  \]
  \\


  \item \textbf{Variance of the process \(X(t)\)}. 

  \[\sigma_\mathrm{x}^2(t) = E\left [\left(X(t)-\mu_\mathrm{x}(t) \right)^2 \right ]
  \]
  
\end{itemize}

If these properties change with time, the process is
referred to as \emph{non-stationary}.

\section{\(2^{nd}\) Order Properties of a Stochastic Process}

The autocorrelation function is an important concept in the study of stochastic processes. It quantifies the degree of linear dependence or correlation between a stochastic process and its past values at different time lags. The autocorrelation function provides insights into the process's temporal dependencies, periodicities, and trends. \\

\textbf{Definition.} 

\[
\begin{split}
    \phi_\mathrm{x_1 x_2}(t_1,t_2) & = E\left [ X(t_1) X(t_2) \right ] = E[X_1 X_2]\\
    \\
    & = \int_{-\infty}^{\infty}\int_{-\infty}^{\infty}x_1x_2f_\mathrm{x_1 x_2}(x_1,x_2) \, dx_1 \, dx_2
\end{split}
\]


\begin{tcolorbox}
    \begin{exmp}{Random binary transmission (non-zero mean).}

Suppose we have a stochastic process $X(t)$  representing the binary transmission illustrated in Figure \ref{fig:random_binary}. Assume that $X(t)$ can take values 0 or 1 with equal probabilities. That is:
\[P(0) = P(1) = \frac{1}{2}\]
This is true at any time \(t\). \\

Find the mean and the autocorrelation function of this random process.
    \end{exmp}
\end{tcolorbox}


First, assume that $X(t)$ can take values 0 or 1 with equal probabilities. That is:
\[P(0) = P(1) = \frac{1}{2}\]\\

This is true at any time \(t\).\\


\begin{figure}[h!]
\centering
\includegraphics[scale=0.4]{figures/Random_binary.png}
\caption{Two sample functions from a random binary transmission (bitstream).}
\label{fig:random_binary}
\end{figure}

Now consider the mean of the process. 
\[
\begin{split}
    E[X(t)] & = \int_{-\infty}^{\infty} x f_\mathrm{x}(x,t) \, dx \\
            & = (0)P(0) + (1)P(1)\\
            & = \frac{1}{2}
\end{split}
\]

The variance of the process is:
\[
\begin{split}
    \sigma_\mathrm{x}^2(t) & = E \left [ (X(t)-\mu_\mathrm{x}(t))^2 \right ]\\
            & = E \left [ X^2(t) \right ] - E\left [ X(t) \right ]^2 \\
            & = (0)^2P(0) + (1^2)P(1) - \left (\frac{1}{2} \right )^2 \\
            & = \frac{1}{4}
\end{split}
\]

Computing the autocorrelation function is a bit more involved. Recall the definition:

\[
    \phi_\mathrm{x_1 x_2}(t_1,t_2) = E\left [ X(t_1) X(t_2) \right ]
\]

We must consider two separate cases:\\

\textbf{Case 1:} \( |t_2-t_1|>T\)\\
Here, \(X(t_1)\) and \(X(t_2)\) are uncorrelated, and so
\[
E\left [ X(t_1) X(t_2) \right ] = E\left [ X(t_1) \right ] E\left [X(t_2)\right ] = \frac{1}{4}
\]

\textbf{Case 2:} \( |t_2-t_1|\le T\)\\
Here, \(X(t_1)\) and \(X(t_2)\) are uncorrelated if \(t_1\) and \(t_2\) fall outside of the same bit and \emph{correlated} if they fall inside the same bit. These conditions are  illustrated in Figure \ref{fig:random_binary_interval}.\\

Therefore, the expected value is;
\[
\begin{split}
    E\left [ X(t_1) X(t_2) \right ] & = P(\text{same interval})E_s[X(t_1) X(t_2)] + P(\text{different interval})E_d[X(t_1) X(t_2)]\\
        & = \left ( 1-\frac{|t_2-t_1|}{T} \right)E_s[X(t_1) X(t_2)] + \left (\frac{|t_2-t_1|}{T} \right)E_d[X(t_1) X(t_2)]
\end{split}
\]

Here, \(E_s[X(t_1) X(t_2)]\) and \(E_d[X(t_1) X(t_2)]\) denote the expected value in the same bit and different bit, respectively. For a random binary transmission, \(E_d[(X_1 X_2]=1/4\), as in case 1.\\

For \(X(t_1)\) and \(X(t_2)\) in the same interval, \(X_1=X_2\) so that

\[ 
E_s[X(t_1) X(t_2)] = P(\left \{ X(t_1)=1 \text{ and } X(t_2)=1 \right \}) = \frac{1}{2}.
\]

Finally, 
\[
\begin{split}
    E\left [ X(t_1) X(t_2) \right ] & = P(\text{same interval})E_s[X(t_1) X(t_2)] + P(\text{different interval})E_d[X(t_1) X(t_2)]\\
        & = \frac{1}{2}\left ( 1-\frac{|t_2-t_1|}{T} \right) + \frac{1}{4}\left (\frac{|t_2-t_1|}{T} \right)\\
        & = \frac{1}{4}\left ( 1-\frac{|t_2-t_1|}{T} \right) + \frac{1}{4}
\end{split}
\]

Now, \( \tau = t_2 - t_1\) and
\[ 
E[X(t_1) X(t_2)] = \begin{cases}
\frac{1}{4}\left ( 1-\frac{|\tau|}{T} \right) + \frac{1}{4} & \text{for } |\tau|<T \\
\frac{1}{4} & \text{for } |\tau| \ge T
\end{cases}
\]

This is depicted in Figure \ref{fig:autocorr_binary}.


\begin{figure}[h!]
\centering
\includegraphics[scale=0.4]{figures/Binary_Interval.png}
\caption{A random binary interval with \(t_1\) and \(t_2\) falling in different (top) and the same (bottom) bit.}
\label{fig:random_binary_interval}
\end{figure}

\begin{figure}[h!]
\centering
\includegraphics[scale=0.5]{figures/Autocorr_binary.png}
\caption{The autocorrelation function of a random binary transmission.}
\label{fig:autocorr_binary}
\end{figure}

Note that:
\begin{itemize}
    \item \( \phi_\mathrm{xx}(0) = E[X(t_i)X(t_i)] = MSV = \sigma_\mathrm{x}^2 + \mu_\mathrm{y}^2 = \frac{1}{2} \)
    \item \( \phi_\mathrm{xx}(\infty) =  \mu_\mathrm{y}^2(t) = \frac{1}{4}\)
\end{itemize}
\bigskip

\begin{tcolorbox}
    \begin{exmp}{Random binary transmission (zero mean).}\\
Suppose we have a stochastic process $X(t)$ representing the binary transmission, as illustrated in Figure \ref{fig:random_binary-zm}. Assume that $X(t)$ can take values -1 or 1 with equal probabilities. That is:
\[P(-1) = P(1) = \frac{1}{2}\]
This is true at any time \(t\). \\

Find the mean and the autocorrelation function of this random process.

\end{exmp}
\end{tcolorbox}



\begin{figure}[h!]
\centering
\includegraphics[scale=0.4]{figures/Random_binary_zm.png}
\caption{A random binary transmission (zero mean).}
\label{fig:random_binary-zm}
\end{figure}

Now consider the mean of the process. 
\[
\begin{split}
    E[X(t)] & = \int_{-\infty}^{\infty} x f_\mathrm{x}(x,t) \, dx \\
            & = (-1)P(-1) + (1)P(1)\\
            & = 0
\end{split}
\]

The variance of the process is:
\[
\begin{split}
    \sigma_\mathrm{x}^2(t) & = E \left [ (X(t)-\mu_\mathrm{x}(t))^2 \right ]\\
            & = E \left [ X^2(t) \right ] - E\left [ X(t) \right ]^2 \\
            & = (1)^2P(-1) + (1^2)P(1) - (0)^2 \\
            & = 1
\end{split}
\]

Computing the autocorrelation function is a bit more involved. Recall the definition:

\[
    \phi_\mathrm{x_1 X_2}(t_1,t_2) = E\left [ X(t_1) X(t_2) \right ]
\]

We must consider two separate cases:\\

\textbf{Case 1:} \( |t_2-t_1|>T\)\\
Here, \(X(t_1)\) and \(X(t_2)\) are uncorrelated, and so
\[
E\left [ X(t_1) X(t_2) \right ] = E\left [ X(t_1) \right ] E\left [X(t_2)\right ] = 0
\]

\textbf{Case 2:} \( |t_2-t_1|\le T\)\\
Here, \(X(t_1)\) and \(X(t_2)\) are uncorrelated if \(t_1\) and \(t_2\) fall outside of the same bit (as illustrated in Figure \ref{fig:random_binary_interval}, and \emph{correlated} if they fall inside the same bit.\\

Therefore, the expected value is;
\[
\begin{split}
    E\left [ X(t_1) X(t_2) \right ] & = P(\text{same interval})E_s[X(t_1) X(t_2)] + P(\text{different interval})E_d[X(t_1) X(t_2)]\\
        & = \left ( 1-\frac{|t_2-t_1|}{T} \right)E_s[X(t_1) X(t_2)] + \left (\frac{|t_2-t_1|}{T} \right)E_d[X(t_1) X(t_2)]
\end{split}
\]

Here, \(E_s[X(t_1) X(t_2)]\) and \(E_d[X(t_1) X(t_2)]\) denote the expected value in the same bit and different bit, respectively. For a random binary transmission, \(E_d[(X_1 X_2]=0\), as in case 1.\\

For \(X(t_1)\) and \(X(t_2)\) in the same interval, the product \(X_1 X_2\) is always equal to 1.
 
\[ 
E_s[X(t_1) X(t_2)] = 1.
\]

Finally, 
\[
\begin{split}
    E\left [ X(t_1) X(t_2) \right ] & = P(\text{same interval})E_s[X(t_1) X(t_2)] + P(\text{different interval})E_d[X(t_1) X(t_2)]\\
        & = \left( 1-\frac{|t_2-t_1|}{T} \right )
\end{split}
\]

Since \(X(t)\) is stationary, then \( \tau = t_2 - t_1\) and
\[ 
E[X(t_1) X(t_2)] = \begin{cases}
\left ( 1-\frac{|\tau|}{T} \right)  & \text{for } |\tau|<T \\
0 & \text{for } |\tau| \ge T
\end{cases}
\]

This is depicted in Figure \ref{fig:autocorr_binary_zero}.


\begin{figure}[h!]
\centering
\includegraphics[scale=0.5]{figures/Autocorr_binary_zero.png}
\caption{The autocorrelation function of a random binary transmission (zero mean).}
\label{fig:autocorr_binary_zero}
\end{figure}

It is trivial to extend the result of the zero mean case to the non-zero mean case since for \(Y = aX + \mu\):

\[
\phi_\mathrm{yy}(\tau) = a^2\phi_\mathrm{xx}(\tau) + \mu^2
\]

The non-zero mean process has 1/2 the amplitude of the zero mean process, and is offset by 1/2. This yields
\[ 
\phi_\mathrm{yy}(\tau) = \begin{cases}
\frac{1}{4}\left ( 1-\frac{|\tau|}{T} \right) + \frac{1}{4} & \text{for } |\tau|<T \\
\frac{1}{4} & \text{for } |\tau| \ge T
\end{cases}
\]

which is the result of the previous example.\\

\begin{tcolorbox}
    \begin{exmp}{Random telegraph signal (zero mean).}
    
Suppose we have a stochastic process $X(t)$ where at any instant, with equal probability, the value either $+A$ or $-A$, as depicted in Figure \ref{fig:random_telegraph}. \\

The probability that $n$ sign changes occur in a time interval \((t_1, t_2)\) is given by the given Poisson distribution:

\[
P(n, t_1, t_2) = \frac{\lambda (t_2 - t_1)^n}{n!} \exp \left (-\lambda \cdot (t_2 - t_1) \right )
\]

where \(\lambda\) is the average rate of sign changes per unit time.\\

Find the mean and the autocorrelation function of this random process.
    
    \end{exmp}
\end{tcolorbox}




\begin{figure}[h!]
\centering
\includegraphics[scale=0.5]{figures/Random_telegraph.png}
\caption{A random telegraph signal.}
\label{fig:random_telegraph}
\end{figure}

Now, let \(X_1 = X(t') \) and \(X_2 = X(t'+\tau)\).

\[ 
E[X(t_1) X(t_2)] = \begin{cases}
A^2 & \text{for } X_1 = X_2 \\
-A^2 & \text{for } X_1 = -X_2
\end{cases}
\]

\[
\begin{split}
    \phi_\mathrm{xx}(t_1,t_2) & = A^2 P(A^2) - A^2 P(-A^2)\\
                    & = A^2 P(X_1=X_2) - A^2P(X_1=-X_2)
\end{split}
\]

If \(n\) sign changes occur between \(t_1\) and \(t_2\), then
\[
P(X_1=X_2) = P(n \text{ even})
\]
\[
P(X_1=-X_2) = P(n \text{ odd})
\]

\[
\begin{split}
P(n \text{ even}) & = P(n = 0, 2, 4, \ldots)\\
    & = \sum_{n=0}^{\infty} \frac{(\lambda \tau) ^n}{n!} \exp(\lambda \tau) \quad \text{for \emph{n} even}\\
    & = \exp(-\lambda \tau) \cosh(\lambda \tau)
\end{split}
\]

\[
\begin{split}
P(n \text{ odd}) & = P(n = 1, 3, 5, \ldots)\\
    & = \sum_{n=0}^{\infty} \frac{(\lambda \tau) ^n}{n!} \exp(\lambda \tau) \quad \text{for \emph{n} odd}\\
    & = \exp(-\lambda \tau) \sinh{}h(\lambda \tau)
\end{split}
\]

\[
\begin{split}
\phi_\mathrm{xx}(\tau) & = A^2 \exp(-\lambda \tau) \left \{ \cosh(\lambda \tau) - \sinh(\lambda \tau)\right \} \\
        & =  A^2 \exp(-\lambda \tau)  \exp(-\lambda \tau)\\
        & = A^2 \exp(- 2 \lambda \tau) \quad \text{for } t_2>t_1\\
    & = \exp(-2 \lambda |\tau|) \quad \text{for all } t_2, t_1
\end{split}
\]

This is illustrated in Figure \ref{fig:random_telegraph_AC}.


\begin{figure}[h!]
\centering
\includegraphics[scale=0.3]{figures/Random_telegraph_AC.png}
\caption{The autocorrelation function of a random telegraph signal.}
\label{fig:random_telegraph_AC}
\end{figure}


\begin{tcolorbox}
    \begin{exmp}{Sinusiod with random phase offset.}
    
Suppose we have a stochastic process \(X(t) = a \cos(\omega_0 t + \theta)\), where \(a\) and \(\omega_0\) are constants, and \(\theta\) is a random variable uniformly distributed over the interval \([0, 2\pi]\) as depicted in Figure \ref{fig:random_cosine}. \\

Find the mean and autocorrelation function of this random process.

\end{exmp}
\end{tcolorbox}


\begin{figure}[h!]
\centering
\includegraphics[scale=0.6]{figures/Random_cosine.png}
\includegraphics[scale=0.4]{figures/theta_dist.png}
\caption{Top: A cosine process with random phase shift. Bottom: Density function of a uniformly distributed phase shift \(\theta\).}
\label{fig:random_cosine}
\end{figure}

\textbf{Mean:}
\[
\begin{split}
\mu_\mathrm{x}(t) = E[X(t)] & = E\left[a\cos(\omega_0 t + \theta)\right] \\
    & = \int_{-\infty}^{\infty} a\cos(\omega_0 t + \theta) f_{\theta}(\theta) d\theta \\
    & = \int_{-\infty}^{\infty} a\cos(\omega_0 t + \theta) \cdot \frac{1}{2\pi} d\theta \\
    & = \frac{a}{2\pi} \int_0^{2\pi} \cos(\omega_0 t + \theta) d\theta \\
    & = 0
\end{split}
\]

\textbf{Variance:}
\[
\begin{split}
\sigma_\mathrm{x}(t)^2 = E[(X(t)-\mu_\mathrm{x}(t))^2] & = E\left[a\cos(\omega_0 t + \theta)\right] \\
    & = E[(X(t)^2] \\
    & = \frac{a^2}{2}
\end{split}
\]

\textbf{Autocorrelation:}
\[
\begin{split}
\phi_\mathrm{xx}(t) = E[X(t_1)X(t_2)] & = E\left[a\cos(\omega_0 t_1 + \theta) a\cos(\omega_0 t_2 + \theta) \right] \\
    & = E\left[\frac{a^2}{2}\cos(\omega_0 (t_2-t_1) ) \right] + E\left[\frac{a^2}{2}\cos(\omega_0 (t_2+t_1) + 2\theta) \right] \\
    & = \frac{a^2}{2}\cos(\omega_0 (t_2-t_1) ) \\
    & = \frac{a^2}{2}\cos(\omega_0 \tau ) 
\end{split}
\]

\begin{figure}[h!]
\centering
\includegraphics[scale=0.4]{figures/cosine_AC.png}
\caption{The autocorrelation function of a cosine process with random phase shift.}
\label{fig:cosine_ac}
\end{figure}


\subsection{Cross-Correlation Function}

 The cross-correlation function between two random processes \(X(t) \) and \(Y(t) \) is a measure of similarity between them as a function of the displacement of one relative to the other. It is defined as:

\[ \phi_\mathrm{xy}(t_1, t_2) = E[X(t_1)Y(t_2)] \]

This is illustrated in Figure \ref{fig:cross_corr_ac}.

\begin{figure}[h!]
\centering
\includegraphics[scale=0.35]{figures/Cross_correlation.png}
\caption{The cross-correlation function of processes \(X(t)\) and \(Y(t)\).}
\label{fig:cross_corr_ac}
\end{figure}

Now, \( X(t) \) and \(Y(t) \) are \emph{uncorrelated} random processes if

\begin{itemize}
    \item \( \phi_\mathrm{xy}(t_1,t_2) = E[X(t_1)]E[Y(t_2)] = \mu_\mathrm{x}(t_1) \mu_\mathrm{y}(t_2)\) \quad \quad \quad for all \(t_1, t_2\)
    \item \( C_\mathrm{xy}(t_1,t_2) = E[X(t_1)]E[Y(t_2)] - \mu_\mathrm{x}(t_1) \mu_\mathrm{y}(t_2)\) = 0 \quad cross-covariance
\end{itemize}

They are \emph{orthogonal} if \( \phi_\mathrm{xy}(t_1,t_2) = 0\).\\

To be \emph{independent} processes:

\[
f_\mathrm{x_1 x_2 \cdots y_1 y_2 \cdots}(x_1, x_2, \ldots, y_1, y_2, \ldots) = f_\mathrm{x_1 x_2 \cdots}(x_1, x_2, \ldots) f_\mathrm{y_1 y_2} \cdots(y_1, y_2, \ldots)
\]\\

In general, all statistics and density functions may depend on time. If so, the process is \emph{non-stationary}.

\section{Stationary Processes}

\textbf{Strict-Sense Stationary (SSS)}: All statistics and density functions must be time-invariant.\\

\textbf{Wide-Sense Stationary (WSS)}: Up to \(2^{nd}\) order statistics (moments) must be time-invariant.

\[
\begin{split}
\mu_\mathrm{x}(t) & \rightarrow \mu_\mathrm{x} \\
\sigma_\mathrm{x}^2(t) & \rightarrow \sigma_\mathrm{x}^2 \\
\phi_\mathrm{xx}(t_1,t_2) & \rightarrow \phi_\mathrm{xx}(t_2-t_1) = \phi_\mathrm{xx}(\tau)
\end{split}
\]
This is a valid criterion for stationarity in many engineering applications.

\subsection{Ensemble  and Time Averages}
\textbf{Ensemble Averages:}

\begin{itemize}
    \item \textbf{Definition:} Ensemble averages involve taking the average of a statistical property (e.g., mean, variance) over an ensemble of multiple realizations (sample functions) of a stochastic process.
    \item \textbf{Calculation:} To calculate ensemble averages, you consider a collection of independent realizations of the stochastic process and average the corresponding values of the property of interest over these realizations.
    \item \textbf{Interpretation:} Ensemble averages provide statistical insights into the behavior of the stochastic process as a whole. They describe how the process behaves on average across different realizations.
\end{itemize}

\textbf{Time Averages:}

\begin{itemize}
    \item \textbf{Definition:} Time averages involve taking the average of a single realization (sample function) of a stochastic process over a specified time interval.
    \item \textbf{Calculation:} For a given time series, one calculates time averages by sampling values of the process at different time points and then averaging them over the chosen time interval (within which the process must be ergodic).
    \item \textbf{Interpretation:} Time averages provide insights into the behavior of a specific realization of a stochastic process. 
\end{itemize}

An illustration of ensemble and time averages is given in Figure \ref{fig:time_ensemble}. 

\begin{figure}[h!]
\centering
\includegraphics[scale=0.3]{figures/time_ensemble.png}
\caption{The calculation of an ensemble and time average mean. The ensemble average is computed across all (preferably infinite) sample functions at a specified time. Here, at \(t=t_1\),  \(\mu_\mathrm{x}(t_1) = E[X(t_1)]\). The time average mean computed on the \(3^{rd}\) sample function is  \(M_x = \frac{1}{T} \int_{0}^{T} x_3(t) dt \). Ideally, \(T\rightarrow \infty\) but clearly, this is not practical.}
\label{fig:time_ensemble}
\end{figure}


With unlimited time to compute the time average statistics of \(X(t)\) we define the mean as:
\[
M_\mathrm{x} = \lim_{T \to \infty} \frac{1}{T} \int_{0}^T X(t) dt
\]

The \emph{time-average autocorrelation function} is 

\[
R_\mathrm{xx}(\tau) = \lim_{T \to \infty} \frac{1}{2T} \int_{-T}^{T} X(t)X(t+\tau) dt
\]


\(M_\mathrm{x}\) and \( R_\mathrm{xx}(\tau)\) are RV because the random process assigns a sample function from which they are computed.

In general:

\[
\mu_\mathrm{x} \neq M_\mathrm{x}
\]

\[
\phi_\mathrm{xx}(T) \neq R_\mathrm{xx}(\tau)
\]

for since the ensemble values are constants/vectors for a WSS process, and the time-average values are random variables.


\subsection{Ergodic Stochastic Processes}

If `time averages' are equal to `ensemble averages,' then $X(t)$ is said to be ergodic. Implications of ergodicity are as follows:

\begin{enumerate}
    \item The time average can be determined from \emph{any} individual sample functions (\emph{i.e.}, each sample function must yield the same result, since \( M_\mathrm{x}=\mu_\mathrm{x}, R_\mathrm{xx}(\tau) = \phi_\mathrm{xx}(\tau)\)).
    \item Ergodicity implies stationarity, since $\mu_\mathrm{x}(t)$, $\sigma_\mathrm{x}(t)$, and $\phi_\mathrm{xx}(\tau)$ must be time-invariant.

\end{enumerate}

\begin{tcolorbox}
    
\begin{exmp}{A stationary random process that is not ergodic.}\\

    Consider a stationary random process, denoted as \(X(t)\), which alternates between two constant values, \(A\) and \(-A\) amongst sample functions, with equal probability. Mathematically, it can be represented as:

\[
X(t) = \begin{cases} 
A, & \text{with probability } 0.5 \\
-A, & \text{with probability } 0.5 
\end{cases}
\]
\end{exmp}
\end{tcolorbox}

This process is stationary because its statistical properties (mean, variance, autocorrelation) do not change with time. However, it is not ergodic because the time averages and ensemble averages are not equal.\\

\begin{itemize}
    \item Time Averages: 
    \begin{itemize}
        \item If you observe a single realization of this process (a single sample function), you will see that it takes the value of either \(A\) and \(-A\) for all time.
        \item The time average of this single realization will be either A or -A.
    \end{itemize}
    \item Ensemble Averages:
    \begin{itemize}
        \item When considering the ensemble of all possible realizations, you will have an equal number of realizations where the time average is \(A\) and where it is \(-A\).
        \item So, the ensemble average will be zero.
    \end{itemize}

\end{itemize}


\subsection{Wide-Sense Ergodicity}

Wide-sense ergodicty requires ergodicty of the mean and autocorrelation function. For the process $X(t)$:

\begin{enumerate}
    \item $\mu_\mathrm{x} = M_\mathrm{x}$ (time-average mean)
    \item $\phi_\mathrm{xx}(\tau) = R_\mathrm{xx}(\tau)$ (time-average autocorrelation).
\end{enumerate}


\textbf{Ergodicity of the mean:}

For infinite averaging time \(T\):
\[
M_\mathrm{x} = \lim_{T \to \infty} \frac{1}{2} \int_{0}^{T} x(t) dt
\]

For finite averaging time \(T\):

\[
M_T =  \frac{1}{T} \int_{0}^{T} x(t) dt
\]

Since \(M_T\) is a random variable, then to be ergodic:
\[
E[M_T] = \frac{1}{T} \int_{0}^{T} E[X(t)] dt = \mu_\mathrm{x}.
\]

The variance of the estimate \(M_T\) decreases as the window size \(T\) increases, as illustrated in Figure \ref{fig:time_avg_mean}.

\begin{figure}[h!]
\centering
\includegraphics[scale=0.3]{figures/time_avg_mean.png}
\caption{The estimate of the time average mean \(M_T\), extracted for a few sample functions \(x_i(t)\), as the window size \(T\) increases. If the process is ergodic, then \(E[M_T] = \mu_\mathrm{x}  \forall     T\) and \(\text{Var} [M_T] \rightarrow 0 \) as  \(T \rightarrow \infty\).}
\label{fig:time_avg_mean}
\end{figure}

The conditions for ergodicity of the mean are:
\begin{itemize}
    \item \(E[M_T] = \mu_\mathrm{x}\) \quad for all \(T\)
    \item \(\text{Var} [M_T] = E[M_T^2] - E[M_T]^2 \rightarrow 0 \) \quad as  \(T \rightarrow \infty\)
\end{itemize}

\bigbreak

\textbf{Ergodicity of the autocorrelation function:}
\[
R_\mathrm{xx}(\tau) = \lim_{T \to \infty} \frac{1}{2T} \int_{-T}^{T} x(t)x(t+\tau) dt
\]
(a random variable).

Now, define

\[
R_T(\tau) = \frac{1}{2T} \int_{-T}^{T} x(t)x(t+\tau) dt
\]
(also a random variable). \\

Conditions for ergodicty of the autocorrelation function:
\begin{itemize}
    \item \(E[R_T(\tau)] = \frac{1}{2T} \int_{-T}^{T} x(t)x(t+\tau) dt = \phi_\mathrm{xx}(\tau)\) \quad for all \(T\)
    \item \(\text{Var} [R_T(\tau)] \rightarrow 0 \) \quad as  \(T \rightarrow \infty\)
\end{itemize}

The variability of \(R_T\) as the window size \(T\) increases from \(T_1\) to \(T_2\) is illustrated in Figure \ref{fig:time_avg_mean}.

\begin{figure}[h!]
\centering
\includegraphics[scale=0.3]{figures/time_avg_ac.png}
\caption{The estimate of \(R_T\) as the window size increases from \(T_1\) to \(T_2\).}
\label{fig:time_avg_ac}
\end{figure}

\subsection{Properties of \(\phi_\mathrm{xx}(\tau)\) and \(\phi_\mathrm{xy}(\tau)\) (Stationary Processes)}

\begin{enumerate}
    \item \(\boldsymbol{\phi_\mathrm{xx}(\tau)}\) \textbf{ is an even function.}\\
    We know that 
     \[ 
    \phi_\mathrm{xx}(\tau) = E[x(t)x(t+\tau)]
    \]
    
    Consider that 
    \[ 
    \phi_\mathrm{xx}(-\tau) = E[x(t)x(t-\tau)]
    \]
    

   
Now, let \(\alpha = t - \tau\), which yields
 \[ 
 \begin{split}
         \phi_\mathrm{xx}(-\tau) & = E[x(\alpha+\tau)x(\alpha)] \\
         \phi_\mathrm{xx}(-\tau) & = \phi_\mathrm{xx}(\tau)
\end{split}
\]

    \item \textbf{The value of }\(\phi_\mathrm{xx}(0)\) \textbf{is the mean squared value.}
    \[
    \phi_\mathrm{xx}(0)\ = E[X^2(t)] = MSV
    \]

    \item \textbf{The value of } \(\phi_\mathrm{xx}(0)\) \textbf{is greater than or equal to all other values of} \(\phi_\mathrm{xx}(\tau)\).\\
    A typical non-periodic autocorrelation function (\emph{e.g.} the autocorrelation function of a random telegraph signal) is depicted in Figure \ref{fig:random_telegraph_AC} which has its peak at \(\tau=0\).   

\bigskip
\textbf{Proof:}
Consider 
\[
E \left [ \left \{ X(t) \pm X(t + \tau) \right \} ^2 \right ] \ge 0
\]

\[
\phi_\mathrm{xx}(0) +  \phi_\mathrm{xx}(0) \pm 2\phi_\mathrm{xx}(\tau)\ge 0
\]

\[
2 \left \{ \phi_\mathrm{xx}(0) \pm \phi_\mathrm{xx}(\tau) \right \} \ge 0
\]

Therefore,

\[
\phi_\mathrm{xx}(0) \ge | \phi_\mathrm{xx}(\tau) \quad \text{ for all } \tau \ge 0.
\]

This is illustrated in Figure \ref{fig:MSV_peak}. A periodic process may have peaks recurring at multiples of the period, as with the cosine process with random phase shift (Figure \ref{fig:cosine_ac}). 


\begin{figure}[h!]
\centering
\includegraphics[scale=0.5]{figures/MSV_Peak.png}
\caption{The value of \( \phi_\mathrm{xx}(0) = MSV \) is the peak for a non-periodic process.}
\label{fig:MSV_peak}
\end{figure}

\item \textbf{The value of }\(\phi_\mathrm{xx} (\tau) \rightarrow \mu_\mathrm{x}^2 \) \textbf{ as }  \( \tau \rightarrow \infty\).

\bigskip
\textbf{Proof:}

\[
\begin{split}
    \phi_\mathrm{xx}(\tau) & = E[X(t) X(t+ \tau)]\\
                    & = \int_{-\infty}^{\infty}\int_{-\infty}^{\infty}x_1 x_2 f_\mathrm{x_1 x_2}(x_1,x_2) \, dx_1 \, dx_2
\end{split}
\]

If, as \(\tau \rightarrow \infty\), \(X_1\) and \(X_2\) are independent, then 
\[
\begin{split}
    \phi_\mathrm{xx}(\tau) & = \int_{-\infty}^{\infty}x_1  f_\mathrm{x_1}(x_1) \, dx_1 \int_{-\infty}^{\infty}x_2 f_\mathrm{x_2}(x_2) \, dx_2 \\
                    & = \mu_\mathrm{x_1}\mu_\mathrm{x_2} \\
                    & = \mu_\mathrm{x}^2
\end{split}
\]

Again, this is illustrated in Figure \ref{fig:MSV_peak}.

\item \textbf{The cross-correlation has reciprocity}.\\

The reciprocity property of cross-correlation (also known as the symmetry property) states that swapping the order of the two signals in a cross-correlation calculation is equivalent to reversing the time-lag (or index) of the resulting correlation function. 

\bigskip
\textbf{Proof:}

\[
\begin{split}
    \phi_\mathrm{xy}(\tau) & = E[X(t)Y(t + \tau)]\\
    \phi_\mathrm{xy}(-\tau) & = E[X(t)Y(t - \tau)] \quad \text{let } \alpha=t-\tau \\
                    & = E[X(\alpha + \tau)Y(\alpha)]\\
                    & = E[Y(\alpha)X(\alpha + \tau)]\\
     \therefore \phi_\mathrm{xy}(-\tau)               & = \phi_\mathrm{yx}(\tau)
\end{split}
\]
\end{enumerate}

\bigskip

\begin{tcolorbox}
    
\begin{exmp}{A sinusoidal process with random phase shift, corrupted by bandlimited white noise.}\\

Consider a stochastic process represented by the equation:
\[Y(t) = X(t) + A \sin(\omega_0 t + \alpha) \]
where:
\begin{itemize}
  \item \(X(t)\) is a noise process.
  \item \(A\) and \(\omega_0\) are constants.
  \item \(\alpha\) is a random variable uniformly distributed over the interval \([0, 2\pi]\) and is independent of \(X(t)\).
\end{itemize}

Find the autocorrelation function $\phi_\mathrm{yy}(\tau)$.
\end{exmp}
\end{tcolorbox}

Now,
\[
\begin{split}    
\phi_\mathrm{yy}(\tau) & = E\left [ \{X(t) + A \sin (\omega_0 t + \alpha)\}\{X(t+\tau) + A \sin (\omega_0 (t + \tau)+ \alpha)\}\ \right ] \\
    & = \phi_\mathrm{xx}(\tau) +  E\left [X(t)A \sin (\omega_0 (t+\tau) + \alpha) \right ] \\
    & \quad  + E\left [X(t+\tau)A \sin (\omega_0 t + \alpha) \right ] \\
    & \quad  + E\left [A^2 \sin (\omega_0 t + \alpha) \sin (\omega_0 (t+\tau) + \alpha)\right ]  \quad  \\
\end{split}
\]

Now using \( \sin A \sin B = 1/2 \{\cos (A-B) - \cos (A+B)\}\),\\

\[
\begin{split}    
\phi_\mathrm{yy}(\tau) & = \phi_\mathrm{xx}(\tau) + \frac{A^2}{2}E\left [ \cos(\omega_0 \tau) - \cos(2 \omega_0 t + \omega_0 \tau + 2 \alpha)\right ] \\
    & = \phi_\mathrm{xx}(\tau) + \frac{A^2}{2} \cos(\omega_0 \tau)
\end{split}
\]

This is illustrated in Figure \ref{fig:Noisy_ACF}. If \(X(t)\) is bandlimited white noise, \(\phi_\mathrm{yy}(\tau)\) becomes equal to that of the periodic signal as the noise becomes uncorrelated with itself as \(\tau \rightarrow \) large. 

\begin{figure}[h!]
\centering
\includegraphics[scale=0.5]{figures/Noisy_ACF.png}
\caption{The autocorrelation function of a sinusoidal process with random phase shift, corrupted by bandlimited white noise. For the autocorrelation function of bandlimited white noise, please see Section \ref{BLWN}.}
\label{fig:Noisy_ACF}
\end{figure}




\subsection{Autocovariance function}

If \(X(t)\) is a stationary process:

\begin{itemize}
    \item \( \phi_\mathrm{xx}(t) \rightarrow \) autocorrelation function
    \item \( C_\mathrm{xx}(t) \rightarrow \) autocovariance function
\end{itemize}

   \[
\begin{split}
   C_\mathrm{xx}(t) & = E\left[(X(t) - \mu_\mathrm{x})(X(t + \tau) - \mu_\mathrm{x})\right] \\
                & = \phi_\mathrm{xx}(t) - \mu_\mathrm{x} E[X(t) - \mu_\mathrm{x}] E[X(t + \tau)] + \mu_\mathrm{x}^2\\
                & = \phi_\mathrm{xx}(t) - \mu_\mathrm{x}^2
\end{split}
\]

Note:
\[
C_\mathrm{xx}(0) = \sigma_\mathrm{x}^2 = MSV - \mu_\mathrm{x}^2 \quad \text{(variance).}
\]

\section{Power Density Spectrum}

The power density spectrum, or power spectral density (PSD) of a stationary process \(X(t)\) is related to its autocorrelation function \(\phi_\mathrm{xx}(\tau)\) by the Fourier transform. The PSD, denoted by \(\Phi_\mathrm{xx}(f)\), is given by:

\[ \Phi_\mathrm{xx}(f) = \int_{-\infty}^{\infty} \phi_\mathrm{xx}(\tau) e^{-j2\pi f \tau} d\tau \]

This relationship provides a frequency-domain representation of the energy distribution in the process. The PSD is an important tool in understanding the frequency content of a stochastic process.

\[
\phi_\mathrm{xx}(\tau) \Leftrightarrow \Phi_\mathrm{xx}(\omega)
\]

\subsubsection{Interpretation of the PSD as a density function}

Consider a random process \(X(t)\) (in volts) driving a resistive circuit, as illustrated in Figure \ref{fig:power_circuit}.
\begin{figure}[h!]
\centering
\includegraphics[scale=0.4]{figures/Power_circuit.png}
\caption{Power dissipated in a resistive circuit}
\label{fig:power_circuit}
\end{figure}




The instantaneous power \(P(t)\) dissipated in \(R\) is a random process
\[
P(t) = \frac{X^2(t)}{R} \quad \text{watts}
\]

without loss of generality, let \(R=1 \Omega\). Therefore, \(P(t) = X^2(t) \text{ volts}^2 \).\\

The average power \(\overline{P} \) is also a random variable:
\[
\begin{split}
\overline{P} & = \lim_{T \to \infty} \frac{1}{2T} \int_{-T}^{T} P(t) \, dt\\
                & = \lim_{T \to \infty} \frac{1}{2T} \int_{-T}^{T} X^2(t) \, dt
\end{split}
\]

Now,
\[
E[\overline{P}] = E \left [ \lim_{T \to \infty} \frac{1}{2T} \int_{-T}^{T} X^2(t) \, dt \right ] = MSV = \phi_\mathrm{xx}(0) \quad \text{volts}^2
\]

But we know that
\[
\phi_\mathrm{xx}(\tau) = \int_{-\infty}^{\infty} \Phi_\mathrm{xx}(f) e^{-j2\pi f \tau} \, df \quad \text{volts}^2
\]

and therefore
\[
\phi_\mathrm{xx}(0) = \int_{-\infty}^{\infty} \Phi_\mathrm{xx}(f) \, df \quad \text{volts}^2
\]

Therefore, \(\Phi_\mathrm{xx}(f)\) has units of \(V^2/Hz\) and is a density function in the frequency domain.

\subsection{Properties of \(\Phi_\mathrm{xx}(f)\)}

\begin{enumerate}
    \item Since \(\phi_\mathrm{xx}(\tau)\) is an even function, \(\Phi_\mathrm{xx}(f)\) is \emph{real} and \emph{even}.
    \item \(\phi_\mathrm{xx}(0) = \int_{-\infty}^{\infty} \Phi_\mathrm{xx}(f) \, df \) is \emph{the average power} of \(X(t)\) - the MSV.
    \item \(\Phi_\mathrm{xy}(f)\) is a \emph{complex} function since \(\phi_\mathrm{xy}(\tau)\) is an \emph{odd} function.
    \item Since \(\phi_\mathrm{xy}(\tau) = \phi_\mathrm{yx}(-\tau)\) it follows that \( P_\mathrm{xy}(f) = P_\mathrm{yx}^*(f)\) - the complex conjugate.

\end{enumerate}


\subsection{White noise process.}
Pure white noise is not a realizable process (infinite frequencies!), but it is a useful theoretical reference. It has equal power density at all (infinite) frequencies, and an autocorrelation function equal to a Dirac delta function at \(\tau=0\), as illustrated in Figure \ref{fig:PSD_WN}.  

\begin{figure}[h!]
\centering
\includegraphics[scale=0.4]{figures/PSD_WN.png}
\caption{The power spectral density and the autocorrelation function of white noise.}
\label{fig:PSD_WN}
\end{figure}



\subsection{Bandlimited white noise process.}\label{BLWN}

Bandlimited white noise contains equal power at all frequencies within a specific frequency range (\(-B\) to \(B\) in the double-sided PSD), but it's limited to that range. Unlike standard white noise, which extends infinitely across all frequencies, bandlimited white noise is confined to a certain bandwidth, as illustrated in Figure \ref{fig:PSD_BLWN}.\\

\begin{figure}[h!]
\centering
\includegraphics[scale=0.4]{figures/PSD_BLWN.png}
\caption{The power spectral density and the autocorrelation function of bandlimited white noise.}
\label{fig:PSD_BLWN}
\end{figure}


The autocorrelation function is expressed as the \emph{sinc} function:

\[
\phi_\mathrm{xx} (\tau) = 2BK\frac{\sin(2 \pi B \tau)}{2 \pi B \tau }
\]
Note the reciprocal relationship between bandwidth in the frequency domain and the width of the autocorrelation function in the time domain. Reducing bandwidth has the effect of imposed more correlation in the time domain (broadening the autocorrelation function).   





\chapter{Estimation}

\section{Introduction}
In this chapter, we are interested in finding a function of random  variables used to estimate an unknown parameter of an underlying  distribution. Such a function is called a \emph{statistic}, and it generally reduces the number of random variables, called samples, to a lower dimension without loss of information about the parameter. The scenario can be summarized:

\begin{itemize}
    \item Let \(X(t)\) be a random process with fixed probability density function that is parameterized by a scalar \(\theta\). Assume that \(\theta\) is not observable: it cannot be sampled directly and information about \(\theta\) is obtained only through samples of \(X(t)\).
    \item The random process is sampled \(N\) times, yielding \(N\) independent and identically distributed (i.i.d.) random variables \(\{X_1, X_2 \ldots X_N\}\). Although the RVs \(\{X_N\}\) contain identical information about the parameter \(\theta\), the specific samples \(\{x_N\}\) generally vary.
    
\end{itemize}

The \(\{X_N\}\) could also be derived from samples of a random sequence \(X[k]\) or they may be just be \(N\) samples of a particular random variable with respect to time. We will use a subscript on \(X_n\) to indicate that it is a sample (random variable) of either the process \(X(t)\), sequence \(X[k]\) or random variable \(X\). The samples need not be collected from the process/sequence with uniform spacing in time; we can simply view \(\{X_N\}\) as a set of i.i.d. random variables with parameter \(\theta\).\\

In order to simplify the equations, we use \(\mathbf{X} \triangleq \{X_1, X_2 \ldots X_N\}\) to represent the \(N\) samples and \(\mathbf{x} \triangleq \{x_1, x_2 \ldots x_N\}\) to represent their specific outcomes.\\

This chapter explores various estimation techniques, ranging from classical frequentist approaches to Bayesian methodologies\footnote{The primary difference between frequentist 
 and Bayesian methods is that the parameter is a considered a fixed quantity versus a random variable, respectively. Refer to \hyperlink{https://towardsdatascience.com/statistics-are-you-bayesian-or-frequentist-4943f953f21b}{this article}{ for theoretical and practical differences between the two approaches.}}, each offering unique insights into parameter estimation. Four key methods that will be examined in detail are Bayesian Estimation, Maximum A Posteriori (MAP) Estimation,  Maximum Likelihood (ML) Estimation, and Linear Minimum Mean Squared Error (LMMSE) Estimation.


\section{Estimator Performance}
To evaluate the efficacy of an estimator, it is important to establish metrics of estimator performance.\\


\subsubsection*{Bias}
    \begin{itemize}
        \item \emph{Definition: }The bias of an estimator \(\hat{\theta}\) with respect to the true parameter \(\theta\) is given by 
        \[
        Bias(\hat{\theta}) = E[\hat{\theta}] - \theta
        \]
        \item \emph{Interpretation: } An estimator is unbiased if \(E[\hat{\theta}] = \theta\), meaning that, on average, the estimator's expected value equals the true parameter value.
    \end{itemize}
   
\subsubsection*{Variance}
\begin{itemize}
    \item \emph{Definition: } The variance of an estimator \(\hat{\theta}\) is given by 
    \[
    \begin{aligned}
        Var(\hat{\theta}) & = E[(\hat{\theta} - E[\hat{\theta}])^2]\\
            & = E[\hat{\theta}^2] - E[\hat{\theta}]^2
    \end{aligned}
    \]
    
   \item \emph{Interpretation: } Variance measures the spread or variability of the estimator around its expected value.
\end{itemize}

\subsubsection*{Mean Squared Error (MSE)}
 \begin{itemize}
    \item \textbf{Definition: } MSE of an estimator \(\hat{\theta}\) is given by 
    \[
    MSE(\hat{\theta}) = E[(\hat{\theta} - \theta)^2]
    \]
    
   \item \textbf{Interpretation: } MSE combines both bias and variance to quantify the average squared difference between the estimated values and the true parameter value. That is:
   \[
    \begin{aligned}
        MSE(\hat{\theta}) &=  E[(\hat{\theta} - \theta)^2]\\
                            & = E[\hat{\theta}^2  - 2 \hat{\theta} \theta + \theta^2] \\
                            & = E[\theta^2] - 2 E[\hat{\theta}] \theta + E[\hat{\theta}]^2 +  E[\hat{\theta}^2] - E[\hat{\theta}]^2 \\
                            & = Bias^2(\hat{\theta}) + Var(\hat{\theta})
    \end{aligned} 
   \]
   \item An estimator is said to be \emph{consistent} if $\text{MSE} \rightarrow$ as \(T\rightarrow \infty\).
\end{itemize}

 

\subsubsection*{Relationships}
\begin{itemize}
    \item \emph{Consistency and Bias:} Consistent estimators are typically asymptotically unbiased, meaning they become unbiased as the sample size increases.
    \item \emph{Bias and Variance in MSE: } MSE can be decomposed into bias squared and variance. Minimizing MSE involves finding a balnce between reducing bias and variance.
\end{itemize}




\section{Bayesian Estimation}
\label{sec:Bayesian_Estimation}

Bayesian estimation is a statistical framework for estimating parameters of a model using Bayes' theorem. In Bayesian statistics, the parameter of interest (\( \theta\)) is treated as a random variable, and its uncertainty is quantified using probability distributions. The estimation process involves updating prior beliefs about the parameter with observed data \(X\) to obtain a posterior distribution.\\

We may think of \(\theta\) and \(X\) as events, in which case we can use probabilities to describe the elements of Bayes estimation. If \(\theta\) and \(X\) are continuous random variables, we can also express the elements of Bayes' theorem in terms of their probability density functions.\\

Here are the key components of Bayesian estimation:

\begin{enumerate}
    \item \textbf{Prior Distribution (Prior):} This is the initial belief or probability distribution about the parameter before observing any data. The prior represents existing knowledge or beliefs about the parameter and is denoted as \( P(\theta) \) (probability of $\theta$) or \(f_\theta(\theta)\) (probability density function of $\theta$).
    \item \textbf{Likelihood Function (Likelihood):} This function represents the probability of observing the data \(X\) given a specific parameter value $\theta$. It is denoted as \( P(X | \theta) \) or \( f_\mathrm{{x}|\theta}({x})\).
    \item \textbf{Posterior Distribution (Posterior):} The posterior distribution is the updated probability distribution of the parameter after taking into account the observed data. It is calculated using Bayes' theorem:

   \[ P(\theta | X) = \frac{P(X | \theta) \cdot P(\theta)}{P(X)} \]

    Here, \( P(X) \) is the marginal likelihood: the probability of observing the data irrespective of the parameter. When \(\theta\) and \(X\) are continuous random variables it is common to use probability density functions:

       \[ f_\mathrm{\theta | x} (\theta) = \frac{f_\mathrm{x | \theta} (x) \cdot f_\theta (\theta)}{f_\mathrm{x}(x)} \]

The posterior distribution summarizes our updated beliefs about the parameter of interest after incorporating the observed data. It combines our prior beliefs with the new evidence provided by the data, using the likelihood function to assess how well different values of the parameter(s) explain the observed data. As more data is collected, the posterior distribution can be continually updated to refine our estimates of the parameter.

\end{enumerate}

\textbf{Vector Notation of the Data}\\
Recall that, in general, the random process my be sampled \(N\) times, yielding \(N\) independent and identically distributed (i.i.d.) random variables such that \(\mathbf{X} \triangleq \{X_1, X_2 \ldots X_N\}\). In this case:
\[
f_\mathrm{x}({x}) = \prod_{n=1}^N f_\mathrm{x_n}(x_n)
\]

and

\[
f_\mathrm{{x}|\theta}({x}) = \prod_{n=1}^N f_\mathrm{x_n|\theta}(x_n)
\]

In the following sections the simpler notation for the data will be \(X\) with the acknowledgement that it can be used interchangeably with the vector \(\mathbf{X}\).\\

\textbf{Posterior Estimation}\\ The parameter estimate in Bayesian statistics is often summarized using the posterior distribution. The point estimate can be the mean, median, or mode of the posterior distribution, and credible intervals (analogous to confidence intervals in frequentist statistics) can be derived to express the uncertainty around the estimate.\\

\textbf{Advantages of Bayesian Estimation}
\begin{itemize}
    \item \emph{Incorporation of Prior Information:} Bayes' theorem allows the incorporation of prior knowledge or beliefs into the estimation process.
    \item \emph{Quantification of Uncertainty:} Provides a full probability distribution for the parameter, allowing for a more comprehensive understanding of uncertainty.
    \item \emph{Flexibility:} Can be applied to a wide range of problems, especially when dealing with complex models or small sample sizes.
\end{itemize}

\textbf{Challenges of Bayesian Estimation}
\begin{itemize}
    \item \emph{Subjectivity of Priors:} The choice of prior can impact the results, and different analysts may have different priors.
    \item \emph{Computational Complexity:} In some cases, obtaining the exact posterior distribution may be computationally challenging, leading to the use of approximation methods like Markov Chain Monte Carlo (MCMC).
\end{itemize}

\bigskip

\begin{tcolorbox}
    \begin{exmp}{Bayesian estimation.}

    Consider using a Bayesian estimator to estimate the mean $\mu$ of a Gaussian process \(X \sim \mathcal{N}(\mu,\sigma)\). Assume  the standard deviation $\sigma$ of the distribution is known.
    \end{exmp}
\end{tcolorbox}

Using a Bayesian approach to estimate the mean \(\mu\) of a Gaussian process \(X \sim N(\mu, \sigma^2)\) involves updating our beliefs about \(\mu\) based on observed data. We'll start with a prior distribution for \(\mu\), then update it using Bayes' theorem after observing data points from the process.\\

Let's denote:
\begin{itemize}
    \item \(X = \{x_1, x_2, ..., x_n\}\) as the observed data points.
    \item \(N\) as the total number of observations.
    \item \(\sigma^2\) as the known variance of the process.
\end{itemize}

After observing data \(X\), the posterior distribution for \(\mu\) can be calculated using Bayes' theorem:

\[ f(\mu | X) = \frac{f(X | \mu) \cdot f(\mu)}{f(X)} \]

Where:
\begin{itemize}
    \item \( f(\mu | X) \) is the posterior distribution for \(\mu\) given the observed data \(X\)
    \item \( f(X | \mu) \) is the likelihood function, representing the likelihood of observing data \(X\) given \(\mu\), which is given by the product of \(N\) Gaussian probability density functions
    \item \( f(\mu) \) is the prior distribution for \(\mu\), assumed to be normal
    \item \( f(X) \) is the marginal likelihood or evidence, which can be calculated by integrating the likelihood function over all possible values of \(\mu\).
\end{itemize}


Once we have the posterior distribution \( f(\mu | X) \), we can summarize it by calculating its mean, median, mode, or credible intervals, which provide a range of plausible values for \(\mu\) given the observed data. \\


First, let's define our prior knowledge about $\mu$. We'll assume that $\mu$ itself follows a Gaussian distribution with known mean $\mu_0$ and standard deviation $\sigma_0$. This is our prior distribution \(f(\mu) = \mathcal{N}\ (\mu_0, \sigma_0)\).\\

After observing some data $X = \{x_1, x_2, ..., x_n\}$, we can calculate the likelihood function $\mu$, representing the likelihood of observing data $X$ given $f(\mu)$. Recall that this is given by the product of $n$ Gaussian probability density functions with known variance $\sigma$. This is computed as the sample estimates of mean and variance from the Gaussian process \(X\) with:
\[
\begin{aligned}
 \hat{\mu} & = \frac{1}{n}\sum_{i=1}^n x_i\\
 \hat{\sigma} & =  \frac{1}{n}\sigma
 \end{aligned}
\]

We can then apply Bayes' theorem to calculate the posterior distribution $P(\mu|X)$, which is proportional to the product of the likelihood and the prior:

$$ f(\mu|X) \propto f(X|\mu)f(\mu) $$

The posterior distribution is also Gaussian, since it is the product of two Gaussian distributions. 

We multiply the Prior kernel by the Likelihood kernel:

$$f(\mu | x) \propto \underbrace{\exp\left( -\frac{(\mu - \mu_0)^2}{2\sigma_0^2} \right)}_{\text{Prior}} \times \underbrace{\exp\left( -\frac{(x - \mu)^2}{2\sigma^2} \right)}_{\text{Likelihood}}$$

Since we are multiplying exponentials, we add the terms inside the brackets:$$f(\mu | x) \propto \exp\left( -\frac{1}{2} \left[ \frac{(\mu - \mu_0)^2}{\sigma_0^2} + \frac{(x - \mu)^2}{\sigma^2} \right] \right)$$


Before we evaluate the mean and variance of the  the product of two exponential functions, let's define \emph{precision} as the inverse of the variance, so that
\[
p = \frac{1}{\sigma^2}
\]

The precision of the posterior distribution is then
\[
\begin{aligned}
    p_{\text{posterior}} & = \frac{1} {\sigma^2_{\text{prior}}} + \frac{1}{\hat{\sigma}} \\
        & = \frac{1} {\sigma^2_0} + \frac{n}{\sigma^2}
\end{aligned}
\]
And,

\[
\sigma^2_{p} = \frac{1}{p_{p}}
\]

The mean of the posterior distribution. After completing the square and normalizing, the new mean $\mu_p$ emerges as a weighted sum:

$$\mu_p = \frac{\text{Prior Precision} \cdot \mu_0 + \text{Data Precision} \cdot \hat{\mu}}{\text{Total Precision}}$$Mathematically, this looks like:
\[
\begin{aligned}
\mu_{p} & = \frac{\frac{\mu_0}{\sigma_0^2} + \frac{\hat{\mu}}{\sigma^2}}{\frac{1}{\sigma_0^2} + \frac{n}{\sigma^2}}
 \\
\end{aligned}
\]

This can be simplified to give
\[ 
\mu_{p} = \frac{\sigma^2\mu_0 + n\sigma_0^2\hat{\mu}}{\sigma^2 + n\sigma_0^2}
\]


The Bayesian estimate is a weighted average of the prior mean $\mu_0$ and the sample mean $\hat{\mu}$, with more weight given to whichever has smaller variance. This reflects the trade-off between prior knowledge and observed data in Bayesian estimation. If we have a lot of data (large $n$), the sample mean will be more influential. If we have strong prior knowledge (small $\sigma_0$), the prior mean will be more influential. \\

This example is evaluated in numerical form in the code \href{https://www.ece.unb.ca/kengleha/Textbook/Code/Bayes_estimate_mu.py}{\texttt{Bayes\_estimate\_mu.py}}. Figure \ref{fig:Bayes_estimate_mu} uses this code to evaluate the posterior density and compute the resulting mean. Here, the prior mean is $\mu_0=0$, and the example uses a prior standard deviation that is large ($\sigma_0=2.0$) and small $\sigma_0=0.5$.  \\

\begin{figure}[h!]
\centering
\includegraphics[scale=0.5]{figures/Bayes_estimate_mu_large_sigma.png}
\includegraphics[scale=0.5]{figures/Bayes_estimate_mu_small_sigma.png}
\caption{Bayesian estimate of the mean of a process. After an update using four new data samples,  the resulting updated posterior density is more affected by the new data for the large prior $\sigma_0=2.0$ than for $\sigma_0=0.5$.}
\label{fig:Bayes_estimate_mu}
\end{figure}

\begin{tcolorbox}
\begin{exmp}{Bayesian estimation of mean response time of a new drug.}

      Let's consider an example involving the estimation of the mean response time of a new drug in a clinical trial. Suppose a new drug is designed to reduce the response time of patients to a specific stimulus. The response times are assumed to follow a normal distribution, and we want to estimate the mean response time (\(\mu\)) of the population. We can use Bayesian estimation to update our belief about \(\mu\) based on observed data. In this case, the response time is updated from a an observed average response time, after adding \(n\) new patients to a study.
      \end{exmp}
\end{tcolorbox}

     

\begin{enumerate}
    \item Prior Distribution:
   We start with a prior distribution representing our initial belief about the mean response time. Let's assume a normal distribution with a mean \(\mu_0\) and a relatively wide standard deviation (\(\sigma_0\)) to reflect our uncertainty before any data is collected.

   \[ \text{Prior: } f(\mu) \sim \mathcal{N}(\mu_0, \sigma_0) \]

    \item Likelihood Function:
   After conducting a clinical trial, we collect response time data from patients. We model the likelihood of the observed data given the parameter \(\mu\), assuming that the response times are normally distributed.

   \[ \text{Likelihood: } f(\text{Data} | \mu) \sim \mathcal{N}(\mu, \sigma) \]

    \item Posterior Distribution:
   Using Bayes' theorem, we update our belief about \(\mu\) based on the prior and likelihood. The posterior distribution reflects our updated knowledge after considering the observed data.

   \[ \text{Posterior: } f(\mu | \text{Data}) \sim \mathcal{N}(\mu_p, \sigma_p) \]

   The posterior parameters (\(\mu_p, \sigma_p\)) are calculated based on the prior parameters (\(\mu_0, \sigma_0^2\)) and the observed data.

    \item Bayesian Estimation:
   We can calculate the posterior mean (\(\mu_p\)) as our Bayesian estimate of the mean response time. The posterior standard deviation (\(\sigma_p\)) provides information about the uncertainty of our estimate.

\end{enumerate}

This example is evaluated in numerical form in the code \href{https://www.ece.unb.ca/kengleha/Textbook/Code/Bayes_clinical_trial.py}{\texttt{Bayes\_clinical\_trial.py}}. Before adding new patients to the study, the response time is \( \sim \mathcal{N}(50,20) \).  Thirty new patients enter the study, drawn from a \( \sim \mathcal{N}(45,15) \) likelihood distribution. Figure \ref{fig:Bayes_clinical_trial} shows the prior density, the new data and likelihood distribution, and the resulting updated posterior density. It is clear that the updated response time is less than the prior.

\begin{figure}[h!]
\centering
\includegraphics[scale=0.4]{figures/Bayes_clinical_trial.png}
\caption{Bayesian estimate of the response time of a new drug, before and after adding 30 new patients to the study.}
\label{fig:Bayes_clinical_trial}
\end{figure}


\subsection{Maximum a Posteriori Estimation }
    

Maximum a Posteriori (MAP) estimation is a Bayesian point estimate of the most probable value of a parameter (the mode of the probability density function) given observed data and prior knowledge. 


Recall that Bayes' theorem relates the posterior \(P(\theta | X)\), likelihood \(P(X | \theta)\), prior \(P(\theta)\), and evidence \(P(X)\):

\[ P(\theta | X) = \frac{P(X | \theta) \cdot P(\theta)}{P(X)} \]

Generally, it is assumed that the  likelihood \(P(X | \theta)\) and the prior \(P(\theta)\) are known or can be estimated. The estimate of the parameter \(\theta\) does not depend directly on the evidence \(P(X)\); consequently it does not affect the maximum and can be ignored. The MAP estimator for \(\theta\) is the value maximizing the posterior pdf \footnote{Note that if the prior is uniform or does not significantly affect the result, the MAP estimate becomes equivalent to Maximum Likelihood Estimation \( \theta_{\text{MLE}} = \arg \max_\theta P(X | \theta) \). This is discussed in the next section.}:
 
\[ \theta_{\text{MAP}} = \arg \max_\theta \left[ P(X | \theta) \cdot P(\theta) \right] \]



Often, to simplify the expression, one can take the logarithm of both sides:

\[ \log P(\theta | X) = \log P(X | \theta) + \log P(\theta) - \log P(X) \]


Maximizing the posterior probability is equivalent to finding the maximum of the logarithm of the posterior:

\[ \theta_{\text{MAP}} = \arg \max_\theta \left[ \log P(X | \theta) + \log P(\theta) \right] \]

The MAP estimation provides a principled way to incorporate prior knowledge into the parameter estimation process, making it particularly useful when dealing with limited data.

\begin{tcolorbox}
    
\begin{exmp}{Example with Gaussian Likelihood and Prior.}

For a Gaussian likelihood \(P(X | \theta) = \mathcal{N}(\theta, \sigma)\) and a Gaussian prior \(P(\theta) = \mathcal{N}(\mu_0, \sigma_0)\), find the MAP estimate of $\theta$.

\end{exmp}
\end{tcolorbox}

The map estimate of $\theta$ is given by:
\[ \log P(\theta | X) = \log P(X | \theta) + \log P(\theta)\]

where \( X = \{x_1, x_2, \ldots, x_n\} \) is a vector of data points. Let's consider first the probability density functions: 

   \[ f(X | \theta) = \prod_{i=1}^{n} \frac{1}{\sqrt{2\pi\sigma^2}} \exp\left(-\frac{(x_i - \theta)^2}{2\sigma^2}\right) \]


   \[ f(\theta) = \frac{1}{\sqrt{2\pi\sigma_0^2}} \exp\left(-\frac{(\theta - \mu_0)^2}{2\sigma_0^2}\right)\]


The joint posterior distribution is given by:

\[ f(\theta | X) \propto f(X | \theta) \cdot f(\theta) \]

Substituting the expressions for the joint likelihood and the joint prior, we have:

\[ f(\theta | X) \propto \prod_{i=1}^{n} \frac{1}{\sqrt{2\pi\sigma^2}} \exp\left(-\frac{(x_i - \theta)^2}{2\sigma^2}\right) \cdot \frac{1}{\sqrt{2\pi\sigma_0^2}} \exp\left(-\frac{(\theta - \mu_0)^2}{2\sigma_0^2}\right) \]

Taking the logarithm of this expression makes it easier to work with:

\[ \ln f(\theta | X) \propto \sum_{i=1}^{n} \left( -\frac{(x_i - \theta)^2}{2\sigma^2} \right) - \frac{(\theta - \mu_0)^2}{2\sigma_0^2} + C \]

Where \( C \) is a constant term.\\

To find the MAP estimate of \( \theta \), we need to maximize \( \log f(\theta | X) \) with respect to \( \theta \). We can do this by taking the derivative of \( \log f(\theta | X) \) with respect to \( \theta \), setting it to zero, and solving for \( \theta \).

\[ \frac{d}{d\theta} \ln f(\theta | X) = \sum_{i=1}^{n} \frac{x_i - \theta}{\sigma^2} - \frac{\theta - \mu_0}{\sigma_0^2} = 0 \]

Solving this equation for \( \theta \), we get the MAP estimate:

\[ \hat{\theta}_\text{MAP} = \frac{\frac{\sum_{i=1}^{n} x_i}{\sigma^2} + \frac{\mu_0}{\sigma_0^2}}{\frac{n}{\sigma^2} + \frac{1}{\sigma_0^2}} \]

This is illustrated in Figure \ref{fig:MAP1_example}.


\begin{figure}[h!]
\centering
\includegraphics[scale=0.6]{figures/MAP1_example.png}
\caption{The prior distribution (green), posterior distribution (red), and a histogram of the observed data (orange). The Maximum A Posteriori (MAP) estimate of the parameter $\theta$
 is indicated by the dashed gray line. The data histogram represents the observed dataset $X$.}
\label{fig:MAP1_example}
\end{figure}

This expression provides the MAP estimate of \( \theta \) given the dataset \( X \), the known standard deviation \( \sigma \) of the likelihood, the prior mean \( \mu_0 \), and the prior standard deviation \( \sigma_0 \).

\bigskip

\newpage
\begin{tcolorbox}
    \begin{exmp}{Bit Error Rate in a Communications System.}

    Let's consider a communication system where we want to estimate the probability of bit errors (\(p\)) in a digital communication channel. 
    \end{exmp}
\end{tcolorbox}

We have some prior information about the bit error probability $p$, which we model as a Beta distribution with parameters \(\alpha\) and \(\beta\). \\

\[ P(p) = \frac{p^{\alpha - 1} (1-p)^{\beta - 1}}{B(\alpha, \beta)} \]
where \( B(\alpha, \beta) \) is the Beta function, which serves as a normalization constant to ensure that the total probability integrates to 1.\\


Now, we conduct a series of experiments to measure the bit error rate in the communication channel. Let \(k\) be the number of observed errors and \(n\) be the total number of bits transmitted in the experiments. The likelihood function for \(p\) given the observed data is given by the Binomial distribution:

\[ P(\text{data} | p) = \text{Binomial}(k | n, p) \]

Now, our goal is to find the Maximum A Posteriori (MAP) estimate for the bit error probability (\(p\)), given the observed data and the prior information. According to Bayes' Theorem, the posterior distribution is proportional to the product of the likelihood and the prior:

\[ P(p | \text{data}) \propto P(\text{data} | p) \cdot P(p) \]

For a Beta prior and a Binomial likelihood, the posterior distribution is also a Beta distribution. This posterior distribution is given by:

\[ P(p | \text{data}) = \text{Beta}(\alpha + k, \beta + n - k) \]

The likelihood function for \( p \) given the observed data (number of observed errors \( k \) and total number of bits transmitted \( n \)) is given by the Binomial distribution:

\[ P(\text{data} | p) = \binom{n}{k} p^k (1-p)^{n-k} \]


The posterior distribution is the product of the likelihood function and the prior distribution (ignoring the terms that simply scale the magnitude in the expression of proportionality):
\[ P(p | \text{data}) \propto p^k (1-p)^{n-k} \cdot p^{\alpha - 1} (1-p)^{\beta - 1} \]

This simplifies to:
\[ P(p | \text{data}) \propto p^{k + \alpha - 1} (1-p)^{n - k + \beta - 1} \]
The mode of this distribution, which is the MAP estimate, is found by taking the derivative of the log of the posterior with respect to \( p \), setting it to zero, and solving for \( p \). This yields:
\[ \hat{p}_\text{MAP} = \frac{k + \alpha - 1}{n + \alpha + \beta - 2} \]




Here, \(\hat{p}_{\text{MAP}}\) represents the MAP estimate of the bit error probability based on the observed data and the prior information.\\

A numerical example is provided in \href{https://www.ece.unb.ca/kengleha/Textbook/Code/MAP_example.py}{\texttt{MAP\_example.py}}, and the results shown in Figure \ref{fig:MAP_example}.

\begin{figure}[h!]
\centering
\includegraphics[scale=0.5]{figures/MAP_example.png}
\caption{MAP example: The prior and posterior distributions of the bit error probability \(p\) in a digital communication channel are shown. The data were simulated using a binomial distribution with a true bit error probability of \(p=0.2\). The prior distribution, modeled as a Beta distribution with parameters \(\alpha = 10000\) and \(\beta = 41000\), and the posterior distribution, obtained using a hybrid Beta-Binomial model, are shown. The vertical dashed line indicates the true bit error probability, while the green dashed line represents the MAP estimate of \(p\).}
\label{fig:MAP_example}
\end{figure}


\subsubsection*{Comparison between Bayesian and MAP Estimation}
Bayesian estimation and MAP estimation are both methods used in statistics for parameter estimation, and they both incorporate prior knowledge about the parameters. However, they differ in how they use this information:

\begin{enumerate}
    \item \emph{Estimate: } The MAP estimate is the mode (peak) of the posterior distribution, which is the value of the parameter that maximizes the posterior probability. On the other hand, the Bayesian estimate is typically the expected value (mean) of the posterior distribution, which takes into account all possible values of the parameter, weighted by their posterior probabilities.

    \item \emph{Uncertainty: } Bayesian estimation provides a full posterior distribution for the parameter, which gives a sense of the uncertainty or variability in the estimate. It tells us how plausible each parameter value is given the data. In contrast, MAP provides a single value without a direct measure of uncertainty.

    \item \emph{Influence of Prior: } In MAP estimation, the prior can have a strong influence on the estimate when there's little data, but its influence diminishes with more data. In Bayesian estimation, the prior always influences the estimate to some extent because it affects the shape of the entire posterior distribution.
\end{enumerate}


In summary, if you're interested in a point estimate and have a good idea of what the parameter should be, MAP might be a good choice. If you want to quantify uncertainty and incorporate prior information in a more nuanced way, Bayesian estimation would be more appropriate.

\section{Maximum Likelihood Estimation}
Consider estimating parameter \(\theta\) where its prior distribution \(f_\theta (\theta)\) is not known, and consequently cannot be used in MAP estimation. In this situation, we can use the MLE estimator which does not require information about the prior distribution. MLE can also be viewed as the MAP point estimator in which the a prior density is assumed uniform over parameter range.\\



The likelihood function, denoted as \( \mathcal{L}( x | \theta ) \), represents the probability of observing the given data \( x \) for different values of the model parameters \( \theta \).

\[
\mathcal{L}(x| \theta) = f_\mathrm{x|\theta}(x)
\]

For \(N\) i.i.d data samples \(\mathbf{X} = \{X_1, X_2, \ldots X_n \} \) this is 
\[
\mathcal{L}(\mathbf{x} | \theta ) = \prod_{n=1}^N f_\mathrm{x_n|\theta}(x_n)
\]

Due to mathematical convenience, the log-likelihood function (log-likelihood) is often used instead of the likelihood function.
 \[
 \ell(\mathbf{x} | \theta ) =  \sum_{n=1}^N \ln \left (f_\mathrm{x_n|\theta}(x_n) \right )\]

MLE seeks to find the parameter values that maximize the likelihood (or log-likelihood) function. Mathematically, it involves solving the optimization problem: 

\[ \hat{\theta}_{\text{MLE}} = \arg\max_{\theta} \ell(\theta | \mathbf{x}) 
\]

\textbf{Steps in MLE}\\

\begin{enumerate}
    \item \emph{Formulate Likelihood Function:} Write down the likelihood function based on the assumed statistical model and the observed data.
    \item \emph{Take Logarithm:} Simplify the likelihood function by taking the natural logarithm, yielding the log-likelihood.
    \item \emph{Differentiate and Set to Zero: } Differentiate the log-likelihood with respect to each parameter and set the derivatives equal to zero to find the critical points.
    \item \emph{Solve for Parameters:} Solve the system of equations to find the parameter values that maximize the log-likelihood.

\end{enumerate}

\begin{tcolorbox}
    \begin{exmp}{ML Estimation of the Mean of a Normal Distribution.}

    Let's consider a simple analytical example to illustrate MLE: estimating the mean of a normal distribution.

\end{exmp}
\end{tcolorbox}



Suppose you have a random sample \(X_1, X_2, \ldots, X_n\) from a normal distribution with an unknown mean \(\mu\) and a known standard deviation \(\sigma\). The likelihood function for this scenario is given by the product of the probability density functions of the normal distribution:

\[ \mathcal{L}( x_1, x_2, \ldots, x_n | \mu ) = \prod_{i=1}^{n} \frac{1}{\sqrt{2\pi}\sigma} e^{-\frac{(x_i - \mu)^2}{2\sigma^2}} \]

The log-likelihood function is:

\[ \ell( x_1, x_2, \ldots, x_n | \mu) = -\frac{n}{2} \log(2\pi\sigma^2) - \frac{1}{2\sigma^2} \sum_{i=1}^{n} (x_i - \mu)^2 \]

To find the maximum likelihood estimate \(\hat{\mu}_{\text{MLE}}\), we differentiate the log-likelihood with respect to \(\mu\) and set the result equal to zero:

\[ \frac{d\ell}{d\mu} = \frac{1}{\sigma^2} \sum_{i=1}^{n} (x_i - \hat{\mu}_{\text{MLE}}) = 0 \]

Solving this equation gives the MLE for \(\mu\):

\[ \hat{\mu}_{\text{MLE}} = \frac{1}{n} \sum_{i=1}^{n} x_i \]

This is the sample mean, which is a well-known result for the MLE of the mean in a normal distribution.\\

\begin{tcolorbox}
    \begin{exmp}{ML Estimation of Regression Parameters.}
    
    Let's consider a more complex example of MLE: estimating the parameters of a simple linear regression model.
    \end{exmp}
\end{tcolorbox}


Suppose we have a dataset \((x_1, y_1), (x_2, y_2), \ldots, (x_n, y_n)\) and assume that the relationship between \(x\) and \(y\) follows a simple linear regression model:

\[ y_i = \beta_0 + \beta_1 x_i + \epsilon_i \]

where \(\epsilon_i\) is a random error term. We assume that \(\epsilon_i \sim N(0, \sigma^2)\) with constant variance.

The likelihood function for the observed data is given by the product of the normal distributions:

\[ \mathcal{L}(x_1, y_1, \ldots, x_n, y_n | \beta_0, \beta_1, \sigma ) = \prod_{i=1}^{n} \frac{1}{\sqrt{2\pi\sigma^2}} e^{-\frac{(y_i - (\beta_0 + \beta_1 x_i))^2}{2\sigma^2}} \]

The log-likelihood function is:

\[ \ell(x_1, y_1, \ldots, x_n, y_n | \beta_0, \beta_1, \sigma) = -\frac{n}{2} \log(2\pi\sigma^2) - \frac{1}{2\sigma^2} \sum_{i=1}^{n} (y_i - (\beta_0 + \beta_1 x_i))^2 \]

To find the maximum likelihood estimates \(\hat{\beta}_0\), \(\hat{\beta}_1\), and \(\hat{\sigma}\), we might consider differentiating the log-likelihood with respect to the parameters and set the results equal to zero. The solutions involve solving a system of equations. In practice, this is very difficult if not impossible, and so numerical optimization methods are often used.\\

 
In routine \href{https://www.ece.unb.ca/kengleha/Textbook/Code/MLE_example.py}{\texttt{MLE\_example.py}} we generate synthetic data from a linear regression model, define the log-likelihood function, and use the `minimize` function from SciPy to find the maximum likelihood estimates. Figure  \ref{fig:MLE_example} shows the true regression line, observed data, and the MLE regression line.

\begin{figure}[h!]
\centering
\includegraphics[scale=0.4]{figures/MLE_regression.png}
\caption{The MLE estimate of the regression parameters \(\hat{\beta}_0\), \(\hat{\beta}_1\), and \(\hat{\sigma}\).}
\label{fig:MLE_example}
\end{figure}


\section{Minimum Mean Square Estimation}
Given a random variable \(X\), we wish to estimate the RV \(Y\) from \(X\) is the MMSE sense. Define
\[
\hat{Y} = g(X)
\]
where \(\hat{Y}\) is the MMSE estimate of \(Y\). In order to minimize the mean square error:

\[
\begin{aligned}
MSE &= E[ (Y - \hat{Y})^2]\\
&= \int_{-\infty}^{\infty} \int_{-\infty}^{\infty} (Y-g(X))^2 f_\mathrm{xy}
(x,y) \, \, dx \, dy \\
&= \int_{-\infty}^{\infty} \underbrace{\left [ \int_{-\infty}^{\infty} (Y-g(X))^2 f_\mathrm{y}(y|x) \, dy \right ]}_{\text{minimize this quantity}} f_\mathrm{x}(x) \, dx 
\end{aligned}
\]
Denote

\[
G = \int_{-\infty}^{\infty} (Y-g(x))^2 f_\mathrm{y}(y|x) \, dy
\]

Minimize with respect to \(\hat{Y}=g(X)\) using the chain rule. Note that \(\frac{d(y-\hat{Y})^2}{d \hat{Y}} = -2(y-\hat{Y})\).

\[
\frac{dG}{d\hat{Y}} = \int_{-\infty}^{\infty} -2y f_\mathrm{y}(y|x) \, dy + \hat{Y}\int_{-\infty}^{\infty}2 f_\mathrm{y}(y|x) \, dy = 0
\]

so that,
\[
2 \hat{Y} - 2 \int_{-\infty}^{\infty} y f_\mathrm{y}(y|x) \, dy = 0
\]

And finally,
\[
\boxed{\hat{Y} =  \int_{-\infty}^{\infty} y f_\mathrm{y}(y|x) \, dy = E[Y|X]}
\]

This is the minimum mean square solution, the conditional mean of the random variable  \(Y\) conditioned upon \(X\).

\subsection{Linear MMSE Estimation}

Let's derive the MMSE solution for the linear estimator \(\hat{Y} = aX + b\).\\

Given the Mean Squared Error (MSE) as:

\begin{align*}
\text{MSE} &= E[(Y - (\hat{Y}))^2] \\
&= E[Y^2 - 2aXY - 2bY + (\hat{Y})^2] \\
&= E[Y^2] - 2aE[XY] - 2bE[Y] + E[a^2X^2 + 2abX + b^2] \\
&= E[Y^2] - 2aE[XY] - 2bE[Y] + a^2E[X^2] + 2abE[X] + b^2 \\
\end{align*}

Now, take the partial derivatives with respect to \(a\) and \(b\) and set them equal to zero:


\begin{align*}
\frac{\partial \text{MSE}}{\partial a} &= -2E[XY] + 2aE[X^2] + 2bE[X] = 0 \\
\frac{\partial \text{MSE}}{\partial b} &= -2E[Y] + 2aE[X] + 2b = 0
\end{align*}

Solve these two equations simultaneously for \(a\) and \(b\). The solutions are:

\begin{align*}
a &= \frac{\text{Cov}(X, Y)}{\text{Var}(X)} \\
b &= E[Y] - aE[X]
\end{align*}

These are the values for \(a\) and \(b\), providing the MMSE estimator for \(Y\):

\[ \hat{Y}_{\text{MMSE}} = aX + b \]


\textbf{Remark:} If $X$ and $Y$ are uncorrelated, i.e., $\text{Cov}(X, Y) = 0$, then setting $a = 0$ solves this equation. In this case, we have $Y = g(x) = b$. The implication is that the guess for $Y$ does not include any information from $X$. Therefore, linear estimation cannot extract any information about $Y$ that can minimize the Mean Squared Error $MSE = E[(Y - b)^2]$.\\


\begin{tcolorbox}
    \begin{exmp}{Signal + Random Noise with MMSE}
    
Consider an observed signal $X(t)$ that consists of a `signal' $Z(t)$ which is a a zero-mean random process with variance $\sigma_z^2$, and additive zero-mean noise $W(t)$ with variance $\sigma_w^2$.

\[ X(t) = Z(t) + W(t) \]

We wish to find the linear MMSE solution \(Z_\text{MMSE}\) from the noisy signal \(X\) using the linear estimator

\[
\hat{Z} = aX + b
\]
    
    \end{exmp}
\end{tcolorbox}

We calculate $a$ using the formula:
\[ a = \frac{\text{Cov}(Z, X)}{\text{Var}(X)} \]

\[
\begin{aligned}
\text{Cov}(Z, X) & = E[(Z-\mu_z)(X-\mu_\mathrm{x})]\\
                & = E[ZX]\\
                & = E[Z(Z+W)]\\
                & = E[ZZ] + E[ZW]
\end{aligned}
\]

If we assume that the signal and the noise are uncorrelated, then

\[
\text{Cov}(Z, X) = \sigma_z^2
\]

Given that $\text{Var}(X) = \sigma_z^2 + \sigma_w^2$, this becomes
\[ a = \frac{\sigma_z^2}{\sigma_z^2 + \sigma_w^2} \]

Since all processes have zero mean, \(b=0\).\\

The linear MMSE estimate of the signal $Z$ is :
\[Z = \frac{\sigma_z^2}{\sigma_z^2 + \sigma_w^2} X(t) \]


The following Matlab code illustrates this example:

\begin{lstlisting}[language=Matlab, caption={MATLAB Code Example}]
% Parameters
sigma_z = 1;  % Variance of Z
sigma_w = 0.5;  % Variance of W
sigma_\mathrm{x} = 0.8;  % Variance of X

% Generate a time vector
t = 0:0.1:10;

% Generate the random process Z
Z = randn(size(t)) * sqrt(sigma_z);

% Generate white Gaussian noise W
W = randn(size(t)) * sqrt(sigma_w);

% Form the observed signal X
X = Z + W;

% Calculate the optimal coefficient A
A = sigma_z^2 / (sigma_z^2 + sigma_w^2);

% Compute the MMSE estimate of Z
Z_hat = A * X;

% Plot the signals
figure;
plot(t, Z, 'b', 'LineWidth', 2, 'DisplayName', 'Z(t)');
hold on;
plot(t, X, 'r', 'LineWidth', 1.5, 'DisplayName', 'X(t)');
plot(t, Z_hat, 'g--', 'LineWidth', 1.5, 'DisplayName', 'Estimated Z(t)');
legend();
xlabel('Time');
ylabel('Amplitude');
title('MMSE Estimation Example');
grid on;
\end{lstlisting}

\begin{figure}[h!]
\centering
\includegraphics[scale=0.4]{figures/MMSE_example.png}
\caption{The MMSE estimate of \(Z\) given noisy \(X\).}
\label{fig:MMSE_example}
\end{figure}



\subsection{Orthogonality Principle}

The orthogonality principle in LMMSE states that the error is orthogonal to the data:
\[
E[(Y-\hat{Y})X] = 0
\]

\bigskip

For example, for \(\hat{Y}=aX + b\) and without loss of generality, let \(\mu_\mathrm{x} = 0\) so that \(b = 0\).

\[
\begin{aligned}
    E[(Y-aX)X] & = E[YX] - E[aX^2]\\
                & = E[YX] -a E[X^2]
\end{aligned}
\]

We know that the LMMSE solution is:

\[
\begin{aligned}
a &= \frac{\text{Cov}(X, Y)}{\text{Var}(X)} \\
    & = \frac{E[XY]}{\sigma_\mathrm{x}^2}
\end{aligned}
\]

Therefore,
\[
E[(Y-aX)X] = E[YX] - \frac{E[XY]}{\sigma_\mathrm{x}^2}(\sigma_\mathrm{x}^2 + \mu_\mathrm{x}^2)
\]
since \(\mu_\mathrm{x}=0\) this becomes
\[
\begin{aligned}
    E[(Y-aX)X] & = E[YX] - \frac{E[XY]}{\sigma_\mathrm{x}^2}\sigma_\mathrm{x}^2 \\
                & = 0
\end{aligned}
\]

Therefore, the error is orthogonal to the data in the LMMSE.\\

The orthogonality of the error to the data is an extremely important result in estimation. In general, given an LMMSE estimate \( Y \) in terms of data \( \bar{X} \):

\[
Y = a_1 X_1 + a_2 X_2 + \ldots + a_n X_n
\]

Then, from the orthogonality principle:

\[
{E}[(Y - \hat{Y}) X_i] = 0 \quad \text{for all } i
\]



\begin{tcolorbox}
\begin{exmp}{Multivariate LMMSE}
Suppose we have a system where the observed data vector is \(X = [X_1, X_2, X_3]\) and you want to estimate a target variable \(Y\). We'll assume a linear relationship between \(X\) and \(Y\) with some noise.

The relationship is given by:

\[ Y = w_1 X_1 + w_2 X_2 + w_3 X_3 + \epsilon \]

where:
\begin{itemize}
    \item \(w_1, w_2, w_3\) are the weights or coefficients
    \item \(\epsilon\) is a random variable representing the noise
\end{itemize}


We want to find the weights \(w_1, w_2, w_3\) that minimize the mean square error.


\end{exmp}
\end{tcolorbox}

\textbf{Steps:}
\begin{enumerate}
    \item \textbf{Define the Linear Model}
   \[ \hat{Y} = w_1 X_1 + w_2 X_2 + w_3 X_3 \]
    \item \textbf{Define the Error}
   \[ \epsilon = Y - \hat{Y} \]

    \item \textbf{Minimize Mean Square Error (MSE)}
   \[ \text{MSE} = {E}[(Y - \hat{Y})^2] \]
   Minimize MSE with respect to \(w_1, w_2, w_3\) to find the optimal weights.

    \item \textbf{Optimal Weights (LMMSE Estimators)}
   \[ w_1^* = \frac{\text{Cov}(X_1, Y)}{\text{Var}(X_1)} \]
   \[ w_2^* = \frac{\text{Cov}(X_2, Y)}{\text{Var}(X_2)} \]
   \[ w_3^* = \frac{\text{Cov}(X_3, Y)}{\text{Var}(X_3)} \]
   
\end{enumerate}

These optimal weights minimize the mean square error, and the linear combination \(\hat{Y} = w_1^* X_1 + w_2^* X_2 + w_3^* X_3\) gives the linear MMSE estimate of \(Y\), shown in Figure \ref{fig:LMMSE_MV}. The code for this example is in \href{https://www.ece.unb.ca/kengleha/Textbook/Code/LMMSE_multivariate.py}{\texttt{LMMSE\_multivariate.py}}.

\begin{figure}[h!]
\centering
\includegraphics[scale=0.4]{figures/LMMSE_MV.png}
\caption{The linear estimate using \(\hat{Y} = w_1 X_1 + w_2 X_2 + w_3 X_3 \).}
\label{fig:LMMSE_MV}
\end{figure}

Note: This example assumes that the variance of the noise term \(\epsilon\) is known, and the covariance and variance terms are finite. In practical applications, you might need to estimate these values from data.





\section{Summary}

\textbf{Bayesian Estimation}

Bayesian estimation is a framework that incorporates prior knowledge about the parameters along with observed data to make statistical inferences. It involves updating prior beliefs using Bayes' theorem.
The posterior distribution, which is the distribution of the parameters given the observed data, is central to Bayesian estimation. The mean of the posterior distribution is a Bayesian point estimate, and this is often used as a Bayesian counterpart to MLE.\\

\textbf{Maximum A Posteriori (MAP) Estimation}

MAP estimation is a specific form of Bayesian estimation where the goal is to find the parameter values that maximize the posterior distribution. Mathematically, it involves maximizing the product of the likelihood and the prior.
MAP estimation can be seen as a compromise between MLE and Bayesian estimation, where it incorporates prior information while still providing a point estimate.\\

\textbf{Maximum Likelihood Estimation}

MLE is a classical frequentist approach to estimation. It seeks to find the parameter values that maximize the likelihood function, which measures how well the observed data fit the assumed statistical model.
MLE is often used when there is no prior information available about the parameters, and it provides point estimates for the parameters.\\




\textbf{Linear Minimum Mean Squared Error Estimation}

LMMSE estimation is a particular form of estimation that minimizes the mean squared error between the true parameter value and the estimated value, given the observed data.
LMMSE estimation is often used in the context of linear models and is particularly relevant when there is a linear relationship between the observed data and the parameters.\\

\textbf{Relationships}

\begin{itemize}
    \item MLE is a special case of Bayesian estimation where the prior distribution is assumed to be non-informative or uniform.
    \item MAP estimation incorporates prior information and reduces to MLE when the prior distribution is non-informative.
    \item LMMSE is not inherently Bayesian because it does not involve specifying prior distributions or updating beliefs based on observed data in the formal Bayesian sense. Instead, it focuses solely on minimizing the mean squared error without considering probabilistic models or uncertainties explicitly.
\end{itemize}

\chapter{Detection}

This chapter develops the fundamentals of detection theory. It follows the chapter on estimation theory as often plays a role in designing detection algorithms by estimating parameters of signals or noise distributions. \\

The basic components of a detection (or decision theory) problem are depicted in Figure \ref{fig:decision_components}.

\begin{figure}[h!]
\centering
\includegraphics[scale=0.3]{figures/decision_components.png}
\caption{The components of a decision theory problem.}
\label{fig:decision_components}
\end{figure}

The first component is a \emph{source} that generates a signal. In the simplest case, this signal is one of two choices. We can refer to these are \emph{hypotheses} and label them $H_0$ and $H_1$, which is a \emph{binary} problem. More generally, the signal might be one of $M$ possible choices which is referred to as an $M$-ary problem. Some typical sources include:
\begin{itemize}
    \item A digital communication system that transmits information by sending zeros ($H_0$) and ones ($H_1$).
    \item An electrocardiogram is examined to determine if a patient has had a heart attack. Here, a feature (or features) of the electrocardiogram will suggest the presence $H_1$ or absence $H_0$ of a heart attack. 
\end{itemize}

The second component is a probabilistic transition mechanism and the third is the observation space. The transition mechanism captures the dynamics of the system and therefore governs how the input will map to the observation space. Generally, with noisy or uncertain inputs (hypotheses) it may be viewed as a device that generates a point in observation space according to some probability law.



\section {Binary Detection}
\label{sec:Binary_Detection}

The binary detection problem often involves making a decision between two possible hypotheses in the presence of \emph{additive white noise}. This problem is often encountered in signal processing, communications, and statistical signal detection. The hypotheses are formulated as follows:

\begin{itemize}
    \item \textbf{Hypothesis \(H_0\):} The observed signal is purely noise (null hypothesis).
    \item \textbf{Hypothesis \(H_1\):} The observed signal contains a signal of interest embedded in white noise (alternative hypothesis).
\end{itemize}


The goal is to design a detector that can make a decision based on the observed signal to determine which hypothesis is more likely to be true. A Bayes detector is a type of detector that minimizes the probability of error by considering the probability of each hypothesis given the observed data.\\

The Bayes detector uses Bayes' rule to compute the posterior probabilities of the hypotheses given the observed data. The decision rule is then based on comparing these posterior probabilities. Mathematically, the Bayes detector can be described as follows:\\

\begin{enumerate}
    \item \textbf{Prior Probabilities:} Assign prior probabilities \(P(H_0)\) and \(P(H_1)\) to the hypotheses, indicating the initial belief in the likelihood of each hypothesis before observing the data.
    \item \textbf{Likelihood Function:} Specify the likelihood function \(P(\mathbf{X} | H_0)\) and \(P(\mathbf{X} | H_1)\), which describe the probability distribution of the observed data \(\mathbf{X}\) given each hypothesis. In general, \(\mathbf{X}\) is a vector of data samples. In the case of white noise, this distribution is typically a normal distribution.
    \item \textbf{Posterior Probabilities:} Use Bayes' rule to compute the posterior probabilities:
   \[ P(H_0 | \mathbf{X}) = \frac{P(\mathbf{X} | H_0) \cdot P(H_0)}{P(\mathbf{X})} \]
   \[ P(H_1 | \mathbf{X}) = \frac{P(\mathbf{X} | H_1) \cdot P(H_1)}{P(\mathbf{X})} \]
   where \(P(\mathbf{X})\) is the probability of the data \(\mathbf{X}\).
    \item \textbf{Decision Rule:} Make a decision by comparing the posterior probabilities:
   \begin{itemize}
       \item If \(P(H_0 | \mathbf{X}) > P(H_1 | \mathbf{X})\), decide in favor of \(H_0\).
       \item If \(P(H_1 | \mathbf{X}) > P(H_0 | \mathbf{X})\), decide in favor of \(H_1\).
      \end{itemize}
   This is equivalent to choosing \(H_1\) \emph{iff} \(P(H_1)P(\mathbf{X}|H_1) > P(H_0)P(\mathbf{X}|H_0)\) and \(H_0\) otherwise.

\end{enumerate}

\subsection{{The Bayesian Binary Detector}}


Consider the components of the  Bayesian binary detector:

\begin{itemize}
    \item The prior probabilities: \( P(H_0) \), \( P(H_1) \)
    \item \( \mathbf{R} \): The received signal; in general, a data vector.
    \item \( Z_0 \): The region for which we conclude \( H_0 \) is true when we receive \( \mathbf{R} \) in the observation space.
    \item \( Z_1 \): The region for which we conclude \( H_1 \) is true when we receive \( \mathbf{R} \) in the observation space.
    \item The channel transition probabilities (or likelihood functions). These are $f(\mathbf{R} | H_0)$ and $f(\mathbf{R} | H_1)$, or equivalently, by their probabilities $P(\mathbf{R} | H_0)$ and $P(\mathbf{R} | H_1)$.
\end{itemize}

The Bayes' binary detector determines a decision rule that divides the observation space into two parts, $Z_0$ and $Z_1$ as shown in Figure \ref{fig:Z0_Z1_area}. Whenever an observation falls in $Z_0$ we say $H_0$ and whenever an observation falls in $Z_1$ we say $H_1$.\\

\begin{figure}[h!]
\centering
\includegraphics[scale=0.3]{figures/Z0_Z1_area.png}
\caption{The decision space of $Z_0$ and $Z_1$.}
\label{fig:Z0_Z1_area}
\end{figure}

The error probability is given by:

\[ P(E) = P(E | H_0)P(H_0) + P(E | H_1)P(H_1) \]

Let's express \( P(E | H_0) \) and \( P(E | H_1) \) in terms of the likelihoods \( f(\mathbf{R} | H_1) \) and \(f(\mathbf{R} | H_0) \):



\[ P(E | H_0) = \int_{Z_1} f(\mathbf{R} | H_0) \, dR \]
\[ P(E | H_1) = \int_{Z_0} f(\mathbf{R} | H_1) \, dR \]
Note that if  \(\mathbf{R} = \{r_1, r_2, \ldots r_N\}\), then \( f(\mathbf{R} | H_i) \) is an \(N^{th}\) order joint density function.\\\\
Consider the decision rule which assigns the regions \( Z_0 \) and \( Z_1 \). The probability of error is given by:

\[ P(E) = P(H_0) \int_{Z_1} f(\mathbf{R} | H_0) \, dR + P(H_1) \int_{Z_0} f(\mathbf{R} | H_1) \, dR \]



This can also be written as:

\[ P(E) = P(H_0) \left \{ 1- \int_{Z_0} f(\mathbf{R} | H_0) \, dR  \right \} + P(H_1) \int_{Z_0} f(\mathbf{R} | H_1) \, dR \]

Which is:

\[ P(E) = 
\begin{cases}
    P(H_0) +  \int_{Z_0} \left \{ P(H_1) f(\mathbf{R} | H_1) 
 - P(H_0) f(\mathbf{R} | H_0) \right \} \, dR  \\
 \text{or} \\
 P(H_1) +  \int_{Z_1} \left \{ P(H_0) f(\mathbf{R} | H_0) 
 - P(H_1) f(\mathbf{R} | H_1) \right \} \, dR
\end{cases}
\]

A decision rule to minimize \(P(E)\) is to assign \(\mathbf{R}\) to \(Z_0\) \emph{iff}:
\[
\int_{Z_0} P(H_1)f(\mathbf{R} | H_1) < \int_{Z_0} P(H_0)f(\mathbf{R} | H_0)
\]

which can also be written as:
\[
P(H_1)P(\mathbf{R}|H_1) < P(H_0)P(\mathbf{R}|H_0)
\]
This can be written to establish an assignment to both $H_0$ and $H_1$ as 

\[
\frac{P(\mathbf{R}|H_1)}{P(\mathbf{R}|H_0)}  
\begin{tabular}{c @{\hspace{2\tabcolsep}} *{5}{c}}

     \(H_1\) \\
     $>$ \\
    $<$ \\
    \(H_0\)
\end{tabular}
\frac{P(H_0)}{P(H_1)} 
\]


We refer to the left side as the \emph{likelihood ratio} (for obvious reasons), denoted as 
\[
\Lambda(\mathbf{R}) = \frac{P(\mathbf{R}|H_1)}{P(\mathbf{R}|H_0)}  
\begin{tabular}{c @{\hspace{2\tabcolsep}} *{5}{c}}

     \(H_1\) \\
     $>$ \\
    $<$ \\
    \(H_0\)
\end{tabular}
\frac{P(H_0)}{P(H_1)} 
\]


This conclusion aligns with the idea that the likelihood ratio test is optimal for minimizing the probability of error in a Bayesian framework. It's also worth noting that this result is consistent with the Neyman-Pearson criterion for hypothesis testing.

\subsection{Binary Detection in  Gaussian White Noise}

The binary detection problem is the task of deciding whether a signal \(\mathbf{S}\) in the presence of white Gaussian noise is present or not in a given time interval \([0,T]\). This signal \(\mathbf{S}\) may be any waveform such as a sinusoidal pulse or a binary (ones and zeros) waveform.  For simplicity, we consider a random binary waveform, as  illustrated in Figure \ref{fig:Binary_Wave_Noisey}.

\begin{figure}[h!]
\centering
\includegraphics[scale=0.7]{figures/Binary_Wave_Noisey.png}
\caption{Sampling a random binary waveform. Here, the binary waveform has a period of \(T=100 ms\) and is sampled every \(10ms.\)}
\label{fig:Binary_Wave_Noisey}
\end{figure}

We denote the received binary waveform as \(\mathbf{R}\) received signal, consisting of \(N\) samples per period (in this case, \(N=10\). This received waveform contains the noise samples \(\mathbf{N} \sim \mathcal{N}(0, \sigma_n^2)\) plus the signal \(\mathbf{S}\) (a pulse of amplitude = 1) when it is present. Consequently, the hypothesis testing problem is:
\[
\begin{aligned}
& H_0: \mathbf{R} = \mathbf{N} \\
& H_1: \mathbf{R} = \mathbf{S} + \mathbf{N}
\end{aligned}
\]

The detector decision rule is:
\[
\Lambda(\mathbf{R}) = \frac{P(\mathbf{R}|H_1)}{P(\mathbf{R}|H_0)}  
\begin{tabular}{c @{\hspace{2\tabcolsep}} *{5}{c}}

     \(H_1\) \\
     $>$ \\
    $<$ \\
    \(H_0\)
\end{tabular}
\frac{P(H_0)}{P(H_1)} 
\]

We now need to evaluate the likelihood functions, \(P(\mathbf{R}|H_0)\) and \(P(\mathbf{R}|H_1)\).  Since we know that the noise samples are \(\sim \mathcal{N}(0, \sigma_n^2)\) and \emph{i.i.d}:

\[
P(\mathbf{R}|H_0) = \prod_{i=1}^N \left [ \frac{1}{\sqrt{2 \pi} \sigma_n} \exp \left ({-\frac{r_i^2}{2 \sigma_n^2}} \right )\right] =  \frac{1}{(2 \pi \sigma^2_n)^{N/2}} \exp \left ({-\sum_{i=1}^N \frac{r_i^2}{2 \sigma_n^2}} \right )
\]

\[
P(\mathbf{R}|H_1) = \prod_{i=1}^N \left [ \frac{1}{\sqrt{2 \pi} \sigma_n} \exp \left ({-\frac{(r_i-s_i)^2}{2 \sigma_n^2}} \right )\right] =  \frac{1}{(2 \pi \sigma^2_n)^{N/2}} \exp \left ({-\sum_{i=1}^N \frac{(r_i-s_i)^2}{2 \sigma_n^2}} \right )
\]

For \(P(\mathbf{R}|H_1)\) we make the distribution zero mean by removing the signal portion of the waveform: \(\mathbf{N=R-S}\). \\

Therefore, the likelihood ratio is:
\[
\Lambda(\mathbf{R}) = \frac{\exp \left ({-\sum_{i=1}^N \frac{(r_i-s_i)^2}{2 \sigma_n^2}} \right )}{\exp \left ({-\sum_{i=1}^N \frac{r_i^2}{2 \sigma_n^2}} \right )}  
\]

Taking the natural logarithm:
\[
\ln \{  \Lambda(\mathbf{R}) \} = \sum_{i=1}^N 2\frac{r_i s_i}{2 \sigma_n^2} - \sum_{i=1}^N \frac{(r_i-s_i)^2}{2 \sigma_n^2}  
\begin{tabular}{c @{\hspace{2\tabcolsep}} *{5}{c}}

     \(H_1\) \\
     $>$ \\
    $<$ \\
    \(H_0\)
\end{tabular}
\ln \left \{\frac{P(H_0)}{P(H_1)} \right \} 
\]

This can be rearranged as:
\[
\underbrace{\sum_{i=1}^N r_i s_i}_{\text{matched filter}}   
\begin{tabular}{c @{\hspace{2\tabcolsep}} *{5}{c}}

     \(H_1\) \\
     $>$ \\
    $<$ \\
    \(H_0\)
\end{tabular}
\underbrace{\sigma_n^2 \ln \left \{\frac{P(H_0)}{P(H_1)} \right \}}_{\text{a priori probabilities}}  - \underbrace{\frac{1}{2}\sum_{i=1}^N s_i^2}_{\frac{1}{2}\text{ signal energy}}
\]

The left-hand side is a discrete-time \emph{matched filter}, with the decision threshold determined by the a priori probabilities and the signal energy.\\

Note: if the noise is not Gaussian, then the matched filter detector will not yield the minimum \(P(E)\), but it will yield the best SNR. \\ 

\begin{tcolorbox}
\begin{exmp}{Numerical Example of Binary Detector.}

    We now consider a random binary waveform with \(P(H_0) = P(H_1) = 0.5\). This corresponds to a log likelihood ratio of \(\Lambda(\mathbf{R}) = 0\). The waveform is subject to Gaussian white noise with $\sim \mathcal{N} (0, \sigma_n)$. The noisy waveform is sampled 10 times per period at the detector.

    \end{exmp}
\end{tcolorbox}

    The detector is implemented in \href{https://www.ece.unb.ca/kengleha/Textbook/Code/BinaryDetectorExample.py}{\texttt{BinaryDetectorExample.py}}. First, consider \(\sigma_n=0.1\). First, consider the likelihood functions \(P(\mathbf{R}|H_0)\) and \(P(\mathbf{R}|H_1)\), which are plotted in Figure \ref{fig:BinaryDetectorLikelihoods} (top). For this relatively low value of noise variance, it is clear that little overlap exists between a transmitted \(H_0\) or \(H_1\). Indeed, this can be observed in the samples from each of samples under \(H_0\) or \(H_1\), shown in the bottom plot of \ref{fig:BinaryDetectorLikelihoods}. \\
    

\begin{figure}[h!]
\centering
\includegraphics[scale=0.4]{figures/BinaryDetectorLikelihoods.png}
\includegraphics[scale=0.4]{figures/BinaryDetectorSamples.png}
\caption{A binary detector with noise variance \(\sigma_n=0.1\). Top: The likelihood functions of the binary detector. Bottom: samples of the noisy signal during transmission of \(H_0\) and \(H_1\).}
\label{fig:BinaryDetectorLikelihoods}
\end{figure}

In this case, we have:\\

\begin{tabular}{c @{\hspace{2\tabcolsep}} *{5}{c}}
 - Log-Likelihood Ratio (H0): & -13.71 & (choose $H_0$) \\
    - Log-Likelihood Ratio (H1): & 77.61 & (choose $H_1$) 
     
\end{tabular}

These are safely below and above the detection threshold of 0.\\

Increasing the noise variance to \(\sigma_n=1.0\) depicts a more challenging situation for the detector, illustrated in Figure \ref{fig:BinaryDetectorLikelihoods2}. In this case, we have:\\

\begin{tabular}{c @{\hspace{2\tabcolsep}} *{5}{c}}
 - Log-Likelihood Ratio (H0): & -6.43 & (choose $H_0$) \\
    - Log-Likelihood Ratio (H1): & 6.60 & (choose $H_1$) 
     
\end{tabular}
   

Even in the presence of substantial noise, and outliers that would be clearly misclassified on their own, the inclusion of 10 samples in the likelihood calculation builds in robustness. 

\begin{figure}[h!]
\centering
\includegraphics[scale=0.4]{figures/BinaryDetectorLikelihoods2.png}
\includegraphics[scale=0.4]{figures/BinaryDetectorSamples2.png}
\caption{A binary detector with noise variance \(\sigma_n=1.0\). Top: The likelihood functions of the binary detector. Bottom: samples of the noisy signal during transmission of \(H_0\) and \(H_1\).}
\label{fig:BinaryDetectorLikelihoods2}
\end{figure}

\subsection{Detector Performance Measures}

In the context of binary hypothesis testing, several performance measures are used to assess the performance of a detector. Consider first the possible outcomes of a detector:

\begin{itemize}
    \item \textbf{True Negative} (\emph{TN}): The outcome is correctly identified as negative (belonging to \(H_0\)) by the detector when it is actually negative (belonging to \(H_0\)). 
    \item \textbf{True {Positive}} (\emph{TP}): The outcome is correctly identified as positive (belonging to \(H_1\)) by the detector when it is actually positive (belonging to \(H_1\)).
    \item \textbf{False Negative} (\emph{FN}): The outcome is incorrectly identified as negative (belonging to \(H_0\)) by the detector when it is actually positive (belonging to \(H_1\)).
    \item \textbf{False Positive} (\emph{FP}): The outcome is correctly identified as positive (belonging to \(H_1\)) by the detector when it is actually negative (belonging to \(H_0\)). 

\end{itemize}

Here,  are some of the key performance measures:

\begin{enumerate}

\item \textbf{True Negative Rate (TNR) or Specificity}
\begin{itemize}
    \item Measures the ability of the detector to correctly determine the absence of the signal when it is not present.
    \item \(P(\text{Decision} = H_0 | H_o \text{ is true})\)
    \item \(TNR = \frac{{TN}}{{TN + FP}}\)
\end{itemize}

\item \textbf{True Positive Rate (TPR) or Sensitivity}
\begin{itemize}
    \item  Measures the probability of the detector correctly detecting the presence of the signal when it is present.
    \item \(P(\text{Decision} = H_1 | H_1 \text{ is true})\)
    \item \(TPR = \frac{{FP}}{{FP + TN}}\)
\end{itemize}

\item \textbf{False Negative Rate (FNR) or Probability of False Negative (\(P_{FN}\))}
\begin{itemize}
    \item  Measures the probability of the detector failing to detect the signal when it is actually present.
    \item \(P(\text{Decision} = H_0 | H_1 \text{ is true})\)
    \item \(FNR = P_{FN} =  \frac{{FN}}{{FN + TP}}\)
\end{itemize}


\item \textbf{False Positive Rate (FPR) or Probability False Alarm (\(P_{FA}\))}
\begin{itemize}
    \item  Measures the probability of the detector incorrectly detecting the presence of the signal when it is not present.
    \item \(P(\text{Decision} = H_1 | H_0 \text{ is true})\).
    \item \(FPR = P_{FA} =  \frac{{FP}}{{FP + TN}}\)
\end{itemize}





\item \textbf{Positive Predictive Value (PPV) or Precision }
\begin{itemize}
    \item Measures the proportion of true positive decisions among all positive decisions made by the detector.
    \item \(PPV = \frac{{TP}}{{TP + FP}}\)
\end{itemize}


\item \textbf{F1 Score}
\begin{itemize}
    \item The F1 score is the harmonic mean of precision and sensitivity. It provides a balanced measure of a detector's accuracy when dealing with imbalanced datasets.
   \item \(F1 \text{ Score} = \frac{2}{PPV^{-1} + TPR^{-1}} = \frac{2 TP}{2 TP + FP + FN}\)
\end{itemize}

\end{enumerate}

\bigskip

\begin{exmp}{Binary Detector Metrics}
    \begin{tcolorbox}
        Suppose we have a binary detector that detects a digital signal (0 or 1) with additive noise. The detector outputs a probability of the signal being 1. We can use a threshold (e.g., 0.5) to make a binary decision (0 or 1). Show the effect of the noise noise on the binary detector metrics.  
    \end{tcolorbox}
\end{exmp}


The result is shown in Figure \ref{fig:BinaryDetectorMetrics}. The code \href{https://www.ece.unb.ca/kengleha/Textbook/Code/BinaryDetectorMetrics.py}{\texttt{BinaryDetectorMetrics.py}} implements this example.


\begin{figure}[h!]
\centering
\includegraphics[scale=0.7]{figures/BinaryDetectorMetrics.png}
\caption{The effect of noise level on binary detector metrics.}
\label{fig:BinaryDetectorMetrics}
\end{figure}

\bigskip

As the noise becomes much larger than the signal, the detector's output becomes essentially random, and the detector is equally likely to detect a TP, TN, FP or FN (if the probability of transmitting true and false are equal). Consequently, the values of each of the metrics converges to 0.5, indicating a complete loss of signal detection capability. 




\chapter{Optimal Filtering}

Optimal filtering is a fundamental concept in signal processing and control theory, aiming to estimate the state of a dynamic system or the parameters of a signal from noisy or uncertain observations. It encompasses a diverse range of techniques designed to extract useful information from observed data while minimizing estimation errors or achieving certain performance criteria. \\

At the heart of optimal filtering lies the principle of optimality, which entails selecting the most suitable estimation algorithm or filter to minimize some measure of error or cost. The choice of optimality criterion depends on the specific application and the desired performance metrics. Common optimality criteria include minimizing mean squared error, maximizing likelihood, minimizing Bayesian risk, or achieving certain performance trade-offs.\\

An important optimal filtering technique is the Wiener filter, which is used for linear estimation problems with stationary stochastic inputs. The Wiener filter minimizes mean squared error by exploiting the statistical properties of the input signal and noise, providing an optimal linear estimate of the desired signal.\\

One of the most well-known and widely used optimal filtering techniques is the Kalman filter. The Kalman filter is particularly powerful for estimating the state of linear dynamic systems subject to Gaussian noise. It recursively updates the state estimate based on a dynamic model of the system and noisy observations, leveraging both past estimates and current measurements to produce an optimal estimate of the system state. \\




\section{Linear Time-Invariant Systems with Stochastic Inputs}

Understanding Linear Time-Invariant (LTI) systems with stochastic inputs is essential for study of optimal filtering. This is because LTI systems are prevalent in engineering applications due to their simplicity and analytical tractability, and many real-world signals are stochastic in nature.\\

When dealing with stochastic inputs, the theory of LTI systems can be extended to include random processes or stochastic signals. A block diagram of an LTI system \(h(t)\) with input process \(X(t)\) and output \(Y(t)\) is depicted in Figure \ref{fig:LTI_system}.

\begin{figure}[h!]
\centering
\includegraphics[scale=0.4]{figures/LTI_system.png}
\caption{A generalized block diagram of a linear time-invariant system.}
\label{fig:LTI_system}
\end{figure}

LTI systems with stochastic inputs have the following characteristics:

\begin{itemize}
    \item \textbf{Linearity:} An LTI system is linear if it satisfies the superposition principle. For stochastic inputs, linearity implies that the system responds to the sum of individual input realizations in the same way it would respond to the sum of the individual responses to each input.
    \item \textbf{Time-invariance:} A system is time-invariant if its behavior or response does not change with time. In the context of stochastic inputs, this means that the system's response characteristics remain constant over time, regardless of when the input is applied.
    \item \textbf{Convolution operation:} The response of an LTI system to a stochastic input is often characterized by the convolution operation. For instance, if \(X(t)\) is the input stochastic process and \(h(t)\) is the impulse response of the LTI system, then the output \(Y(t)\) is given by \(Y(t) = X(t) * h(t)\), where '*' denotes convolution.
    \item \textbf{Stationarity:} When dealing with stochastic processes, assumptions about the stationarity of the input process become important. Wide-sense stationarity or strict-sense stationarity conditions are generally assumed.
\end{itemize}

With stochastic inputs, an LTI system  may be thought of as providing a one-to-one mapping between corresponding sample functions, as depicted in Figure \ref{fig:LTI_mapping}.\\

\begin{figure}[h!]
\centering
\includegraphics[scale=0.5]{figures/LTI_mapping.png}
\caption{One-to-one mapping of input to output sample functions in an LTI system.}
\label{fig:LTI_mapping}
\end{figure}



\subsection{Cross-correlation and Auto-correlation}
Now, we turn to the first and second-order moments and PSD of the output process \(Y(t)\) in terms of \(h(t)\) and the moments of \(X(t)\). \\
\[
\begin{split}
    Y(t) &  = \int_{-\infty}^\infty X(\upsilon)h(t-\upsilon) d\upsilon\\
    & = \int_{-\infty}^\infty X(t-\upsilon)h(\upsilon) d\upsilon
\end{split}
\]

\[
    Y(t+\tau)  = \int_{-\infty}^\infty X(t+\tau-\upsilon)h(\upsilon) d\upsilon
\]

Consider first the cross-correlation function:
\[
\begin{split}
    \phi_\mathrm{xy}(\tau) & = E [ X(t) Y(t+ \tau)]\\
    & = E \left [ X(t) \int_{-\infty}^\infty X(t+\tau-\upsilon)h(\upsilon) d\upsilon \right ] \\
    & = \int_{-\infty}^\infty E \left [ X(t)  X(t+\tau-\upsilon) \right ] h(\upsilon) d\upsilon  \\
    & = \int_{-\infty}^\infty \phi_\mathrm{xx}(\tau-\upsilon) h(\upsilon) d \upsilon
\end{split}
\]

Therefore:
\[
\boxed{\phi_\mathrm{xy}(\tau) = \phi_\mathrm{xx}(\tau) * h(\tau)}
\]
which is a convolution with \(h(t)\).\\

Now the autocorrelation function:
\[
\begin{split}
    \phi_\mathrm{yy}(\tau) & = E [ Y(t) Y(t+ \tau)]\\
    & = E \left [ \int_{-\infty}^\infty X(t-u)h(u) du \int_{-\infty}^\infty X(t+\tau-\upsilon)h(\upsilon) d\upsilon \right ] \\
    & = \int_{-\infty}^\infty \int_{-\infty}^\infty E \left [  X(t-u) X(t+\tau-\upsilon) \right ] h(u)   h(\upsilon) du d\upsilon  \\
    & = \int_{-\infty}^\infty \phi_\mathrm{xx}(\tau-\upsilon+u) h(u) h(\upsilon) d \upsilon
\end{split}
\]

Therefore:
\[
\boxed{\phi_\mathrm{yy}(\tau) = \phi_\mathrm{xx}(\tau) * h(\tau) * h(-\tau)}
\]

which is a double convolution with \(h(t)\).\\

\subsection{Cross- and Auto-Spectral Density}
Consider the the cross-spectral density:
\[
\begin{split}
    \Phi_\mathrm{xy}(f) & = \mathcal{F} \left \{ \phi_\mathrm{xy}(\tau)\right \}\\
                & = \int_{-\infty}^\infty \phi_\mathrm{xy}(\tau) e^{-j2\pi f \tau} d \tau\\
                & = \int_{-\infty}^\infty \left [ \int_{-\infty}^\infty \phi_\mathrm{xx}(\tau - \upsilon) h(\upsilon) d \upsilon \right ] e^{-j2\pi f \tau} d \tau\ \\
                & = \int_{-\infty}^\infty \left [ \int_{-\infty}^\infty \phi_\mathrm{xx}(\tau) e^{-j2\pi f \tau} e^{-j2\pi f \upsilon} d \tau\ \right ] h(\upsilon) d \upsilon \\
                 & = \int_{-\infty}^\infty \left [ \Phi_\mathrm{xx}(f) e^{-j2\pi f \upsilon} \, d f \right ] h(\upsilon) d \upsilon \\
                 & = \Phi_\mathrm{xx}(f)  \int_{-\infty}^\infty  h(\upsilon) e^{-j2\pi f \upsilon} d \upsilon
\end{split}
\]


Therefore:
\[
\boxed{\Phi_\mathrm{xy}(f)  = \Phi_\mathrm{xx}(f) H(f)}
\]

which is the Fourier transform pair of the cross-correlation \(\phi_\mathrm{xy}(\tau)  = \phi_\mathrm{xx}(\tau) * h(\tau)\).\\

Now consider the the auto-spectral density:
\[
\begin{split}
    \Phi_\mathrm{xy}(f) & = \mathcal{F} \left \{ \phi_\mathrm{yy}(\tau)\right \}\\
                & = \int_{-\infty}^\infty \phi_\mathrm{yy}(\tau) e^{-j2\pi f \tau} d \tau\\
                & = \int_{-\infty}^\infty \left [ \int_{-\infty}^\infty \phi_\mathrm{xx}(\tau - \upsilon +u) h(u) h(\upsilon) du d \upsilon \right ] e^{-j2\pi f \tau} d \tau\ \\
\end{split}
\]

Using the same methods as above:
\[
\Phi_\mathrm{yy}(f)  = \Phi_\mathrm{xx}(f) H(f) H^{*}(f)
\]

\[
\boxed{\Phi_\mathrm{yy}(f)  = \Phi_\mathrm{xx}(f) \left |H(f) \right |^2}
\]

\subsection{The Mean}

\[
\begin{split}
    E[Y(t)] & = E \left [ \int_{-\infty}^\infty X(t-\upsilon) h(\upsilon) d \upsilon \right ]\\
                & =  \int_{-\infty}^\infty E \left [ X(t-\upsilon) \right ] h(\upsilon) d \upsilon \\
                & =  \mu_\mathrm{x} \int_{-\infty}^\infty  h(\upsilon) d \upsilon
\end{split}
\]

\[
\boxed{\mu_\mathrm{y}  = \mu_\mathrm{x} H(0)}
\]

where \(H(0)\) is the DC gain of \(H(f)\). \\


\textbf{Note:} If \(X(t)\) is WSS, then \(Y(t)\) is also WSS.\\

Observation:
\[
\begin{split}
    Y(t) & = \int_{-\infty}^{\infty} X(t) h(t-u) du\\
        & \approx \sum_{i=-\infty}^\infty X(i \Delta) h(t-i \Delta) \Delta
\end{split}
\]

 If \(h(t)\) is lowpass with respect to \(X(t)\),

\[
Y(t) \approx h(t) \sum_{i=-\infty}^\infty X(i \Delta)  \Delta
\]

If the \(X(i \Delta)\) are roughly independent, the output will be Gaussian (central limit theorem).  

\section{Shaping Filters}

Given a process \(X(t)\) with some \(\Phi_\mathrm{xx}(f)\), generate a process
\(Y(t)\) with some \(\Phi_\mathrm{yy}(f)\).\\

Recall that:

\[ Y(t) = h(t) \ast X(t) \]

and 
\[ \Phi_\mathrm{yy}(f) = |H(f)|^2 \Phi_\mathrm{xx}(f) \]

so that
\[ |H(f)|^2 = \frac {\Phi_\mathrm{yy}(f)} {\Phi_\mathrm{xx}(f)} \]

\begin{tcolorbox}
    
\begin{exmp}{Shaping a white noise process.}\\
Suppose \(X(t)\) is white with \(\Phi_\mathrm{xx}(f) = K\) and we want

\[
\Phi_\mathrm{yy}(\omega) = \frac{1}{\omega^2 + \alpha^2}
\]

Find the required shaping filter \(h(t)\).

\end{exmp}
\end{tcolorbox}

Therefore,
\[ 
\begin{split}
    |H(f)|^2    & =   \frac{\Phi_\mathrm{yy}(f)}{\Phi_\mathrm{xx}(f) } \\
                & = \frac{1}{K(\omega^2 + \alpha^2)}
\end{split}
\]

To determine the transfer function, we  factor:
\[
|H(f)|^2 = \frac{1}{\sqrt{K}(\alpha+j \omega)} \cdot \frac{1}{\sqrt{K}(\alpha-j \omega)}
\]

Taking the causal part:
\[
H(\omega) = \frac{1}{\sqrt{K}(\alpha+j \omega)}
\]

Therefore:
\[
h(t) = \frac{1}{\sqrt{K}}e^{-\alpha t} u(t)
\]

which is a $1^{st}$ order LPF, depicted in Figure \ref{fig:LPF}.

\begin{figure}[h!]
\centering
\includegraphics[scale=0.6]{figures/LPF.png}
\caption{The impulse response of a first-order LPF \(h(t) = \frac{1}{\sqrt{K}}e^{-\alpha t} u(t)\) for \(K = 2.0 \) and \(\alpha=0.5\).}
\label{fig:LPF}
\end{figure}


\section{Poles and Zeros of Random Spectra}
In the context of random processes, the terms "poles" and "zeros" are often used in relation to the power spectral density.\\

\textbf{Poles}
\begin{itemize} 
    \item Poles are locations where the denominator of PSD becomes zero (and the PSD approaches infinity).
    \item The presence of poles in the spectrum indicates the characteristic frequencies at which the process tends to oscillate or exhibit significant energy.
\end{itemize}

\textbf{Zeros}
\begin{itemize}
    \item Zeros are locations where the numerator of the PSD becomes zero (and the the PSD becomes zero).
    \item The presence of zeros in the spectrum indicates frequencies at which certain components of the signal are absent or have minimal contribution.
\end{itemize}

Since the PSD of a random spectrum real and even, the poles and zeros must occur in groups of 4 or 2, as illustrated in Figure \ref{fig:poles_zeros}. \\

\begin{figure}[h!]
\centering
\includegraphics[scale=0.4]{figures/Poles_Zeros.png}
\caption{A typical pole-zero plot of the PSD of a random process.}
\label{fig:poles_zeros}
\end{figure}

For any \(\Phi_\mathrm{xx}(\omega)\) we can factor, grouping top-half plane (THP)  poles and zeros and lower-half plane (LHP) poles and zeros.
\[
\Phi_\mathrm{xx}(\omega) = \underbrace {\Phi_\mathrm{xx}^+(\omega)}_{\text{THP}} \underbrace {\Phi_\mathrm{xx}^-(\omega)}_{\text{LHP}}
\]

\begin{tcolorbox}
    \begin{exmp}{First-order lowpass system.}
    
Suppose that \(\phi_\mathrm{xx}(\tau) = e^{-a |\tau|}\) as illustrated in Figure \ref{fig:LP_phi}. Find $\Phi_\mathrm{xx}(\omega)$.
\end{exmp}
\end{tcolorbox}

\begin{figure}[h!]
\centering
\includegraphics[scale=0.4]{figures/LP_phi.png}
\caption{Autocorrelation function \(\phi_\mathrm{xx}(\tau) = e^{-a |\tau|}\).}
\label{fig:LP_phi}
\end{figure}



The Fourier transform of $\phi_\mathrm{xx}(\tau)$, denoted as $\Phi(\omega)$, is given by
\[
\Phi_\mathrm{xx}(\omega) = \int_{-\infty}^{\infty} \phi_\mathrm{xx}(\tau) e^{-j\omega\tau} \, d\tau.
\]

Substituting $\phi_\mathrm{xx}(\tau) = e^{-a|\tau|}$ into the integral, we get
\[
\begin{aligned}
\Phi_\mathrm{xx}(\omega) &= \int_{-\infty}^{\infty} e^{-a|\tau|} e^{-j\omega\tau} \, d\tau \\
&= \underbrace{\int_{-\infty}^{0} e^{(a-j\omega)\tau} \, d\tau}_{\text{-ve time}} + \underbrace{\int_{0}^{\infty} e^{(-a-j\omega)\tau} \, d\tau} _{\text{+ve time}}
\end{aligned}
\]

Integrating each term separately, we have
\[
\begin{aligned}
\Phi_\mathrm{xx}(\omega) &= \left. \frac{e^{(a-j\omega)\tau}}{a-j\omega} \right|_{-\infty}^{0} + \left. \frac{e^{(-a-j\omega)\tau}}{-a-j\omega} \right|_{0}^{\infty} \\
&= \underbrace{\frac{1}{a-j\omega}}_{\text{pole at }-ja} + \underbrace{\frac{1}{a+j\omega}}_{\text{pole at }+ja}.
\end{aligned}
\]

This is illustrated in Figure \ref{fig:LP_PZ}.\\

\begin{figure}[h!]
\centering
\includegraphics[scale=0.5]{figures/LP_PZ.png}
\caption{Pole-zero plot of $\Phi_\mathrm{xx}(\omega) = \frac{2a}{a^2 + \omega^2}$.}
\label{fig:LP_PZ}
\end{figure}

Note that positive time is mapped to THP poles and zeros, while negative time is mapped to LHP poles and zeros.\\

To simplify further, multiply the numerator and denominator of the second term by the complex conjugate of the denominator:
\[
\begin{aligned}
\Phi_\mathrm{xx}(\omega) &= \frac{a+j\omega}{a^2 - (j\omega)^2} + \frac{a-j\omega}{a^2 - (j\omega)^2} \\
&= \frac{2a}{a^2 + \omega^2}.
\end{aligned}
\]


This is illustrated in \ref{fig:PSD_omega} for \(a=2\).

\begin{figure}[h!]
\centering
\includegraphics[scale=0.5]{figures/PSD_omega.png}
\caption{The PSD of $\Phi_\mathrm{xx}(\omega) = \frac{2a}{a^2 + \omega^2}$ for \(a=2\).}
\label{fig:PSD_omega}
\end{figure}



\section{Wiener Filters}\label{Wiener}

\subsection{The General Case}

The general components of an optimal filtering system are shown in Figure 6.8. Let \(x(t)\) be the observed data and \(y(t)\) be the output of an LTI system. The desired signal is \(z(t)\), and the goal of designing a filter is to estimate \(z(t)\) that is optimum in some sense, that is, to apply some minimization criterion to the
error signal \(e(t)\). This criterion will depend on the application and its goals.

\begin{figure}[h!]
\centering
\includegraphics[scale=0.6
]{figures/Cont_LMMSE.png}
\caption{The general case for the optimal filtering. Here, \(x(t)\) is the input (observed data), \(z(t)\) is the desired output of the LTI system, \(y(t) = \hat{z(t)}\) is the actual output which is an estimate of \(z(t)\), and \(e(t)\) is the error.}
\label{fig:Cont_MMSE}
\end{figure}

In the most general sense, we have 
\[
y(t) = \hat{z}(t + \alpha)
\]

If $\alpha = 0$, we refer to this as \emph{filtering}. If $\alpha < 0$, it is  \emph{smoothing}. For $\alpha > 0$, we are performing  \emph{prediction}.\\

\begin{description}
    \item[Filtering ($\alpha = 0$)] \hfill \\
    Filtering is the process of estimating the \textit{current} state of a signal using all measurements available up to the present time $t$.
    \begin{itemize}
        \item \textbf{Goal:} To ``filter'' out noise and extract the best possible estimate of what is happening right now.
        \item \textbf{Example:} A noise-canceling headset estimating the current ambient noise to phase it out in real-time.
    \end{itemize}

    \item[Smoothing ($\alpha < 0$)] \hfill \\
    Smoothing is the process of estimating a \textit{past} state of a signal using data collected up to time $t$. Because you have access to data that occurred after the time point of interest ($t + \alpha$), you can create a much cleaner, ``smoother'' estimate than is possible in real-time.
    \begin{itemize}
        \item \textbf{Goal:} To obtain a high-accuracy historical record by utilizing ``future'' information (relative to the point being estimated).
        \item \textbf{Example:} Post-processing GPS data from a hike to create a clean map of where you walked, removing the ``zig-zags'' caused by momentary signal drops.
    \end{itemize}

    \item [Prediction ($\alpha > 0$)] \hfill \\
    Prediction is the process of estimating a future state of a signal using only the measurements available up to the present time $t$. Since you are trying to guess what hasn't happened yet, this is typically the "noisiest" and most difficult form of estimation.
    \begin{itemize}
    \item \textbf{Goal:} To anticipate future values or trends to allow for proactive adjustments or "look-ahead" logic.
    \item \textbf{Example:} A thermostat that "predicts" the outdoor temperature drop to start heating the house early, or a radar system predicting where a plane will be in 5 seconds to help with tracking.
    \end{itemize}
\end{description}

The Wiener filter is essentially the "optimal engine" that makes all three processes (filtering, smoothing, and prediction) work. While the terms we just discussed describe what you are trying to do with time, the Wiener filter provides the mathematical recipe for how to do it perfectly—specifically by minimizing the Mean Square Error (MSE).

\subsection{The Wiener filter}
Consider again that the general components of an optimal filtering system are shown in Figure 6.8. Let \(x(t)\) be the observed data and \(y(t)\) be the output of an LTI system. The desired signal is \(z(t)\), and the goal of designing a filter is to estimate \(z(t)\) that is optimum in some sense, that is, to apply some minimization criterion to the
error signal \(e(t)\). \\

If all signals are sample functions from stochastic processes and are Wide-Sense Stationary, and the optimization crtirion is mean squared error, then the optimum filter is the \emph{Wiener filter}.\\

As an example, consider a signal \(s(t)\) that we want to estimate from a measurable input \(x(t)\). It might be necessary to estimate \(s(t)\) if, for example, the input \(x(t)\) is composed of the signal \(s(t)\) corrupted by additive white noise \(\omega(t)\), as illustrated in Figure \ref{fig:Noise_cancel}.


\begin{figure}[h!]
\centering
\includegraphics[scale=0.6
]{figures/Noise_cancel.png}
\caption{The general case for noise cancellation.}
\label{fig:Noise_cancel}
\end{figure}

In this case, 
\[
z(t) = s(t)
\]
and
\[
y(t) = \hat{s(t)}
\]





\subsubsection*{Derivation}


First, we start with the orthogonality principle:

\[
\begin{aligned}
    E[e(t) x(u)] & = 0 \quad \forall u \\
    E[(z(t)-y(t))x(u)] & = 0 \quad \forall u\\
    E[z(t)x(u)] - E[y(t)x(u)] & = 0 \quad \forall u\\
    \phi_\mathrm{zx}(u-t) - \phi_\mathrm{yx}(u-t) & = 0  \quad \forall u
\end{aligned}
\]
Therefore,

\[
\phi_\mathrm{zx}(u-t) - \int_{-\infty}^\infty \phi_\mathrm{xx}(u-t-\nu)h(-\nu)\, d \nu = 0 \quad \forall u
\]

\[
\phi_\mathrm{xz}(t-u) - \int_{-\infty}^\infty \phi_\mathrm{xx}(t-u-\nu)h(\nu)\, d \nu = 0 \quad \forall u
\]

Let \(\tau = t-u\),

\[
\phi_\mathrm{xz}(\tau) - \int_{-\infty}^\infty \phi_\mathrm{xx}(\tau-\nu)h(\nu)\, d \nu = 0 \quad \forall \tau
\]
Therefore,
\[
\boxed{\phi_\mathrm{xz}(\tau) = \phi_\mathrm{xx}(\tau)*h_o(\tau)}
\]

This is the \emph{continuous Wiener-Hopf equation} where $h_o(t)$ is the optimal Wiener filter. Note also that since the cross correlation function is an odd function:
\[
\phi_\mathrm{zx}(\tau) = \phi_\mathrm{xx}(\tau)*h_o(-\tau)
\]

Also, since \(z(t) = y(t) - e(t)\):
\[
\phi_\mathrm{xz}(\tau) = \phi_\mathrm{xy}(\tau)
\]

This is because the error \(e(t)\) is orthogonal to the data \(x(t)\). It means that in an optimal system, the input $x(t)$ "sees" the hidden signal $z(t)$ exactly the same way it "sees" the estimate $y(t)$. If $\phi_{xz}(\tau) \neq \phi_{xy}(\tau)$, it would mean there is still some information about the signal $z(t)$ left in the error $e(t)$ that the filter failed to extract. By making these correlations equal, the filter has extracted every bit of useful information from the input $x(t)$ to create the estimate $y(t)$. 

\subsection{Non-Causal Solution to the Wiener-Hopf Equation}

Taking the Fourier transform of the continuous Wiener-Hopf equation:
\[
\mathcal{F} \left \{ \phi_\mathrm{xz}(\tau) - 
 \phi_\mathrm{xx}(\tau)*h_o(\tau)\right \}=  0
\]
Now,
\[
\Phi_\mathrm{xz}(\omega) - \Phi_\mathrm{xx}(\omega)\cdot H_o(\omega)=  0
\]
Therefore, 
\[
\boxed{H_o(\omega) = \frac{\Phi_\mathrm{xz}(\omega)}{\Phi_\mathrm{xx}(\omega)}}
\]

We will see that this is impractical in most cases, but a useful performance reference for causal systems.

\subsection{Causal Solution to the Wiener-Hopf Equation}

The causality constraint is defined as:

\[
h(t) = \begin{cases}
    0, & t<0 \\
    h(t) & t\ge 0
\end{cases}
\]

Without causality constraints, 
\[
\begin{aligned}
    E[e(t) x(u)] & = 0\\
    E[(z(t)-y(t))x(u)] & = 0 \\
    E \left [z(t)x(u) - \left \{ \int x(t-\nu)h(\nu) \, d \nu \right \} x(u) \right ] & = 0 \quad \forall u
\end{aligned}
\]

Now, introduce the causality constraints. The orthogonality condition $E[e(t)x(u)] = 0$ only needs to hold for the data we have actually observed. In a real-time (causal) system, we only have access to the input $x(u)$ up to the current time $t$. 
\[
 E \left [z(t)x(u) - \left \{ \int x(t-\nu)h(\nu) \, d \nu \right \} x(u) \right ] = 0 \quad \forall u \le t
\]


From before, this is:
\[
\phi_\mathrm{xz} (t-u) - \int_{-\infty}^{\infty}\phi_\mathrm{xx}(t-u-\nu)h(\nu) \, d \nu = 0 \quad \forall u \le t 
\]

If we define the lag as $\tau = t - u$, the condition must hold for all $\tau \ge 0$.

\[
\phi_\mathrm{xz} (\tau) - \int_{-\infty}^{\infty}\phi_\mathrm{xx}(\tau-\nu)h(\nu) \, d \nu = 0 \quad \forall \tau \ge 0 
\]

But, we cannot simply take the Fourier transform of this expression, since the RHS=0 \emph{only} for \(\tau \ge 0\). For \(\tau < 0\), RHS is undefined.\\

Consider a function:
\[
f(\tau) = 
\begin{cases}
    0, & \tau \ge 0 \\
    f(\tau), & \tau < 0
\end{cases}
\]

Therefore we can set the RHS=0:
\[
\phi_\mathrm{xz} (\tau) - \int_{-\infty}^{\infty}\phi_\mathrm{xx}(\tau-\nu)h(\nu) \, d \nu = f(\tau) \quad \forall \tau  
\]

Now, take the Fourier transform:

\[
\Phi_\mathrm{xz}(\omega) - H(\omega)\Phi_\mathrm{xx}(\omega) = F(\omega)
\]

Since \(f(\tau)=0\) for \(\tau \ge 0\), the poles and zeros of \(F(\omega)\) are all BHP (negative time).\\

Assume that $\Phi_\mathrm{xz}(\omega)$ and  $H(\omega)\Phi_\mathrm{xx}(\omega)$ are ratios of polynomials in $\omega$:

\[
\Phi_\mathrm{xx}(\omega) =  \underbrace{\Phi_\mathrm{xx}^-(\omega)}_{\text{BHP}}\underbrace{\Phi_\mathrm{xx}^+(\omega)}_{\text{THP}}
\]

where BHP and THP designates bottom and top half plane poles and zeros, respectively.\\

Therefore,
\[
\underbrace{\Phi_\mathrm{xx}^-(\omega)}_\text{BHP } \left [\underbrace{\frac{\Phi_\mathrm{xz}(\omega)}{\Phi_\mathrm{xx}^-(\omega)}}_\text{THP, BHP} - \underbrace{\Phi_\mathrm{xx}^+(\omega) H(\omega)}_\text{THP}\right ] = \underbrace{F(\omega)}_\text{BHP } 
\]

Since the RHS has only BHP poles and zeros, so does the left. This means that the THP poles and zeros of $\frac{\Phi_\mathrm{xz}(\omega)}{\Phi_\mathrm{xx}^-(\omega)}$ must cancel those of $\Phi_\mathrm{xx}^+(\omega) H(\omega)$ to leave only BHP poles and zeros\\

To do this we expand $\frac{\Phi_\mathrm{xz}(\omega)}{\Phi_\mathrm{xx}^-(\omega)}$ into partial fractions:

\[
\frac{\Phi_\mathrm{xz}(\omega)}{\Phi_\mathrm{xx}^-(\omega)} = \left [\frac{\Phi_\mathrm{xz}(\omega)}{\Phi_\mathrm{xx}^-(\omega)} \right ]^- +  \left [\frac{\Phi_\mathrm{xz}(\omega)}{\Phi_\mathrm{xx}^-(\omega)} \right ]^+
\]

In order for the THP poles and zeros in $\frac{\Phi_\mathrm{xz}(\omega)}{\Phi_\mathrm{xx}^-(\omega)}$ to cancel those of $\Phi_\mathrm{xx}^+(\omega) H(\omega)$, 
\[
\Phi_\mathrm{xx}^+(\omega) H(\omega) = \left [\frac{\Phi_\mathrm{xz}(\omega)}{\Phi_\mathrm{xx}^-(\omega)} \right ]^+
\]

Solving for the filter:

\[
\boxed{H_o(\omega) = \frac{1}{\Phi_\mathrm{xx}^+(\omega) }\left [\frac{\Phi_\mathrm{xz}(\omega)}{\Phi_\mathrm{xx}^-(\omega)} \right ]^+}
\]

This is the solution for the \emph{causal Wiener filter}.

\bigskip

\begin{tcolorbox}
\begin{exmp}{A Simple Wiener Filter.}
\end{exmp}

Consider a signal \(s(t)\) defined by its PSD:
\[\Phi_\mathrm{ss}(\omega) = \frac{3}{1+\omega^2}
\]
The observed data \(x(t)\) consists of the signal \(s(t)\) corrupted by additive zero-mean white noise \(w(t)\):
\[
\Phi_\mathrm{ww}(\omega) = 1
\]

Find the causal Wiener filter to recover \(s(t)\) from \(x(t)\).
\end{tcolorbox}

\bigskip
We assume that the signal and the noise are uncorrelated:
\[
E[s(t)w(t)] = E[s(t)]E[w(t)]
\]

Since $E[w(t)] = 0$ then 
\[
E[s(t)w(t)] = 0
\]

Since \(x(t) = s(t) + w(t)\), then 
\[
\Phi_\mathrm{xx}(\omega) = \Phi_\mathrm{ss}(\omega) + \Phi_\mathrm{sw}(\omega) + \Phi_\mathrm{ws}(\omega) + \Phi_\mathrm{ww}(\omega)
\]

Knowing that $E[s(t)w(t)] = 0$:

\[
\Phi_\mathrm{xx}(\omega) = \Phi_\mathrm{ss}(\omega) + \Phi_\mathrm{ww}(\omega)
\]

From the known signal and noise characteristics:
\[
\begin{aligned}
   \Phi_\mathrm{xx}(\omega) & = \frac{1}{1+\omega^2}+ 1\\
                        & = \frac{4+\omega^2}{1+\omega^2}\\
                        & = \frac{(2+j\omega)(2-j\omega)}{(1+j\omega)(1-j\omega)}
\end{aligned}
\]

Decomposing into THP and BHP:
\[
\Phi_\mathrm{xx}^+(\omega) = \frac{2+j\omega}{1+j\omega}
\]

\[
\Phi_\mathrm{xx}^-(\omega) = \frac{2-j\omega}{1-j\omega}
\]

Since $s(t)$ and $w(t)$ are uncorrelated:
\[
\Phi_\mathrm{xz}(\omega) \equiv \Phi_\mathrm{xs}(\omega) = \Phi_\mathrm{ss}(\omega) = \frac{3}{1+\omega^2}
\]
Here, \[
\Phi_\mathrm{xs}(\omega) = \Phi_\mathrm{ss}(\omega)
\]




\[
\begin{aligned}
   \text{Proof: \,\,\,\,} & \\
   \Phi_\mathrm{xs}(\omega) & = \mathcal{F} \left \{ E[x(t)s(t+\tau)] \right \} \\
   & = \mathcal{F} \left \{ E[(s(t) + w(t))s(t+\tau)] \right \}\\
   & = \mathcal{F} \left \{ E[(s(t) + s(t+ \tau)] + E[s(t+\tau) + w(t)] \right \}\\
   & = \Phi_\mathrm{ss}(\omega)
\end{aligned}
\]


\bigskip
Now,

\[
\begin{aligned}
    H(\omega)   & = \frac{1+j\omega}{2+j \omega} \left [ \frac {\frac{3}{1+ \omega^2}}{\frac{2-j \omega}{1-j \omega}} \right ]^+\\
            & = \frac{1+j\omega}{2+j \omega} \left [ \frac{3}{(2-j \omega)(1 + j \omega)} \right ]^+\\
            \text{partial fraction expansion} & \\
            & = \frac{1+j\omega}{2+j \omega} \left [ \underbrace{\frac{1}{2 - j \omega}}_{\text{BHP}} + \underbrace{\frac{1}{1 + j \omega}}_{\text{THP}} \right ] ^+\\
            & = \left (\frac{1+j\omega}{2+j \omega} \right ) \left (\frac{1}{1+j\omega} \right ) \\
            & = \frac{1}{2+j \omega}
\end{aligned}
\]

Therefore, the causal Wiener filter is 
\[
\begin{aligned}
H_o(\omega) & = \frac{1}{2+j \omega}  \\
h_o(t) & = e^{-2t} u(t)
\end{aligned}
\]
Which is a first-order lowpass filter, with an impulse function shown in Figure \ref{fig:LPF_impulse}.
\begin{figure}[h!]
\centering
\includegraphics[scale=0.3]{figures/LPF_impulse.png}
\caption{The impulse response of the optimal Wiener filter.}
\label{fig:LPF_impulse}
\end{figure}

\subsection{Evaluating the Error for a Wiener Filter}

Consider the mean square error:
\[
\begin{aligned}
    E[e^2(t)] &= E[\{z(t) - y(t) \}^2 ]\\
        & = \phi_\mathrm{zz}(0) + \phi_\mathrm{yy}(0) + 2\phi_\mathrm{zy}(0) \\
\end{aligned}
\]

But we know that 
\[
\begin{aligned}
\phi_\mathrm{yy}(0) & = \int \phi_\mathrm{xx}(\tau - u -\nu) h(u) h(\nu) \, du \, d \nu |_{\tau = 0} \\
    & = \int \int \phi_\mathrm{xx}(u -\nu) h(u) h(\nu) \, du \, d \nu \\
    & = \int \phi{_\mathrm{xy}}(u) h(u)\, du
\end{aligned}
\]
and also that 

\[
\begin{aligned}
\phi_\mathrm{zy}(\tau) & = \phi_\mathrm{zx}(\tau)* h(\tau)\\
        & = \int \phi_\mathrm{zx}(\tau - u) h(u) \, du \\
\phi_\mathrm{zy}(0) & = \int \phi_\mathrm{zx}(- u) h(u) \, du\\
        & = \int \phi_\mathrm{xz}(u) h(u) 
\end{aligned}
\]

Substituting these results:
\begin{align*}
E[e^2(t)] & = \phi_\mathrm{zz}(0)
              + \int\!\!\int \phi_\mathrm{xx}(u-\nu)\, h(u)\, h(\nu)\, du\, d\nu \\
           & \qquad - 2 \int \phi_\mathrm{xz}(u)\, h(u)\, du \\
           & = \phi_\mathrm{zz}(0)
              + \int h(u) \left[\int \phi_\mathrm{xx}(u-\nu)\, h(\nu)\, d\nu \right] du \\
           & \qquad - 2 \int \phi_\mathrm{xz}(u)\, h(u)\, du \\
           & = \phi_\mathrm{zz}(0)
              + \int \phi_\mathrm{xy}(u)\, h(u)\, du
              - 2 \int \phi_\mathrm{xz}(u)\, h(u)\, du \\
\intertext{\quad since $h(t) = h_0(t)$, then
           $\phi_\mathrm{xz}(\tau) = \phi_\mathrm{xy}(\tau)$, and}
           & = \phi_\mathrm{zz}(0)
              + \int \phi_\mathrm{xz}(u)\, h(u)\, du
              - 2 \int \phi_\mathrm{xz}(u)\, h(u)\, du \\
           & = \phi_\mathrm{zz}(0) - \int \phi_\mathrm{xz}(u)\, h(u)\, du \\
           & = \phi_\mathrm{zz}(0)
              - \int\!\!\int \phi_\mathrm{xx}(u-\nu)\, h(u)\, h(\nu)\, du\, d\nu \\
           & = \phi_\mathrm{zz}(0)
              - \phi_\mathrm{xx}(\tau) * h(\tau) * h(-\tau)\big|_{\tau=0}
\end{align*}

Taking the Fourier transform:

\[
\boxed{\text{MSE} = \frac{1}{2 \pi} \int_{-\infty}^\infty \Phi_\mathrm{zz}(\omega) - \Phi_\mathrm{xx}(\omega)|H_o(\omega)|^2  \, d \omega}
\]

\bigskip

\begin{tcolorbox}
\begin{exmp}{Continue with the previous example and evaluate the MSE.}
\end{exmp}
\end{tcolorbox}

\[
H_o(\omega) = \frac{1}{2+j \omega}, \quad \Phi_\mathrm{xx}(\omega) = \frac{4+\omega^2}{1+\omega^2}, \quad \Phi_\mathrm{zz}(\omega) = \Phi_\mathrm{ss}(\omega) = \frac{3}{1 + \omega^2}
\]

\[
\begin{aligned}    
\text{MSE} & = \frac{1}{2 \pi} \int_{-\infty}^\infty \left (\frac{3}{1 + \omega^2} \right )  - \left ( \frac{4+\omega^2}{1+\omega^2} \right) \left | \frac{1}{2+j \omega} \right |^2  \, d \omega \\
        & = \frac{1}{2 \pi} \int_{-\infty}^\infty \left (\frac{3}{1 + \omega^2} \right )  - \left ( \frac{4+\omega^2}{1+\omega^2} \right) \left | \frac{1}{4 + \omega^2} \right |  \, d \omega \\
        & = \frac{1}{2 \pi} \int_{-\infty}^\infty \frac{1}{1 + \omega^2 }\, d \omega\\
        & = 1
\end{aligned}
\]

Therefore, the \emph{percentage} error is the MSE normalized by the energy (mean square value) in the signal:
\[
\begin{aligned}
    \phi_\mathrm{ss}(0) & = \mathcal{F}^{-1}{\Phi_\mathrm{ss}}(\omega)|_{\tau=0}\\
            & = \frac{1}{2 \pi} \int_{-\infty}^\infty \frac{3}{1 + \omega^2} \, d \omega \\
            & = \frac{3}{2}
\end{aligned}
\]

Therefore:
\[
\frac{\text{MSE}}{\phi_\mathrm{ss}(0)} = \frac{2}{3} = 66\%
\]

\bigskip
For the non-causal filter:

\[
H_o(\omega) = \frac{\Phi_\mathrm{xz}(\omega)}{\Phi_\mathrm{xx}(\omega)} = \frac{3}{4 + \omega^2}
\]

\[
h_o(t) = \mathcal{F}^{-1} \left \{ \frac{3}{4 + \omega^2} \right \} = \mathcal{F}^{-1} \left \{ \frac{3/4}{2 + j\omega} + \frac{3/4}{2 - j\omega} \right \} = \frac{3}{4}e^{-2|t|}
\]

\[
\begin{aligned}    
\text{MSE} & = \frac{1}{2 \pi} \int_{-\infty}^\infty \left (\frac{3}{1 + \omega^2} \right )  - \left ( \frac{4+\omega^2}{1+\omega^2} \right) \left | \frac{3}{4 + \omega^2} \right |^2  \, d \omega \\
        & = \frac{3}{4}
\end{aligned}
\]

Therefore:
Therefore:
\[
\frac{\text{MSE}}{\phi_\mathrm{ss}(0)} = \frac{3}{4}\cdot \frac{2}{3} = \frac{1}{2}= 50\%
\]

This is an intuitive result, as the non-causal filter should always outperform a causal filter.

\bigskip

\subsection{The General Causal Case}
\label{sec:general-causal}

% ── Shorthand used throughout this section ─────────────────────────────────
% \providecommand is safe even if \F is already defined in the preamble.
\providecommand{\F}{\mathcal{F}}      % Fourier-transform operator

The general causal Wiener filter estimates a \emph{shifted} version of the
signal:
\[
  z(t) = s(t+\alpha), \qquad \alpha \in \mathbb{R}.
\]
The parameter $\alpha$ determines which of three classical estimation problems
is being solved, as illustrated in Figure~\ref{fig:alpha_timeline}.

\begin{figure}[htbp]
  \centering
  \includegraphics[width=\linewidth]{alpha_timeline.png}
  \caption{The role of the shift parameter $\alpha$ in the causal Wiener
    filter.  When $\alpha=0$ the filter estimates the signal at the current
    instant (\emph{filtering}); when $\alpha<0$ the desired signal
    $z(t)=s(t+\alpha)$ lies in the past (\emph{smoothing}); when $\alpha>0$
    it lies in the future (\emph{prediction}).  In every case the filter
    $h(t)$ must be causal, i.e.\ $h(t)=0$ for $t<0$.}
  \label{fig:alpha_timeline}
\end{figure}

The general filter structure is shown in Figure~\ref{fig:general_block}.
The causal filter $h(t)$ processes the noisy observation $x(t)$ to produce
the estimate $y(t)\approx z(t)=s(t+\alpha)$; the error
$e(t)=z(t)-y(t)$ is minimised in the mean-square sense.

\begin{figure}[htbp]
  \centering
  \begin{tikzpicture}[>=Stealth, node distance=1.8cm, auto,
      block/.style={draw, rectangle, rounded corners=2pt,
                    minimum width=2.0cm, minimum height=0.9cm,
                    fill=boxbg, font=\small},
      sum/.style  ={draw, circle, minimum size=0.75cm, fill=boxbg},
      wire/.style ={->, thick}]

    % Nodes
    \node (src)  [font=\small]                    {$x(t)$};
    \node (H)    [block,  right=1.6cm of src]     {$H(\omega)$};
    \node (plus) [sum,    right=1.8cm of H]       {$+$};
    \node (out)  [font=\small, right=1.6cm of plus] {$e(t)$};
    \node (zref) [font=\small, above=1.1cm of plus] {$z(t)=s(t+\alpha)$};

    % Wires
    \draw[wire] (src) -- (H);
    \draw[wire] (H)   -- node[midway, above, font=\small] {$y(t)$} (plus);
    \draw[wire] (zref) -- node[right, font=\small] {$-$} (plus);
    \draw[wire] (plus) -- (out);

    % Annotation below filter block
    \node [below=0.15cm of H, font=\footnotesize, text=gray]
          {causal: $h(t)=0,\ t<0$};
  \end{tikzpicture}
  \caption{Block diagram of the general causal Wiener filter.
    The causal filter $H(\omega)$ processes the noisy observation
    $x(t)=s(t)+n(t)$ to form the estimate $y(t)$ of the desired signal
    $z(t)=s(t+\alpha)$.  The error $e(t)=z(t)-y(t)$ is minimised in the
    mean-square sense subject to the causality constraint $h(t)=0$ for
    $t<0$.}
  \label{fig:general_block}
\end{figure}

% ── Spectral setup ─────────────────────────────────────────────────────────

\medskip\noindent
\textit{Continuing the running example} with
\[
  \Phi_{ss}(\omega) = \frac{3}{1+\omega^2}, \qquad
  \Phi_{nn}(\omega) = 1 \quad \text{(white noise)},
\]
the observed-data PSD and its spectral factors are
\[
  \Phi_{xx}(\omega)
    = \Phi_{ss}(\omega)+\Phi_{nn}(\omega)
    = \frac{4+\omega^2}{1+\omega^2}
    = \underbrace{\frac{2+j\omega}{1+j\omega}}_{\Phi_{xx}^+(\omega)}
      \underbrace{\frac{2-j\omega}{1-j\omega}}_{\Phi_{xx}^-(\omega)}.
\]

The cross-spectrum between the observation and the desired signal is
\begin{align*}
  \phi_{xz}(\tau)
    &= E\!\bigl[x(t)\,z(t+\tau)\bigr]
     = E\!\bigl[(s(t)+n(t))\,s(t+\tau+\alpha)\bigr]
     = \phi_{ss}(\tau+\alpha),\\[4pt]
  \Phi_{xz}(\omega)
    &= \F\!\bigl[\phi_{ss}(\tau+\alpha)\bigr]
     = \Phi_{ss}(\omega)\,e^{j\omega\alpha}
     = \frac{3\,e^{j\omega\alpha}}{1+\omega^2}.
\end{align*}

The causal Wiener--Hopf solution is therefore
\begin{equation}
  \boxed{H(\omega)
    = \frac{1}{\Phi_{xx}^{+}(\omega)}
      \left[\frac{\Phi_{xz}(\omega)}{\Phi_{xx}^{-}(\omega)}\right]^{\!+}}
  \label{eq:wiener-hopf}
\end{equation}
where $[\,\cdot\,]^+$ denotes the \emph{causal part} (inverse Fourier
transform, truncate to $t\ge 0$, then transform back).

Substituting the spectral factors:
\[
  \frac{\Phi_{xz}(\omega)}{\Phi_{xx}^{-}(\omega)}
    = \frac{3\,e^{j\omega\alpha}}{1+\omega^2}\cdot\frac{1-j\omega}{2-j\omega}
    = \frac{3\,e^{j\omega\alpha}}{(1+j\omega)(2-j\omega)}.
\]

There are three distinct cases depending on the sign of $\alpha$, treated below.

% ════════════════════════════════════════════════════════════════════════════
\subsection*{Case 1: $\alpha = 0$ \quad (Filtering Problem)}
\label{ssec:filtering}

When $\alpha=0$ the desired signal is the \emph{current} signal value,
$z(t)=s(t)$.  The causal-part operation is straightforward because no shift
complicates the partial-fraction structure.  A partial-fraction expansion gives
\[
  \frac{\Phi_{xz}(\omega)}{\Phi_{xx}^{-}(\omega)}\bigg|_{\alpha=0}
  = \frac{3}{(1+j\omega)(2-j\omega)}
  = \underbrace{\frac{1}{2-j\omega}}_{\text{anti-causal (BHP)}}
    +\underbrace{\frac{1}{1+j\omega}}_{\text{causal (THP)}},
\]
so $[\cdots]^+ = 1/(1+j\omega)$, and from~\eqref{eq:wiener-hopf}:
\[
  H_o(\omega) = \frac{1+j\omega}{2+j\omega}\cdot\frac{1}{1+j\omega}
              = \frac{1}{2+j\omega}
  \quad\Longrightarrow\quad
  h_o(t) = e^{-2t}\,u(t).
\]
This is a single-pole low-pass filter: it attenuates high-frequency noise
while passing the signal.  The impulse response is shown in
Figure~\ref{fig:case1_impulse}.

\begin{figure}[htbp]
  \centering
  \begin{tikzpicture}[xscale=1.6, yscale=1.4]
    % Axes
    \draw[->] (-0.3,0) -- (3.0,0) node[right] {$t$};
    \draw[->] (0,-0.12) -- (0,1.35) node[above] {$h_o(t)$};
    % Curve  h(t) = exp(-2t), t >= 0
    \draw[thick, color=accent, domain=0:2.6, samples=100]
          plot (\x, {exp(-2*\x)});
    % Initial value marker
    \filldraw[accent] (0,1) circle (2pt);
    \draw[dashed, gray] (0,1) node[left, black, font=\small] {$1$} -- (0.05,1);
    % Time-constant marker at t = 0.5 (time constant = 1/2)
    \draw[dashed, gray] (0.5, 0) node[below, font=\footnotesize] {$\tfrac{1}{2}$}
                                 -- (0.5, {exp(-1)});
    \draw[dashed, gray] (0, {exp(-1)}) node[left, font=\footnotesize] {$e^{-1}$}
                                       -- (0.5, {exp(-1)});
    % Label
    \node[accent, font=\small, right] at (1.1, 0.35) {$h_o(t)=e^{-2t}u(t)$};
    % t<0 region: h=0
    \draw[thick, color=accent] (-0.28,0) -- (0,0);
    \filldraw[accent] (0,1) circle (2pt);
    \filldraw[white]  (0,0) circle (2pt);
    \draw[accent,thick] (0,0) circle (2pt);
  \end{tikzpicture}
  \caption{Case~1 ($\alpha=0$, filtering): impulse response of the optimal
    causal Wiener filter $H_o(\omega)=1/(2+j\omega)$.  The filter is a
    first-order low-pass system with time constant $\tau=1/2$~s and unity
    gain at $\omega=0$.  For $t<0$, $h_o(t)=0$ (causality).}
  \label{fig:case1_impulse}
\end{figure}

% ════════════════════════════════════════════════════════════════════════════
\subsection*{Case 2: $\alpha < 0$ \quad (Smoothing Problem)}
\label{ssec:smoothing}

When $\alpha<0$ the desired signal $z(t)=s(t+\alpha)=s(t-|\alpha|)$ is a
\emph{past} value of the signal.  Even though the signal of interest lies in
the past, the causality constraint $h(t)=0,\ t<0$ still applies to the
filter — only past values of the observation $x(\cdot)$ may be used.

The exponential shift $e^{j\omega\alpha}$ prevents an immediate partial-fraction
factorisation, so the causal-part operation is performed in the time domain.
Using the inverse transform of $3/[(1+j\omega)(2-j\omega)]$:
\[
  \F^{-1}\!\left\{\frac{3}{(1+j\omega)(2-j\omega)}\right\}
  = e^{-t}u(t)+e^{2t}u(-t),
\]
delaying by $|\alpha|$ (i.e.\ replacing $t$ by $t+\alpha$) and retaining only
$t\ge 0$:

\begin{align}
  \left[\frac{\Phi_{xz}(\omega)}{\Phi_{xx}^{-}(\omega)}\right]^{\!+}
  &= \int_0^{-\alpha}e^{2(t+\alpha)}\,e^{-j\omega t}\,dt
   + \int_{-\alpha}^{\infty}e^{-(t+\alpha)}\,e^{-j\omega t}\,dt \notag\\[4pt]
  &= \frac{3\,e^{j\omega\alpha}}{(2-j\omega)(1+j\omega)}
   - \frac{e^{2\alpha}}{2-j\omega}.
   \label{eq:causal-part-smooth}
\end{align}

Substituting into~\eqref{eq:wiener-hopf}:
\begin{align}
  H_o(\omega)
  &= \frac{1+j\omega}{2+j\omega}\left[
       \frac{3\,e^{j\omega\alpha}}{(1+j\omega)(2-j\omega)}
      -\frac{e^{2\alpha}}{2-j\omega}\right] \notag\\[4pt]
  &= \frac{3\,e^{j\omega\alpha}}{4+\omega^2}
   - \frac{(1+j\omega)\,e^{2\alpha}}{4+\omega^2}.
   \label{eq:Ho-smooth}
\end{align}

Inverting~\eqref{eq:Ho-smooth} via standard Fourier-transform pairs yields the
time-domain impulse response:
\begin{equation}
  \boxed{h_o(t)
    = \frac{e^{-2|t+\alpha|}}{4}
      \Bigl[3 - e^{2\alpha} + 2\,e^{2\alpha}\operatorname{sgn}(t+\alpha)\Bigr]\,u(t)}
  \label{eq:ho-smooth-time}
\end{equation}
The structure of~\eqref{eq:ho-smooth-time} reflects the smoothing geometry:
for $0\le t < -\alpha$ (i.e.\ $t+\alpha<0$) the filter weight grows as
$e^{2(t+\alpha)}$, reaching a peak at $t=-\alpha$, then decays as
$e^{-(t+\alpha)}$ for $t>-\alpha$.  A jump discontinuity of magnitude
$e^{2\alpha}$ occurs at $t=-\alpha$ as a result of the causality constraint.

\medskip\noindent
\textbf{Sanity checks.}
\begin{itemize}
  \item $\alpha=0$: $h_o(t)=(e^{-2|t|}/4)[3-1+2\operatorname{sgn}(t)]
        = e^{-2t}u(t)$, consistent with Case~1. \checkmark
  \item $\alpha\to-\infty$: $e^{2\alpha}\to 0$, so
        $h_o(t)\to\tfrac{3}{4}e^{-2|t+\alpha|}$ — the bilateral
        (non-causal) filter shifted by $|\alpha|$ seconds, which becomes
        entirely causal as $|\alpha|\to\infty$. \checkmark
\end{itemize}

The impulse responses for three smoothing delays are shown in
Figure~\ref{fig:smoothing_impulse}.

\begin{figure}[htbp]
  \centering
  \includegraphics[width=\linewidth]{smoothing_wiener.png}
  \caption{Case~2 (smoothing, $\alpha<0$): impulse response
    $h_o(t)$ given by~\eqref{eq:ho-smooth-time} for three smoothing delays
    $\alpha\in\{-0.5,-1,-2\}$.
    Each response is zero for $t<0$ (causality) and peaks at $t=-\alpha$.
    The left branch ($0\le t < -\alpha$) grows as $e^{2(t+\alpha)}$; the
    right branch ($t>-\alpha$) decays as $e^{-(t+\alpha)}$.  Open circles
    mark the left-hand limit at the peak; filled circles mark the right-hand
    (adopted) value.  The jump $\Delta=e^{2\alpha}$ diminishes as the delay
    $|\alpha|$ increases, and as $\alpha\to-\infty$ the response converges
    to the non-causal bilateral exponential $\tfrac{3}{4}e^{-2|t+\alpha|}$.}
  \label{fig:smoothing_impulse}
\end{figure}

% ── Limit: infinite smoothing ──────────────────────────────────────────────
\subsubsection*{Limit: infinite smoothing delay}

As $\alpha\to-\infty$, equation~\eqref{eq:Ho-smooth} gives
\[
  H_o(\omega) \;\xrightarrow{\alpha\to-\infty}\; \frac{3\,e^{j\omega\alpha}}{4+\omega^2}
  \quad\Longrightarrow\quad
  h_o(t) = \tfrac{3}{4}\,e^{-2|t+\alpha|}.
\]
This is the \emph{non-causal} bilateral Wiener filter (Section~6.4.4)
shifted $|\alpha|$ seconds into the future.  Because every point of the
bilateral filter eventually falls in the causal half-line as
$|\alpha|\to\infty$, the smoother asymptotically achieves the non-causal MSE
performance.  The conclusion is:

\begin{quote}
  \itshape Introducing a sufficiently long smoothing delay allows a causal
  filter to approach the performance of an ideal non-causal filter.
\end{quote}

% ════════════════════════════════════════════════════════════════════════════
\subsection*{Case 3: $\alpha > 0$ \quad (Prediction Problem)}
\label{ssec:prediction}

When $\alpha>0$ the desired signal $z(t)=s(t+\alpha)$ is a \emph{future}
value of the signal.  The prediction geometry is illustrated in the left
panel of Figure~\ref{fig:prediction}.

The shifted time-domain function is $e^{-(t+\alpha)}u(t+\alpha)$; for $\alpha>0$
this function is positive for all $t>-\alpha$, so the portion with $t\ge 0$
is entirely contained in the causal half-line:
\[
  \left[\frac{\Phi_{xz}(\omega)}{\Phi_{xx}^{-}(\omega)}\right]^{\!+}
  = \F\!\Bigl[\bigl\{e^{-(t+\alpha)}\,u(t+\alpha)\bigr\}u(t)\Bigr]
  = \F\!\Bigl[e^{-\alpha}\,e^{-t}\,u(t)\Bigr]
  = \frac{e^{-\alpha}}{1+j\omega}.
\]

Substituting into~\eqref{eq:wiener-hopf}:
\[
  H_o(\omega)
    = \frac{1+j\omega}{2+j\omega}\cdot\frac{e^{-\alpha}}{1+j\omega}
    = \frac{e^{-\alpha}}{2+j\omega}
  \quad\Longrightarrow\quad
  \boxed{h_o(t) = e^{-\alpha}\,e^{-2t}\,u(t)}.
\]

The predictor has the same \emph{shape} as the filtering solution (Case~1)
but with its gain scaled down by the factor $e^{-\alpha}<1$.  This is
equivalent to running the filtering solution and attenuating the output:

\begin{center}
\begin{tikzpicture}[>=Stealth, node distance=2.4cm,
    block/.style={draw, rectangle, rounded corners=2pt,
                  minimum width=2.6cm, minimum height=0.9cm,
                  fill=boxbg, font=\small, align=center}]
  \node (in)   [font=\small] {$x(t)$};
  \node (filt) [block, right=of in]
               {$h(t)=e^{-2t}u(t)$\\[1pt]\footnotesize Causal filter\\[-2pt]\footnotesize ($\alpha=0$)};
  \node (gain) [block, right=of filt]
               {$\times\,e^{-\alpha}$\\[1pt]\footnotesize Prediction\\[-2pt]\footnotesize gain};
  \node (out)  [font=\small, right=of gain] {$\hat{s}(t+\alpha)$};
  \draw[->, thick] (in)   -- (filt);
  \draw[->, thick] (filt) -- (gain);
  \draw[->, thick] (gain) -- (out);
\end{tikzpicture}
\end{center}

The impulse responses for several prediction horizons $\alpha$ are shown in
the right panel of Figure~\ref{fig:prediction}.

\begin{figure}[htbp]
  \centering
  \includegraphics[width=\linewidth]{prediction_shift_diagram.png}
  \caption{Case~3 (prediction, $\alpha>0$).
    \emph{Left:} Timeline of the prediction problem.  At current time $t$
    only the shaded (past) observations $x(\tau),\ \tau\le t$ are available;
    the goal is to estimate the signal $s(t+\alpha)$ that will occur $\alpha$
    seconds in the future (red star).
    \emph{Right:} Impulse responses $h_o(t)=e^{-\alpha}e^{-2t}u(t)$ for
    $\alpha\in\{0,\,0.5,\,1.0,\,2.0\}$.  All responses share the same
    exponential decay shape but the initial amplitude is scaled by
    $e^{-\alpha}$: as the prediction horizon grows the filter trusts the
    observation less and its output tends to zero.  As
    $\alpha\to\infty,\;e^{-\alpha}\to 0$ and $\hat{s}(t+\alpha)\to 0$,
    consistent with the fact that the AR(1)-type process $s(t)$ becomes
    unpredictable over long horizons.}
  \label{fig:prediction}
\end{figure}

% ════════════════════════════════════════════════════════════════════════════
\subsection*{Summary: All Three Cases}

Figure~\ref{fig:three_cases} collects the impulse responses for a
representative value of $|\alpha|=1$ in each case.

\begin{figure}[htbp]
  \centering
  \includegraphics[width=\linewidth]{wiener_three_cases.png}
  \caption{Comparison of the optimal causal Wiener filter impulse responses
    for the three estimation problems (same signal and noise spectra in each
    case).
    \emph{Left} (filtering, $\alpha=0$): purely decaying exponential
    $h_o(t)=e^{-2t}u(t)$ starting at unity.
    \emph{Centre} (smoothing, $\alpha=-1$): two-sided exponential starting
    at a reduced amplitude, growing to its peak at $t=1=-\alpha$, then
    decaying; the jump discontinuity at the peak is $\Delta=e^{2\alpha}$.
    \emph{Right} (prediction, $\alpha=1$): purely decaying exponential with
    reduced initial amplitude $e^{-\alpha}=e^{-1}\approx 0.368$.
    In all cases the filter is causal ($h_o(t)=0$ for $t<0$).}
  \label{fig:three_cases}
\end{figure}

The key qualitative conclusions are:
\begin{itemize}
  \item \textbf{Filtering} ($\alpha=0$): the filter optimally trades off
        signal preservation against noise rejection via a first-order
        low-pass structure.

  \item \textbf{Smoothing} ($\alpha<0$): the delay $|\alpha|$ allows the
        filter to ``look ahead'' in its buffer; it uses data \emph{after}
        the desired time instant (but still in its past relative to the
        output) to form a better estimate.  Performance approaches the
        non-causal limit as $|\alpha|\to\infty$.

  \item \textbf{Prediction} ($\alpha>0$): the filter shape is identical to
        the filtering case, but the overall gain decreases exponentially with
        horizon $\alpha$.  For a process with finite spectral bandwidth,
        long-range prediction is inherently limited.
\end{itemize}



% ══════════════════════════════════════════════════════════════════════════════



\subsection{The Discrete Wiener Filter}
The discrete Wiener filter is realized as a finite impulse response (FIR) filter. The filter tap weights are optimized using the discrete realization of the Wiener-Hopf equation, evaluated using time-series estimate of the input autocorrelation and input/output cross-correlation functions. \\

Consider some observed data \( \mathbf{x} = \{x(1), x(2), \ldots x(N)\} \) consisting of a  discrete signal \(\mathbf{s} = \{s(1), s(2), \ldots s(n) \} \) corrupted by additive white noise \(\mathbf{w} = \{w(1), w(2), \ldots w(n) \} \). For convenience, we will use the symbols \(\mathbf{x} = x(n)\), \(\mathbf{s} = s(n)\), \(\mathbf{w} = w(n)\) and \(\mathbf{z} = z(n)\) interchangeably.\\


As with the continuous Wiener the desired signal \(z(n)\) is the signal \(s(n)\), as illustrated in Figure \ref{fig:Noise_cancel_discrete}. \\

\begin{figure}[h!]
\centering
\includegraphics[scale=0.6
]{figures/Noise_cancel_discrete.png}
\caption{The discrete Wiener filter.}
\label{fig:Noise_cancel_discrete}
\end{figure}

Since the noise is white, it is uncorrelated with the signal:
\[
E[s(n)w(n)] = 0
\]

The output of the filter is:
\[
y(n) = \hat{s}(n) = \sum_{k = 0}^\infty h(k)x(n-k)
\]


Now, consider the orthogonality principle:
\[
\begin{aligned}
    E[e(n) x(n-i)] & = 0 \quad \forall i \ge 0,\\
    E \left [\left \{ s(n) - \sum_{k=0}^\infty h(k) x(n-k) \right \} x(n-i) \right ] & = 0 \quad \forall i \ge 0,
\end{aligned}
\]

Therefore:
\[
\boxed{  \phi_\mathrm{xs}(i) - \sum_{k=0}^\infty h(k) \phi_\mathrm{xx}(i-k) = 0, \quad \forall i \ge 0}
\]

This is the \emph{discrete Wiener-Hopf equation}.\\


Now consider a finite data sequence \( \mathbf{x} = \{x(1), x(2), \ldots x(N)\} \). Using the finite data sequence means that the discrete Wiener-Hopf equation is an \emph{estimate}:
\[
\phi_\mathrm{xs}(i) - \sum_{k=0}^N h(k) \phi_\mathrm{xx}(i-k) \approx 0, \quad \forall i \ge 0
\]

In matrix form this is:
\[
\underbrace{\boldsymbol{\phi}_\mathrm{xs}}_{N\times 1} -   \underbrace{\boldsymbol{\phi}_\mathrm{xx}}_{N\times  N} \underbrace{\boldsymbol{h}^T(k)}_{N \times 1} =  \underbrace{\varnothing}_{\text{null}}
\]

Note that $\boldsymbol{\phi}_\mathrm{xx}$ is formed as a Toeplitz matrix. Let's expand the matrices:

\begin{equation*}
    \begin{bmatrix}
        \phi_\mathrm{xs}(0) \\
        \phi_\mathrm{xs}(1) \\
        \vdots \\
        \phi_\mathrm{xs}(N-1)
    \end{bmatrix}
    =
    \begin{bmatrix}
        \phi_\mathrm{xx}}(0) & \phi_\mathrm{xx}(1) & \cdots & \phi_\mathrm{xx}(N-1) \\
       \phi_\mathrm{xx}(1) & \phi_\mathrm{xx}(0) & \cdots & \phi_\mathrm{xx}(N-2)  \\
        \vdots & \vdots & \ddots & \vdots \\
        \phi_\mathrm{xx}(N-1) & \phi_\mathrm{xx}(N-2) & \cdots & \phi_\mathrm{xx}(0) \\
    \end{bmatrix}
    \begin{bmatrix}
        h(0)\\
        h(1)\\
        \vdots \\
        h(N-1)
    \end{bmatrix}
\end{equation*}

Note that what would normally be negative indices are positive, since \(\phi_\mathrm{xx}(i) \) is defined only for \(i\ge 0\). \\

Now, consider limiting the length of the correlations to \(M\) lags:  

\begin{equation*}
    \begin{bmatrix}
        \phi_\mathrm{xs}(0) \\
        \phi_\mathrm{xs}(1) \\
        \vdots \\
        \phi_\mathrm{xs}(M)
    \end{bmatrix}
    =
    \begin{bmatrix}
        \phi_\mathrm{xx}(0) & \phi_\mathrm{xx}(1) & \cdots & \phi_\mathrm{xx}(M) \\
       \phi_\mathrm{xx}(1) & \phi_\mathrm{xx}(0) & \cdots & \phi_\mathrm{xx}(M-1)  \\
        \vdots & \vdots & \ddots & \vdots \\
        \phi_\mathrm{xx}(M) & \phi_\mathrm{xx}(M-1) & \cdots & \phi_\mathrm{xx}(0) \\
    \end{bmatrix}
    \begin{bmatrix}
        h(0)\\
        h(1)\\
        \vdots \\
        h(M)
    \end{bmatrix}
\end{equation*}

To solve this, we rearrange the discrete Wiener-Hopf equation:
\[
\boldsymbol{h}^T(k)= \boldsymbol{\phi}_\mathrm{xx}^{-1} \boldsymbol{\phi}_\mathrm{xs}
\]

\begin{tcolorbox}
    \begin{exmp}{Simple Example of a Discrete Wiener Filter.}
    
    Consider observed data corrupted by additive white noise:
    \[
    x(n) = s(n) + w(n)
    \]
    
    The signal's autocorrelation function is defined as
\[
\boldsymbol{\phi}_\mathrm{ss} = 
    \begin{bmatrix}
        \phi_\mathrm{ss}(0) \\
        \phi_\mathrm{ss}(1) \\
        \phi_\mathrm{ss}(2)
    \end{bmatrix}
    = 
    \begin{bmatrix}
        \rfrac{3}{4}\\
        \rfrac{1}{2}\\
        \rfrac{1}{4}
    \end{bmatrix}
\]
and \(\phi\mathrm{ss}(i) = 0 \) for \(i \ge 3\). The additive white noise has:

\[
\boldsymbol{\phi}_\mathrm{ww} = 
    = 
    \begin{bmatrix}
        \rfrac{1}{4}\\
        0\\
        0
    \end{bmatrix}
\]
and 0 otherwise.

    \end{exmp}

Find the discrete Wiener filter to recover \(s(n)\) from \(x(n)\).

\end{tcolorbox}

We'll assume that the signal is uncorrelated with the noise:
\[
E[s(n)w(k)] = E[s]E[w] = 0
\]

Therefore
\[
\boldsymbol{\phi}_\mathrm{xx} =  \boldsymbol{\phi}_\mathrm{ss}+ \boldsymbol{\phi}_\mathrm{ww} = 
    \begin{bmatrix}
        1\\
        \rfrac{1}{2}\\
        \rfrac{1}{4}
    \end{bmatrix}
\]
and
\[
\boldsymbol{\phi}_\mathrm{xs} =  \boldsymbol{\phi}_\mathrm{ss}
\]

\[
\boldsymbol{\phi}_\mathrm{xs} = \boldsymbol{\phi}_\mathrm{xx} \boldsymbol{h}^T(k)= 
\]

\begin{equation}
  \begin{bmatrix}
        \rfrac{3}{4}\\
        \rfrac{1}{2}\\
        \rfrac{1}{4}
    \end{bmatrix}
    =
    \begin{bmatrix}
        1 & \rfrac{1}{2} & \rfrac{1}{4} \\
       \rfrac{1}{2} & 1 & \rfrac{1}{2}  \\
        \rfrac{1}{4} & \rfrac{1}{2}  & 1 \\
    \end{bmatrix}
    \begin{bmatrix}
        h(0)\\
        h(1)\\
        h(2)
    \end{bmatrix}
\end{equation}

\[
\boldsymbol{h}^T(k)= \boldsymbol{\phi}_\mathrm{xx}^{-1} \boldsymbol{\phi}_\mathrm{xs} =      = 
       \begin{bmatrix}
        \rfrac{2}{3}\\
        \rfrac{1}{6}\\
        0
    \end{bmatrix}
\]

This is illustrated in Figure \ref{fig:Tapped_delay}.
\begin{figure}[h!]
\centering
\includegraphics[scale=0.15]{figures/Tapped_delay2.png}
\caption{The tapped delay line representation of the resulting FIR Wiener filter.}
\label{fig:Tapped_delay}
\end{figure}


\begin{tcolorbox}
    \begin{exmp}{Recovery of a Noisy Sinusoid.} 
    
    Consider a sinuoisal signal of length \(N=1000\) corrupted by white Gaussian noise. Construct a discrete Wiener filter that will recover the sinusoid with the least mean square error.   

    \end{exmp}


 \end{tcolorbox}

The code \href{https://www.ece.unb.ca/kengleha/Textbook/Code/FIRWienerFilter.py}{\texttt{FIRWienerFilter.py}} implements this filter. The key components are:

\begin{itemize}
    \item Assign the length of the data window and the filter length.\\
    \texttt{\# Parameters\\
    N = 1000 \quad \# Number of samples\\
    M = 200  \quad \,   \# Filter order (one less than the matrix size)\\
    }
    \item Compute the autocorrelation matrix Rxx with dimensions 21x21\\
    \texttt{Rxx = toeplitz(np.correlate(x, x, mode='full')[:M+1])}
    \item Compute the cross-correlation vector Rxy\\
    \texttt{Rxy = np.correlate(x, y, mode='full')[:M+1]}
    \item Solve the Wiener-Hopf equations to find the filter coefficients\\
    \texttt{h = np.dot(pinv(Rxx), Rxy)}
    \item Apply the Wiener filter to the noisy signal\\
    \texttt{filtered\_signal = filtfilt(h, 1.0, y)}
\end{itemize}

The noisy sinusoid and the estimated signal using the Wiener filter are shown in \ref{fig:FIR_sinsusoid}.
\begin{figure}[h!]
\centering
\includegraphics[scale=0.4]{figures/FIRSinusoid.png}
\caption{The noisy sinusoid and the estimated signal using the Wiener filter.}
\label{fig:FIR_sinsusoid}
\end{figure}



    
\subsection{Wiener Filters: Fixed Plant Component}
Consider the scenario where an LTI system exists in the system that cannot be modified (fixed plants). In this case, the Wiener filter be be optimized with regard to the output of the fixed plant, as illustrated in Figure \ref{fig:Fixed_plant}.

\begin{figure}[h!]
\centering
\includegraphics[scale=0.5]{figures/Fixed_plant.png}
\caption{Optimal filter design with a fixed plant $h_F(t)$.}
\label{fig:Fixed_plant}
\end{figure}

That is, given $h_F(t)$, find $h(t)$ such that the overall system is optimal.

\[
H_o(\omega) = \frac{1}{\Phi_\mathrm{x_1 x_1}^+(\omega) }\left [\frac{\Phi_\mathrm{zx_1}(\omega)}{\Phi_\mathrm{x_1 x_1}^-(\omega)} \right ]^+
\]

In terms of the input $X(t)$:
\[
\Phi_\mathrm{x_1 x_1}(\omega) = \Phi_\mathrm{xx}(\omega)\left | H_F(\omega) \right |^2
\]
\[
\Phi_\mathrm{z_1x_1}(\omega) = \Phi_{zx}(\omega) H_F^*(\omega)
\]

Therefore,

\[
H_o(\omega) = \frac{1}{\left [ \Phi_\mathrm{xx}(\omega)\left | H_F(\omega) \right |^2 \right ]^+ } \left [\frac{\Phi_\mathrm{zx}(\omega) H_F^*(\omega)}{\left [ \Phi_\mathrm{xx}(\omega)\left | H_F(\omega) \right |^2 \right ]^- }  \right ]^+
\]
\section{Kalman Filter}
Kalman filtering uses a system's dynamic model (e.g., physical laws of motion), known control inputs to that system, and multiple sequential measurements (such as from sensors) to form an estimate of the system's varying quantities (its state) that is better than the estimate obtained by using only one measurement alone. As such, it is a common sensor fusion and data fusion algorithm.\\


Imagine we have a car equipped with GPS and an inertial measurement unit (IMU) consisting of accelerometers and gyroscopes. The GPS provides periodic but somewhat noisy measurements of the car's position, while the IMU provides continuous but subject to drift measurements of the car's velocity and orientation.\\

Now, let's say we want to accurately estimate the car's position, velocity, and orientation at all times. We can use a Kalman filter to fuse the information from the GPS and IMU to obtain a more accurate estimate than using either sensor alone. \\

Here's how the Kalman filter would work in this scenario:

\begin{itemize}
    \item \textbf{System Model}: We define a dynamic model of the car's motion based on the laws of physics, taking into account factors like acceleration, velocity, and orientation changes. This model predicts how the car's state (position, velocity, orientation) evolves over time given control inputs such as throttle, steering angle, and brake.
    \item \textbf{Control Inputs}: We incorporate information about the control inputs to the car, such as throttle position and steering angle. These inputs affect the predicted state of the car according to the dynamic model.
    \item \textbf{Measurements:} We receive periodic GPS measurements of the car's position and potentially noisy continuous measurements from the IMU about its velocity and orientation.
    \item \textbf{Kalman Filtering:} We use the Kalman filter to combine the predictions from the dynamic model with the measurements from the GPS and IMU. The filter calculates a weighted average of the predicted state (based on the dynamic model) and the measured state (from the sensors), with the weights determined by the relative uncertainties of the predictions and measurements. This process provides an optimal estimate of the car's state that minimizes the mean squared error.
    \item \textbf{State Estimation:} The output of the Kalman filter is an estimate of the car's position, velocity, and orientation that is more accurate than using just the GPS or IMU alone. The filter continually updates this estimate as new measurements arrive, effectively filtering out noise and correcting for drift in the IMU measurements.
\end{itemize}


By fusing information from multiple sensors using the Kalman filter, we obtain a more reliable and accurate estimate of the car's state.\\

The Kalman filter produces an estimate of the state of the system as an average of the system's predicted state and of the new measurement using a weighted average. The purpose of the weights is that values with better (i.e., smaller) estimated uncertainty are "trusted" more. This process is repeated at every time step, with the new estimate and its covariance informing the prediction used in the following iteration. This means that Kalman filter works recursively and requires only the last "best guess", rather than the entire history, of a system's state to calculate a new state.\\


\subsection{The System Dynamic Model}
Kalman filtering is based on linear dynamic systems discretized in the time domain. The state of the target system refers to the ground truth (yet hidden) system configuration of interest, which is represented as a vector of real numbers. At each discrete time increment, a linear operator is applied to the state to generate the new state, with some noise mixed in, and optionally some information from the controls on the system. Then, another linear operator mixed with more noise generates the measurable outputs (i.e., observation) from the true ("hidden") state.\\ 

In order to use the Kalman filter to estimate the internal state \(\mathbf{x}_k\) of a process at time step \(k\) given only a sequence of noisy observations \(\{\mathbf{z}_k, \mathbf{z}_{k-1}\ldots \}\), one must model the process in accordance with the following framework. This means specifying the matrices, for each time-step k:

\begin{itemize}
    \item $F_k$: the state-transition model;
    \item $H_k$: the observation model;
    \item $Q_k$: the covariance of the process noise;
    \item $R_k$: the covariance of the observation noise;
    \item Sometimes $B_k$: the control-input model as described below; if $B_k$ is included, then there is also
    \item $u_k$: the control vector, representing the controlling input into the control-input model.
\end{itemize}

This model is illustrated in Figure \ref{fig:Kalman_model}.\\

\begin{figure}[h!]
\centering
\includegraphics[scale=0.7]{figures/Kalman_filter_model_2.png}
\caption{The model underlying the Kalman filter. Squares represent matrices. Ellipses represent multivariate normal distributions (with the mean and covariance matrix enclosed). Unenclosed values are vectors.}
\label{fig:Kalman_model}
\end{figure}


The Kalman filter model assumes the true state at time $k$ is evolved from the state at $(k - 1)$ according to:
\[
\mathbf{x}_k = \mathbf{F}_k \mathbf{x}_{k-1} + \mathbf{B}_k \mathbf{u}_k + \mathbf{w}_k
\]
where:
\begin{itemize}
    \item $\mathbf{F}_k$: the state transition model applied to the previous state $\mathbf{x}_{k-1}$;
    \item $\mathbf{B}_k$ is the control-input model which is applied to the control vector $\mathbf{u}_k$;
    \item $\mathbf{w}_k$ is the process noise, which is assumed to be drawn from a zero-mean multivariate normal distribution with zero mean and covariance $\mathbf{Q}_k$. That is: $\mathbf{w}_k \sim \mathcal{N}(0,\mathbf{Q}_k)$.
\end{itemize}


At time $k$, an observation (or measurement) $\mathbf{z}_k$ of the true state $\mathbf{x}_k$ is made. The relationship between the measurement and the true state space is: 

\[
\mathbf{z}_k = \mathbf{H}_k \mathbf{x}_k + \mathbf{v}_k
\]

where
\begin{itemize}
    \item $\mathbf{H}_k$ is the observation model, which maps the true state space into the observed space; and
    \item $\mathbf{v}_k$ is the observation noise, which is assumed to be zero-mean Gaussian white noise with covariance $\mathbf{R}_k$. That is: $\mathbf{v}_k \sim \mathcal{N}(0,\mathbf{R}_k)$.

\end{itemize}
The initial state, and the noise vectors at each step $\{\mathbf{x}_0, \mathbf{w}_1, \ldots, \mathbf{w}_k, \mathbf{v}_1, \ldots, \mathbf{v}_k\}$ are all assumed to be mutually independent.

\subsection{Formulation}
The Kalman filter is a recursive estimator. This means that only the estimated state from the previous time step and the current measurement are needed to compute the estimate for the current state. In contrast to batch estimation techniques, no history of observations and/or estimates is required. In what follows, the notation $\hat{\mathbf{x}}_{n\mid m}$ represents the estimate of $\mathbf{x}$ at time $n$ given observations up to and including time $m \leq n$.

The state of the filter is represented by two variables:
\begin{itemize}
    \item $\mathbf{x}_{k\mid k}$, the \emph{a posteriori} state estimate mean at time $k$ given observations up to and including time $k$;
    \item $\mathbf{P}_{k\mid k}$, the \emph{a posteriori} state estimate covariance matrix (a measure of the estimated accuracy of the state estimate).

\end{itemize}

The Kalman filter  is most often conceptualized as two distinct phases: `Predict' and `Update'.

\begin{itemize}
    \item The \textbf{predict phase} uses the state estimate from the previous timestep \(k-1\) to produce an estimate of the state at the current timestep \(k\): $\mathbf{\hat{x}}_{k\mid k-1}$. This predicted state estimate is also known as the \emph{a priori} state estimate because, although it is an estimate of the state at the current timestep, it does \emph{not} include observation information from the current timestep ($\mathbf{z}_{k}$).
    \item In the \textbf{update phase}, the \emph{innovation} is computed. This is the difference between the current \emph{a priori} prediction of the observation and the current observation: \(\mathbf{\Tilde{y}}_k = \mathbf{z_k} - \mathbf{H_k}\mathbf{\hat{x}}_{k\mid k-1}\). This is multiplied by the optimal Kalman gain $\mathbf{K}_k$ and combined with the previous state estimate to refine the state estimate: $\mathbf{x}_{k\mid k} = \mathbf{\hat{x}}_{k\mid k-1} + \mathbf{K}_k \mathbf{\Tilde{y}}_k$. This improved estimate based on the current observation is termed the \emph{a posteriori} state estimate.
\end{itemize} 

Typically, the two phases alternate, with the prediction advancing the state until the next scheduled observation, and the update incorporating the observation. \\

The complete set up update and predict equations are as follows: \\

\begin{tabular}{ll}
\textbf{Predict} & \\
Predicted (\emph{a priori}) state estimate      & $\mathbf{\hat{x}}_{k\mid k-1}  = \mathbf{F}_k \mathbf{{x}}_{k-1\mid k-1} + \mathbf{B}_k \mathbf{u}_k$\\
Predicted (\emph{a priori}) estimate covariance     & $\mathbf{\hat{P}}_{k\mid k-1} =  \mathbf{F}_k \mathbf{{P}}_{k-1\mid k-1} \mathbf{F}^{\text{T}}_k + \mathbf{Q}_k$\\
    &   \\
\textbf{Update} & \\   
{Innovation} & \(\mathbf{\Tilde{y}}_k = \mathbf{z_k} - \mathbf{H_k}\mathbf{\hat{x}}_{k\mid k-1}\) \\
{Innovation covariance} & $\mathbf{S}_k = \mathbf{H_k}\mathbf{\hat{P}}_{k\mid k-1} \mathbf{H}^{\text{T}}_k + \mathbf{R}_k$\\
{Optimal Kalman gain} & $\mathbf{K}_k = \mathbf{\hat{P}}_{k\mid k-1} \mathbf{H}^{\text{T}}_k \mathbf{S}^{-1}_k$\\
{Updated (\emph{a posteriori}) state estimate} & $\mathbf{\hat{x}}_{k\mid k}  = \mathbf{\hat{x}}_{k\mid k-1} + \mathbf{K}_k \mathbf{\Tilde{y}}_k$\\
{Updated (\emph{a posteriori}) estimate covariance} & $\mathbf{P}_{k\mid k} = (\mathbf{I - \mathbf{K}_k \mathbf{H}_k)}\mathbf{\hat{P}}_{k\mid k-1}$\\
{Measurement post-fit residual} & $\mathbf{\Tilde {y}}_{k \mid k} =  \mathbf{z}_k - \mathbf{H}_k \mathbf{\Tilde{x}}_{k \mid k }$
\end{tabular}


\bigskip
The formula for the the optimal Kalman gain \(\mathbf{K}_k \) is that which minimizes the residual error \(\Tilde{y}_{k \mid k}\). Derivation of the optimal Kalman gain and the \emph{a posteriori} estimate covariance matrix is provided in Appendix A.\\

An alternative (intuitive) way to express the updated state estimate (\(\mathbf{x}_{k\mid k}\)) is:

\[
\mathbf{x}_{k\mid k}  = ( \mathbf{I} - \mathbf{K}_k \mathbf{H}_k ) \mathbf{\hat{x}}_{k\mid k-1} + \mathbf{K}_k \mathbf{z}_k
\]

This expression reminds us of a linear interpolation, \(y = (1-\alpha) a + \alpha b \), for \(\alpha\) between 0 and 1. In our case:

\begin{itemize}

    \item \(a \rightarrow \mathbf{\hat{x}}_{k\mid k-1}\), the internal state estimated from the model.
    \item \(b \rightarrow \mathbf{H}^{^1}_k \mathbf{z}_k\) is the internal state estimated from the measurement.
    \item \(\alpha \rightarrow \mathbf{K}_k \mathbf{H}_k\) can range from 0 (high error in the measurement) \(\mathbf{z}_k\) to 1 (low error).
\end{itemize}

Therefore, the Kalman gain (proportional to \(\alpha \)) can be seen to provide a weighted average of the model and the measurement, depending on their relative uncertainty (error).


\bigskip

\begin{tcolorbox}
    \begin{exmp}{Intuitive Illustration of  Kalman Update and Predict.}

    By definition, the values (and distribution) of the state \(\mathbf{x}\) are hidden (and thus unobservable). If we assume that the state and the measurements are normally distributed, we can illustrate the influence of the predict and update stages for a single step.  We show the effect of sensor measurement that has low and high variance (uncertainty), relative to the uncertainty in the state model. Note that for this example, there is no control signal \(\mathbf{u}_k\).
        
    \end{exmp}
\end{tcolorbox}

We consider a system with the prior state \(\mathbf{x}_{k-1 \mid k-1} \sim \mathcal{N}(\mu_{p},\sigma_p)\) We will assume that the state transition model \(\mathbf{F}_k\) is a constant velocity (per step) with inherent noise \( \sim \mathcal{N}(\mu_{w},\sigma_w) \). This yields a  predicted or \emph{a priori} state estimate \( \mathbf{\hat{x}}_{k\mid k-1} \sim \mathcal{N}(\mu_{ap},\sigma_{ap})\) with 

\[
\begin{aligned}
   \mu_{ap} & = \mu_p + \mu_w \\
   \sigma_{ap} & = \sqrt{\sigma_p^2 + \sigma_w^2}
\end{aligned}
\]

Next, the measurement \(\mathbf{z}_k 
 \sim \mathcal{N}(\mu_z,\sigma_z)\) is used to update the state estimate via the product of the innovation and the Kalman gain: $\mathbf{\hat{x}}_{k\mid k}  = \mathbf{\hat{x}}_{k\mid k-1} + \mathbf{K}_k \mathbf{\Tilde{y}}_k$.\\

Recall that the Kalman gain can range from 0 (high error in the measurement, relative to the process)  to 1 (low error).
 Figure \ref{fig:Kalman__pdf_example} illustrates two cases: one where the measurement error is low (thus having a marked effect on the  \emph{a posteriori} state estimate), and one where measurement error is high (having relatively little effect). Note that the noisy measurement not only introduces a minimal shift in the mean from the \emph{a priori} state model to the \emph{a posteriori} state estimate but also incorporates a significant level of uncertainty into the \emph{a posteriori} state estimate.

 
\begin{figure}[h!]
\centering
\includegraphics[scale=0.6]{figures/Kalman_good_sensor.png}
\includegraphics[scale=0.6]{figures/Kalman_noisy_sensor.png}
\caption{An illustration of the state and measurement variable distributions for a single  Kalman predict/update step. Here, the fixed variables are \(\mu_p=5.0, \, \sigma_p=1.0, \, \mu_{ap} = \mu_{p} + 1.0, \, \sigma_w = 0.5, \) and  \(\mu_z = 5.2\). The figures depicts two cases. Top: measurement noise that is low relative to the noise inherent in the state model ($\sigma_z=0.4$); this yields $\mathbf{K}_k=0.9$.  Bottom: measurement noise that is high ($\sigma_z=1.4$); this yields $\mathbf{K}_k=0.4$.}
\label{fig:Kalman__pdf_example}
\end{figure}

\begin{tcolorbox}
    \begin{exmp}{The Kalman filter's application in allowing GPS to be robust.}\\

\textbf{Problem Statement:}

Suppose we have a GPS receiver in a moving vehicle, and you want to accurately estimate the vehicle's position and velocity while driving through urban canyons with tall buildings that can block or reflect GPS signals, leading to noisy and unreliable measurements.\\

\textbf{Kalman Filter Application:}

A Kalman filter can be employed to improve the accuracy of position and velocity estimates by fusing noisy GPS measurements with a dynamic motion model of the vehicle.

    \end{exmp}
\end{tcolorbox}

\textbf{System Model:}

\begin{enumerate}
    \item \textbf{State Variables:} The state vector includes the position (latitude and longitude) and velocity (speed and direction) of the vehicle:\\ 
   - State Vector, $\mathbf{x}_k$ = [latitude, longitude, speed, direction]
   \item \textbf{State Transition Model (Predict):} The state evolves with time based on the vehicle's dynamics (e.g., constant velocity model). This is the prediction step.\\
      - Predicted State, $\mathbf{\hat{x}}_{k|k-1} = \mathbf{F}_k \mathbf{x}_{k-1|k-1}$\\
   - $\mathbf{F}_k$: State transition model
   
    \item \textbf{Observation Model (Update):} GPS provides measurements of the vehicle's position (latitude and longitude) with some noise.\\
   - Measurement: \(\mathbf{z}_k\) = [measured\_latitude, measured\_longitude]\\
   - Measurement Equation: $\mathbf{z}_k = \mathbf{H}_k \mathbf{x}|_k + \mathbf{v}_k$\\
   - $\mathbf{H}_k$: Observation model\\
   - $\mathbf{v}_k$: Measurement noise (GPS measurement error)

\end{enumerate}

\bigskip


\textbf{Kalman Filter Steps:}

\begin{enumerate}
    \item \textbf{Initialization:} Initialize the Kalman filter with an initial state estimate and covariance matrix.

    \item \textbf{Prediction} (Time Update):
\begin{align*}
\text{Predicted State, } \mathbf{\hat{x}}_{k|k-1} &= \mathbf{F}_k \cdot \mathbf{x}_{k-1 \mid k-1} \\
\text{Predicted Covariance, } \mathbf{P}_{k|k-1} &= \mathbf{F}_k \mathbf{P}_{k-1} \mathbf{F}_k^T + \mathbf{Q}_k
\end{align*}
Where:
\(\mathbf{Q}_k\): Process noise covariance

\item \textbf{Measurement Update} (Correction): Update the state estimate using GPS measurements.

\begin{align*}
\text{Kalman Gain, } \mathbf{K}_k &= \mathbf{P}_{k|k-1} \mathbf{H}_k^T  (\mathbf{H}_k \mathbf{P}_{k|k-1} \mathbf{H}_k^T + \mathbf{R}_k)^{-1} \\
\text{Updated State Estimate, } \hat{\mathbf{x}}_{k|k} &= \hat{\mathbf{x}}_{k|k-1} + \mathbf{K}_k  (\mathbf{z}_k - \mathbf{H}_k  \hat{\mathbf{x}}_{k|k-1}) \\
\text{Updated Covariance, } \mathbf{P}_{k|k} &= (\mathbf{I} - \mathbf{K}_k \mathbf{H}_k)  \mathbf{P}_{k|k-1}
\end{align*}
Where:
\(\mathbf{R}_k\): Measurement noise covariance (GPS measurement error)

\end{enumerate}

\begin{tcolorbox}
    \begin{exmp}{Numerical (1D) Example of Kalman filter with GPS }
        In this example, we'll create a simple 1D Kalman filter to estimate the position of a moving vehicle using noisy GPS measurements. In this  this example we assume constant velocity motion along a straight path.
    \end{exmp}
\end{tcolorbox}
For the first 20 time steps, the GPS has not triangulated and so no measurements are available. Whereas the vehicle is moving at constant velocity, assume the transition model of the vehicle's position is subject to a constant drift. This drift can be seen in the first 20 time steps of Figure \ref{fig:Kalman_GPS}. No measurement updates are available during this time.  Once the GPS location data are available (with some inherent noise), the drift is corrected. It is also clear that the Kalman filter allows the GPS measurements to reduce the variance in the process model (the \emph{a priori} estimate of position). 
\begin{figure}[h!]
\centering
\includegraphics[scale=0.5]{figures/Kalman_GPS_position.png}
\includegraphics[scale=0.5]{figures/Kalman_GPS_offset.png}
\caption{The position (top) and the offset from the true position (bottom).}
\label{fig:Kalman_GPS}
\end{figure}

The code for this example is available as \href{https://www.ece.unb.ca/kengleha/Textbook/Code/Kalman_Example_GPS.py}{\texttt{Kalman\_Example\_GPS.py}}.\\

A second example compares a GPS with a Kalman filter to a state estimate using only a system model (using IMUs), and in the case where the IMU sensors drift away from the true position (a constant velocity profile). Figure \ref{fig:KalmanDrift} illustrates the position estimation with and without the use of GPS data. Clearly, even though the GPS measurements are noisy, they correct the offset in position due to drift in the system model.\\

\begin{figure}[h!]
\centering
\includegraphics[scale=0.6]{figures/KalmanDrift.png}
\caption{The effect of estimate drift without a Kalman filter.}
\label{fig:KalmanDrift}
\end{figure}

The code for this example is available as \href{https://www.ece.unb.ca/kengleha/Textbook/Code/KalmanDrift.py}{\texttt{KalmanDrift.py}}.\\

\chapter{Analysis of Time-Series Data}
In most real-world scenarios, the estimation of ensemble statistics is not practical as an ensemble of sample functions is not available. Rather, only a single sample function exists. Having only one sample function. we must make \emph{time-averages} rather than \emph{ensemble averages}. Therefore, it is required that the process be wide-sense ergodic (which implies wide-sense stationarity).\\

\section{Tests of Self-Stationarity}
The assertion of true wide-sense stationarity requires ensemble statistics. If we have only one sample function \(x_i(t)\), we must consider self-stationarity of that sample function.\\

\subsubsection{Trend Testing}
To assess self-stationarity in the wide-sense, we must consider the stationarity of the time-average mean and time-average autocorrelation function. From a single sample function \(x_i(t)\) of length \(T\) seconds, we can divide the entire time record into \(N\) intervals of length $\Delta$. This gives \(N\) epochs \(\{ 0, T_1, \ldots T_N\), where \(T_N = T\). Figure \ref{fig:Time_stationarity} depicts this scenario for \(N=5\). \\

\begin{figure}[h!]
\centering
\includegraphics[scale=0.4]{figures/Time_stationarity.png}
\caption{The use of epochs to assess self-stationarity.}
\label{fig:Time_stationarity}
\end{figure}

Using these epochs of length $\Delta$, we can compute an ensemble of \(N\) time-average means and autocorrelation functions. \\

First, the ensemble of time-average means:    

\[
\begin{aligned}
    \hat{\mu}_{1} & = \frac{1}{\Delta}\int_{0}^{T_1}x_i(t) \, dt \\
    \vdots \\
    \hat{\mu}_{N} & = \frac{1}{\Delta}\int_{T_{N-1}}^{T_N}x_i(t) \, dt 
\end{aligned}
\]

And, the ensemble of time-average autocorrelation functions:    

\[
\begin{aligned}
    R_\mathrm{xx_1}(\tau) & = \frac{1}{\Delta}\int_{0}^{T_1}x_i(t) x_i(t + \tau)\, dt \\
    \vdots \\
    R_{\mathrm{xx_1}}(\tau) & = \frac{1}{\Delta}\int_{T_{N-1}}^{T_N}x_i(t) x_i(t+\tau)\, dt 
\end{aligned}
\]

To test for trends in time-series data, one can use statistical tests or visual inspection. Common methods include:
\begin{itemize}
    \item Augmented Dickey-Fuller Test: It is a common test for unit roots in time series data. If the null hypothesis is rejected, it suggests the presence of a trend.
    \item Kendall's Tau or Spearman's Rank Correlation: These non-parametric tests can identify trends in data.
    \item Visual Inspection: Plot the time series data and look for patterns that indicate trends, such as linear or nonlinear trends.
\end{itemize}

\bigskip
\textbf{Challenge: The choice of the epoch length, $\Delta$}
\begin{itemize}
    \item \(\Delta << T \) to get a suffucnet number if epochs in the ensemble to test for trends.
    \item $\Delta$ must be large enough to give a time-average estimate with sufficiently low variance. 
\end{itemize}

\section{Criteria for Ergodicity}
To demonstrate WS ergodicity we must show that as \(T \rightarrow \infty\):

\begin{itemize}
    \item \(E[\hat{\mu}_\mathrm{x}] \rightarrow \mu_\mathrm{x} \)
    \item \(Var[\hat{\mu}_\mathrm{x}] \rightarrow 0\)
     \item \(E[R_\mathrm{xx}](\tau) \rightarrow \mu_\mathrm{x} \)
    \item \(Var[R_\mathrm{xx}(\tau)] \rightarrow 0\)
\end{itemize}

\textbf{Theorem 6.1}: if $X(t)$ is a Gaussian, WSS process with $\mu_\mathrm{x}$ and $C_\mathrm{xx}(\tau)$, a sufficient condition for \(X(t)\) to be WS ergodic is:

\begin{enumerate}
    \item \( \int_{-\infty}^\infty C_\mathrm{xx}(\tau) \, d \tau < \infty \)
    \item \( \int_{-\infty}^\infty \ | \tau | C_\mathrm{xx}(\tau) \, d \tau < \infty \)
        \item \( \int_{-\infty}^\infty C^2_\mathrm{xx}(\tau) \, d \tau < \infty \)
    \item \( \int_{-\infty}^\infty \ | \tau | C^2_\mathrm{xx}(\tau) \, d \tau < \infty \)
\end{enumerate}

\section{Estimation Errors}

As with ensemble estimates, estimation errors exist in time-average estimates. Consider the estimate of a true parameter \(\theta\) given by  \(\hat{\theta}\). The relevant estimation errors are:

\begin{enumerate}
    \item \textbf{Bias}.  \(\text{b}[\hat{\theta}] = E[\hat{\theta}] - \theta\). 
    \begin{itemize}
        \item If \(\text{b}[\hat{\theta}] = 0\) then the estimator is said to be \emph{unbiased}.
    \end{itemize}
    

    \item \textbf{Variance}.  \(\text{Var}[\hat{\theta}] = E[(\hat{\theta} - E[\hat{\theta}])^2]\).

    \item \textbf{Mean square error}. \(\text{MSE}[\hat{\theta}] = E[(\hat{\theta} - \theta)^2]\).
    \begin{itemize}
        \item  An estimator is said to be \emph{consistent} if $\text{MSE} \rightarrow$ as \(T\rightarrow \infty\).
        \item This can also be expressed as  \(\text{MSE}(\hat{\theta}) = \text{Var}[\hat{\theta}] + \text{b}^2[\hat{\theta}]\) where $\text{Var}[\hat{\theta}]$ is the AC component of the MSE and \(\text{b}^2[\hat{\theta}]\) is the DC component.
    \end{itemize}
   
    \item \textbf{Normalized standard error}. (for \(\text{b}(\hat{\theta}) = 0)\)
    \[
    \varepsilon = \frac{\sqrt{\text{Var}[\hat{\theta}]}}{\theta}
    \]

    This is often expressed as a percentage.
\end{enumerate}

\subsection{Estimate of the Mean, $\hat{\mu}_\mathrm{x}$}

\[
\hat{\mu}_\mathrm{x} = \frac{1}{T}\int_{0}^T
X(t) \, dt
\]
\subsubsection*{Bias error}

\[
\text{b}[\hat{\mu}_\mathrm{x}] = E[\hat{\mu}_\mathrm{x}] -  \mu_\mathrm{x}
\]
\[
\begin{aligned}
    E[\hat{\mu}_\mathrm{x}]  & = E \left [ \frac{1}{T} \int_0^T X(t) \, dt \right ] \\
    & =  \frac{1}{T} \int_0^T E \left [ X(t) \right ] \, dt \\ 
    & = \mu_\mathrm{x}
\end{aligned}
\]

Therefore
\[
\text{b}[\hat{\mu}_\mathrm{x}] = 0
\]

and the estimator is \emph{unbiased}.

\subsubsection{Variance of the estimate of the mean}
\[
\begin{aligned}
 \text{Var}[\hat{\mu}_\mathrm{x}] = & E \left [ (\hat{\mu}_\mathrm{x} - \mu_\mathrm{x})^2 \right ] \\
        & = E [ \hat{\mu}^2_\mathrm{x}] - E^2[\hat{\mu}_\mathrm{x}] \\
        & = E^2[\hat{\mu}_\mathrm{x}] - {\mu}^2_\mathrm{x}\\
        & = E \left [ \frac{1}{T} \int_0^T x(u) \, du \cdot \frac{1}{T} \int_0^T x(\nu) \, d \nu \right ] - {\mu}^2_\mathrm{x}\\
        & = \frac{1}{T^2}  \int_0^T  \int_0^T E \left [  x(u) x(\nu) \right ] \, du \, d \nu   - {\mu}^2_\mathrm{x}\\
        & = \frac{1}{T^2}  \int_0^T  \int_0^T \phi_\mathrm{xx}(\nu-u) \, du \, d \nu   - {\mu}^2_\mathrm{x}
\end{aligned}
\]

Since \[
\mathrm{C}_\mathrm{xx}(\nu-u) = \mathrm{C}_\mathrm{xx}(\nu-u) - \mu^2_\mathrm{x}
\]

Then,

\[
\text{Var}[\hat{\mu}_\mathrm{x}] = \frac{1}{T^2}  \int_0^T  \int_0^T \mathrm{C}_\mathrm{xx}(\nu-u) \, du \, d \nu
\]

Let \( \tau = \nu - u\):
\[
\text{Var}[\hat{\mu}_\mathrm{x}] = \frac{1}{T^2}  \int_{-T}^T \mathrm{C}_\mathrm{xx} \left ( T - |\tau| \right ) \, d \tau  \quad \text{(Papoulis, p. 222)}
\]

Rearranging:
\[
\text{Var}[\hat{\mu}_\mathrm{x}] = \frac{1}{T}  \int_{-T}^T \mathrm{C}_\mathrm{xx} \left (1 - \frac{|\tau|}{T}  \right ) \, d \tau
\]

Now \emph{iff}:

\[
\int_{-\infty}^\infty \mathrm{C}_\mathrm{xx} (\tau) \, d \tau < \infty
\]

and

\[
\int_{-\infty}^\infty |\tau| \mathrm{C}_\mathrm{xx} (\tau) \, d \tau < \infty
\]

then $\text{Var}[\hat{\mu}_\mathrm{x}] \rightarrow 0$ as \(T \rightarrow \infty\), and the estimator is \emph{consistent}. These two results correspond to criteria 1 and 2 in Theorem 6.1.\\

\emph{Note:} for \(T\) large (compared to values of \(\tau\) over which $\mathrm{C}_\mathrm{xx} (\tau)$ is significant), then:
\[
\boxed{\text{Var}[\hat{\mu}_\mathrm{x}] \approx \frac{1}{T}  \int_{-T}^T \mathrm{C}_\mathrm{xx}  \, d \tau}
\]

The normalized standard error is:
\[
\varepsilon \% = \frac{\sqrt{\mathrm{Var}[\hat{\mu}_\mathrm{x}]}}{\mu_\mathrm{x}} \times 100\%
\]
This does not reduce to any convenient form. 
\begin{tcolorbox}
    \begin{exmp}{The Estimate the Mean for Bandlimited White Noise.}
    
    Consider The autocorrelation function and the power spectral density of pure bandlimited white noise, as illustrated in Figure \ref{fig:BLWN}. Find the analytic form of the variance of the estimate of the mean.
    \end{exmp}
\end{tcolorbox}
\begin{figure}[h!]
\centering
\includegraphics[scale=0.6]{figures/BLWN_t.png}
\includegraphics[scale=0.6]{figures/BLWN_f.png}
\caption{The autocorrelation function (top) and the power spectral density (bottom) of pure bandlimited white noise.}
\label{fig:BLWN}
\end{figure}

\[
\phi_\mathrm{xx}(\tau) = B N_o \frac{\sin(2 \pi B \tau)}{2 \pi B \tau} + \mu^2_\mathrm{x} = B N_o \, \mathrm{sinc}( 2 B N_o) + \mu^2_\mathrm{x}
\]

\[
\mathrm{C}_\mathrm{xx}(\tau) =  B N_o \frac{\sin(2 \pi B \tau)}{2 \pi B \tau} 
\]

For the estimator
\[
\hat{\mu}_\mathrm{x} = \frac{1}{T}\int_0^T x(t) \, dt \quad \rightarrow \quad  \text{choose } T>> \frac{1}{2B}
\]

Therefore,
\[
\begin{aligned}
    \text{Var}[\hat{\mu}_\mathrm{x}] & \approx \frac{1}{T}  \int_{-\infty}^\infty \mathrm{C}_\mathrm{xx}  \, d \tau \\
        & \approx \frac{1}{T}  \int_{-\infty}^\infty B N_o \frac{\sin(2 \pi B \tau)}{2 \pi B \tau}  \, d \tau \\
        & \approx \frac{\sigma^2_\mathrm{x}}{2B} \, \mathrm{V^2}
\end{aligned}
\]

Therefore, 
$\text{Var}[\hat{\mu}_\mathrm{x}] \rightarrow 0$ as \(T \rightarrow \infty\), and the estimator is \emph{consistent}.\\

Consider the standard error for this estimator:
\[
\varepsilon \% = \frac{\sqrt{\frac{\sigma^2_\mathrm{x}}{2 B T}}}{\mu_\mathrm{x}} \times 100\%
\]

\subsection{Estimate of the Variance, \(\hat{\sigma^2_\mathrm{x}}\)}


First, assume that \(\mu_\mathrm{x} = 0.\) Consider the estimator:
\[
\hat{\sigma}^2_\mathrm{x} = \frac{1}{T}\int_0^T x^2(t) \, dt
\]

\subsubsection*{Bias Error}
First, investigate the bias of the estimator:
\[
\mathrm{b}[\hat{\sigma}^2_\mathrm{x}] = E[\hat{\sigma}^2_\mathrm{x}] - \sigma^2_\mathrm{x}
\]

\[
E[\hat{\sigma}^2_\mathrm{x}] = E \left [  \frac{1}{T} \int_0^T x^2(t) \, dt \right ] = \frac{1}{T} \int_0^T E [x^2(t)] \, dt = \sigma^2_\mathrm{x}
\]

Therefore, the estimate of the variance is unbiased.\\

\subsubsection{Variance of the estimate of the variance}
Now consider the variance of \(\hat{\sigma}^2_\mathrm{x}\):

\[
\begin{aligned}
    \mathrm{Var}[\hat{\sigma}^2_\mathrm{x}] & = E [ (\hat{\sigma}^2_\mathrm{x} - {\sigma}^2_\mathrm{x})^2] \\
    & = E [\hat{\sigma}^4_\mathrm{x}] - 2\hat{\sigma}^4_\mathrm{x} + {\sigma}^4_\mathrm{x} \\
    & = E [\hat{\sigma}^4_\mathrm{x}] - {\sigma}^4_\mathrm{x} \\
    & = E \left [  \frac{1}{T} \int_0^T x^2(u) \, du \cdot \frac{1}{T} \int_0^T x^2(\nu) \, d \nu\right ] -{\sigma}^4_\mathrm{x} \\
    & = \frac{1}{T^2} \int_0^T \int_0^T E[x^2(u) x^2(\nu) ] \, du \,d \nu - {\sigma}^4_\mathrm{x} 
\end{aligned}
\]
To go any further, we must assume some distribution for \(X(t)\). Assume it is Gaussian with \(\mu_\mathrm{x} = 0\):
\[
\begin{aligned}
    E[x^2(u) x^2(\nu)] & = 2 \phi^2_\mathrm{xx}(u-\nu) + \sigma^4_\mathrm{x} \quad \text{from Piersol, p. 94}\\
    & = 2 \mathrm{C}^2_\mathrm{xx}(u-\nu) + \sigma^4_\mathrm{x} 
\end{aligned}
\]

\[
\begin{aligned}
    \mathrm{Var}(\hat{\sigma}^2_\mathrm{x}) & = \frac{1}{T^2} \int_0^T \int_0^T 2 \mathrm{C}^2_\mathrm{xx}(u-\nu) + \sigma^4_\mathrm{x}  \, du \, d \nu - \sigma^4_\mathrm{x} \\
    & = \frac{1}{T^2} \int_0^T \int_0^T 2 \mathrm{C}^2_\mathrm{xx}(u-\nu) \, du \, d \nu
\end{aligned}
\]

Let \(\tau = u - \nu\):
\[
\mathrm{Var}[\hat{\sigma}^2_\mathrm{x}] = \frac{2}{T}\int_{-T}^T \left ( 1 - \frac{|\tau|}{T}\right ) \mathrm{C}^2_\mathrm{xx}(\tau) \, d \tau
\]

Now \emph{iff}:

\[
\int_{-\infty}^\infty \mathrm{C}^2_\mathrm{xx} (\tau) \, d \tau < \infty
\]

and

\[
\int_{-\infty}^\infty \mathrm{C}^2_\mathrm{xx} |\tau|(\tau) \, d \tau < \infty
\]

then \(\text{Var}[\hat{\sigma}^2_\mathrm{x}] \rightarrow 0\) as \(T \rightarrow \infty\), and the estimator is \emph{consistent}. These two results correspond to criteria 3 and 4 in Theorem 6.1.\\

\emph{Note:} for \(T\) large (compared to values of \(\tau\) over which $\mathrm{C}_\mathrm{xx} (\tau)$ is significant), then: 

\[
\boxed{\mathrm{Var} [\hat{\sigma}^2_\mathrm{x}] \approx \frac{2}{T}\int_{-\infty}^\infty \mathrm{C}^2_\mathrm{xx}(\tau) \, d \tau}
\]

The normalized standard error is:
\[
\varepsilon \% = \frac{\sqrt{\mathrm{Var}[\hat{\sigma}^2_\mathrm{x}]}}{\sigma^2_\mathrm{x}} \times 100\%
\]
which does not reduce to any convenient form.

\begin{tcolorbox}
    \begin{exmp}{Continue the bandlimited white noise example for the estimate of the variance.}
   
    \end{exmp}
\end{tcolorbox}

For our process \(X(t)\):
\[
\mathrm{C}_\mathrm{xx}(\tau) =  B N_o \frac{\sin(2 \pi B \tau)}{2 \pi B \tau} 
\]

Note that the magnitude of the peak of the autocorrelation function is \(N_0 B\), which is also equal to the MSV ( and since \(\mu_\mathrm{x} = 0\), is equal to the variance \(\sigma^2_\mathrm{x}\).

Therefore:
\[
N_o B = \sigma^2_\mathrm{x}
\]

So 
\[
\mathrm{C}_\mathrm{xx}(\tau) =  \sigma^2_\mathrm{x} \frac{\sin(2 \pi B \tau)}{2 \pi B \tau} 
\]

and the magnitude of \(\Phi_\mathrm{xx} (\omega) = \frac{\sigma^2_\mathrm{x} }{2 B}\).\\

Using
\[
{\mathrm{Var} [\hat{\sigma}^2_\mathrm{x}] \approx \frac{2}{T}\int_{-\infty}^\infty \mathrm{C}^2_\mathrm{xx}(\tau) \, d \tau}
\]

We have
\[
\begin{aligned}
    \text{Var}[\hat{\sigma}^2_\mathrm{x}] & \approx \frac{2}{T}  \int_{-\infty}^\infty \mathrm{C}^2_\mathrm{xx}  \, d \tau \\
        & \approx \frac{2}{T}  \int_{-\infty}^\infty \sigma^4_\mathrm{x} \left ( \frac{\sin(2 \pi B \tau)}{2 \pi B \tau} \right ) ^2 \, d \tau \\
        & \approx \frac{\sigma^4_\mathrm{x}}{BT} 
\end{aligned}
\]

Consider the standard error for this estimator:
\[
\begin{aligned}
    \varepsilon \% & = \frac{\sqrt{\frac{\sigma^4_\mathrm{x}}{B T}}}{\sigma^2_\mathrm{x}} \times 100\% \\
        & = \frac{1}{\sqrt{BT}} \times 100\% 
\end{aligned}
\
\]

If \(T\) is a parameter over which we have control, choose 
\[
T = \frac{1}{B \varepsilon^2}
\]

to achieve a desired error performance.

\subsection{Estimate of the Autocorrelation Function, \(\hat{\phi}_\mathrm{xx} (\tau)\)}


Consider the estimator:
\[
\mathrm{R}_\mathrm{xx} (\tau) = \hat{\phi}_\mathrm{xx} (\tau) = \frac{1}{T}\int_0^T x(t) x(t + \tau)\, dt
\]

\subsubsection*{Bias Error}

\[
E \left [ \hat{\phi}_\mathrm{xx} (\tau) \right ] = \frac{1}{T} \int_0^T x(t) x(t+\tau) \, dt = \phi_\mathrm{xx} (\tau)
\]
Therefore, the estimator is unbiased.

\subsubsection*{Variance of the estimate \( \hat{\phi}_\mathrm{xx} (\tau) \)}
\[
\begin{aligned}
    \text{Var}[\hat{\phi}_\mathrm{xx} (\tau)] & = E \left [ \hat{\phi}^2_\mathrm{xx} (\tau) \right ] - \phi^2_\mathrm{xx} (\tau) \quad \text{since } E\left [ \hat{\phi}_\mathrm{xx} (\tau) \right ] = \phi_\mathrm{xx}  (\tau) \\
    & = E \left [ \frac{1}{T} \int_0^T x(u) x(u+\tau) \, du \cdot \frac{1}{T} \int_0^T x(\nu) x(\nu+\tau) \, d \nu\right ] - \phi^2_\mathrm{xx} (\tau) \\
    & = \frac{1}{T^2} \int_0^T \int_0^T x(u) x(u+\tau) x(\nu) x(\nu+\tau) \, du \, d \nu - \phi^2_\mathrm{xx} (\tau)
\end{aligned}
\]

To go any further, we must assume some distribution for \(X(t)\). Assume it is Gaussian with \(\mu_\mathrm{x} = 0\):

\[
    \text{Var}[\hat{\phi}_\mathrm{xx} (\tau)] = \frac{1}{T^2} \int_0^T \int_0^T \left \{ \phi^2_\mathrm{xx} (\tau) + \phi^2_\mathrm{xx} (u - \nu) + \phi_\mathrm{xx} (u - \nu + \tau) \phi_\mathrm{xx} (u - \nu - \tau) \right \}   \, du \, d \nu  - \phi^2_\mathrm{xx} (\tau)
\]

Let \(\gamma = u - \nu\):

\[
    \text{Var}[\hat{\phi}_\mathrm{xx} (\tau)] = \frac{1}{T} \int_{-T}^T \left \{ \phi^2_\mathrm{xx} (\gamma) + \phi_\mathrm{xx} (\gamma + \tau) \phi_\mathrm{xx} (\gamma - \tau) \right \} \left \{ 1 - \frac{|\gamma|}{T}\right \} \, d \gamma
\]

Note that as \(T \rightarrow \infty, \text{Var}[\hat{\phi}_\mathrm{xx} (\tau)] \rightarrow 0\) \emph{iff} the integrals are finite, in which case the estimator is consistent.\\

For \(T\) large (with respect to the values of \(\phi_\mathrm{xx} (\tau) \) that are significant): 

\[
    \text{Var}[\hat{\phi}_\mathrm{xx} (\tau)] \approx \frac{1}{T} \int_{-\infty}^\infty  \phi^2_\mathrm{xx} (\gamma) + \phi_\mathrm{xx} (\gamma + \tau) \phi_\mathrm{xx} (\gamma - \tau)  \, d \gamma
\]

\begin{tcolorbox}
    \begin{exmp}{Continue the bandlimited white noise example for the estimate of the autocorrelation function.}
   
    \end{exmp}
\end{tcolorbox}

\[
\begin{aligned}
    \text{Var}[\hat{\phi}_\mathrm{xx} (\tau)]  & = \left (\frac{\sigma^4_\mathrm{x}}{2 B T} \right )\left (  1 + \frac{\sin^2 2 \pi B \tau}{(2 \pi B \tau)^2}\right ) \\
    & = 
    \begin{cases}
        \frac{\sigma^4_\mathrm{x}}{B T} & , \tau = 0 \\
        \frac{\sigma^4_\mathrm{x}}{2 B T} & , \tau \rightarrow \infty 
    \end{cases}
\end{aligned}
\]

Note that at \( \tau = 0\) the expression for \(\text{Var}[\hat{\phi}_\mathrm{xx} (\tau)] \) is that of \(\text{Var}[\hat{\sigma}_\mathrm{x}] \) and at \( \tau \rightarrow \infty\) it is that of \( \text{Var}[\hat{\mu}^2_\mathrm{x} ] \).

\begin{figure}[h!]
\centering
\includegraphics[scale=0.6]{figures/Var_ACF.png}
\caption{The variance of $\hat{\phi}_\mathrm{xx} (\tau)$ as a function of \(\tau\). Here, two examples are shown, both with  $B = 2 Hz$ and $T = 10 s$, but with $\sigma_\mathrm{x} = 1.0 $ and $\sigma_\mathrm{x} = 1.1$}
\label{fig:Var_ACF}
\end{figure}

The normalized standard error is:

\[
\begin{aligned}
\varepsilon \% & = \frac{\sqrt{\mathrm{Var}[\hat{\phi}_\mathrm{xx}]}}{{\phi}_\mathrm{xx}} \\
    & = \frac{1}{2 B T} \left ( 1 + \frac{\phi^2_\mathrm{xx}(0)}{{\phi^2}_\mathrm{xx}(\tau)} \right )
\end{aligned}
\]

which equals \(\frac{1}{BT}\) at \(\tau=0\) and goes to \(\infty\) as \(\tau \rightarrow \infty\).

\section{Power Density Spectrum Estimation}

We will consider the implementation and performance of a few types of PSD estimators. 

\subsection{Analog PSD Estimator}

Analog estimation of the PSD may be considered as a sliding narrowband bandpass filter of width \(B_e\), providing estimates of the power density at each centre frequency. The filter width \(B_e\) can be interpreted as the resolutio of teh PSD estimator. At each frequency, we can consider the amplitude of the PSD as an estimate of the MSV of the signal, as illustrated in Figure \ref{fig:PSD_analog}.

\begin{figure}[h!]
\centering
\includegraphics[scale=0.6]{figures/PSD_analog.png}
\caption{An illustration of the stages of an analog PSD estimator.}
\label{fig:PSD_analog}
\end{figure}

\[
\hat{\mathrm{MSV}}(f,B_e) = \frac{1}{T}\int_0^T x^2(t,f,B_e) \, dt
\]

\subsubsection*{Estimate of the Mean of \( \hat{\Phi}_\mathrm{xx} (f) \)}
To generate a power density spectrum estimate, we divide the estimate of the MSV by the area under the narrowband filter, \(2 B_e\) (accounting for both sides of the PSD):

\[
\begin{aligned}
    \hat{\Phi}_\mathrm{xx} (f) & = \frac{1}{2 B_e T} \int_0^T x^2(t,f,B_e) \, dt \\
    & = \frac{\hat{\mathrm{MSV}}(f,B_e)}{2 B_e}
\end{aligned}
\]

We also know that 
\[
\mathrm{MSV}(f,B_e) = 2 \int_{f-\frac{B_e}{2}}^{f+\frac{B_e}{2}} \Phi_\mathrm{xx} (f) \, df
\]

Consider now the expected value of the PSD estimate:

\[
\begin{aligned}
 E \left [ \hat{\Phi}_\mathrm{xx} (f) \right ] & = \frac{1}{2 B_e} \underbrace{\frac{1}{T}\int_0^T E [x^2(t,f,B_e)] \, dt}_{\mathrm{MSV}(f,B_e)} \\   
 & = \frac{1}{2 B_e}  \int_{f-\frac{B_e}{2}}^{f+\frac{B_e}{2}} \Phi_\mathrm{xx} (f) \, df
\end{aligned}
\]

We can see that in general, 
\[
E \left [ \hat{\Phi}_\mathrm{xx} (f) \right ] \ne \Phi_\mathrm{xx} (f)
\]

unless 
\[
\frac{1}{2 B_e}  \int_{f-\frac{B_e}{2}}^{f+\frac{B_e}{2}} \Phi_\mathrm{xx} (f) \, df \, = \, \frac{\hat{\mathrm{MSV}}(f,B_e)}{2 B_e}
\]

That is, the area under \(\Phi_\mathrm{xx} (f)\) between \(f \pm \frac{B_e}{2}\) is equal to the rectangle of height \(\hat{\mathrm{MSV}}(f,B_e)\). This is is true if \(\Phi_\mathrm{xx} (f)\) is relatively constant over \(f \pm \frac{B_e}{2}\). The larger the value of \(B_e\), the less likely this is true.

\begin{figure}[h!]
\centering
\includegraphics[scale=0.4]{figures/PSD_analog_bias.png}
\caption{The bias inherent in an analog PSD estimate. Small filter bandwidth (\(Be_1\)) introduces smaller bias than larger filter bandwidth (\(Be_2\)).}
\label{fig:PSD_analog_bias}
\end{figure}


\subsubsection*{Estimate of the Variance of \( \hat{\Phi}_\mathrm{xx} (f) \)}

We will neglecting the bias error to simplify the derivation:


\[
\begin{aligned}
    \hat{\Phi}_\mathrm{xx} (f) & = \frac{\hat{\mathrm{MSV}}(f,B_e)}{2 B_e}
\end{aligned}
\]

\[
\begin{aligned}
    \mathrm{Var}[\hat{\Phi}_\mathrm{xx} (f)] & = \frac{\mathrm{Var}[\hat{\mathrm{MSV}}(f,B_e)]}{4 B_e^2}
\end{aligned}
\]

\[
\begin{aligned}
    \mathrm{Var}[\hat{\mathrm{MSV}}(f,B_e)] & =  \frac{\sigma^4_\mathrm{x}(f,B_e)}{B_e T} \quad \text{from the proof of } \mathrm{Var}[ \hat{\sigma^2_\mathrm{x}}]\\
\end{aligned}
\]

Therefore,
\[
\begin{aligned}
    \mathrm{Var}[\hat{\Phi}_\mathrm{xx} (f)] & = \frac{\sigma^4_\mathrm{x}(f,B_e)}{4 B_e^3 T} \\
\end{aligned}
\]

Since \(\sigma^2(f,B_e) = 2 B_e \Phi_\mathrm{xx} (f)\):
\[
\boxed{
    \mathrm{Var}[\hat{\Phi}_\mathrm{xx} (f)] = \frac{\Phi^2_\mathrm{xx} (f)}{B_e T} \\
}
\]

and the normalized standard error is:
\[
\varepsilon = \frac{1}{\sqrt{B_e T}}
\]

\subsection{PSD Estimator using the Autocorrelation Function}

We know that 
\[
{\Phi}_\mathrm{xx} (f) = \mathcal{F} \left \{{\phi}_\mathrm{xx} (f)\right \}
\]

Estimating the autocorrelation function requires using a finite data window. Let's assume that the autocorrelation function is truncated at \(\tau = [-\tau_m,\tau_m]\). This is equivalent to multiplying the true autocorrelation function by a rectangular window \(p(\tau) = \mathrm{rect}[-\tau, \tau]\):
\[
\phi_\mathrm{xx_T} (\tau) = \phi_\mathrm{xx_T} (\tau) \cdot p(\tau)
\]
and 

\[
\Phi_\mathrm{xx_T} (f) = \Phi_\mathrm{xx_T} (f) \ast  P(f)
\]

The frequency spectrum \(P(f)\) of the truncation window takes the form a a sinc function with zero-crossing of the main lobe at \(\frac{1}{2 \tau_m}\), as illustrated in Figure \ref{fig:P(f)}.

\begin{figure}[h!]
\centering
\includegraphics[scale=0.6]{figures/P(f).png}
\caption{The frequency spectrum of the autocorrelation function truncation window.}
\label{fig:P(f)}
\end{figure}

Due to the fact that the spectrum \(\phi_\mathrm{xx_T} (\tau)\) is convolved with this since function, \(P(f)\) effectively "smears" the true spectrum \(\phi_\mathrm{xx} (\tau)\). Also, it introduces a spectral resolution roughly equal to the wide of the main lobe:
\[
B_e \approx \frac{1}{\tau_m}
\]

If we use all of the data toward the estimatio of the autocorrelation function, then \(t_m = T\) and 
\[
B_e \approx \frac{1}{T}
\]

A spectral estimate \(\hat{\Phi}_\mathrm{xx_T} (f) \) taken from an estimate of the autocorrelation function \(\hat{\phi}_\mathrm{xx_T} (\tau) \) will be subject to this smearing and resolution limitation.\\

\subsection{FFT Estimate}
The FFT operates on a finite data sequence \(x_T(t)\) on \([0,T]\) as illustrated in Figure \ref{fig:FFT_rand}.


\begin{figure}[h!]
\centering
\includegraphics[scale=0.6]{figures/FFT_rand.png}
\caption{The data window \(x_T(t)\) used in the FFT estimate of the PSD.}
\label{fig:FFT_rand}
\end{figure}

The PSD estimate is:

\[
\hat{\Phi}_\mathrm{xx_T} (f) = \frac{|X_T(f)|^2}{T}
\]

where \[
X_T(f) = \mathrm{FFT}[x_{_T}(t)]
\]

Consider now the expected value of the PSD estimate:
\[
\begin{aligned}
    E \left [ \hat{\Phi}_\mathrm{xx_T} (f) \right ] & =\frac{1}{T} E \left [ |X_T(f)|^2 \right ]\\
        & = \frac{1}{T} E \left [ \left |\int_0^T x_{_T}(t) e^{-j 2 \pi f t} \, dt \right |^2 \right ]\\
        & = \frac{1}{T} \int_0^T \int_0^T E \left [ x_{_T}(u) x_{_T}(\nu) \right ] e^{-j 2 \pi f u} e^{-j 2 \pi f \nu} \, du \, d \nu  \\
        & = \frac{1}{T} \int_0^T \int_0^T  {\phi}_\mathrm{xx} (\nu - u) e^{-j 2 \pi f u} e^{-j 2 \pi f \nu} \, du \, d \nu  \\
        & = \int_{-T}^T  {\phi}_\mathrm{xx} (\tau) e^{-j 2 \pi f \tau} \left [ 1 - \frac{|\tau|}{T}\right ] d \tau
\end{aligned}
\]
Thus, in general:
\[
E \left [ \hat{\Phi}_\mathrm{xx_T} (f) \right ] \ne \Phi_\mathrm{xx} (f)
\]

However, for \(T\) large with respect to values of \( \tau \) for which \(\phi_\mathrm{xx}\) is significant:
\[
E \left [ \hat{\Phi}_\mathrm{xx_T} (f) \right ] \approx \int_{-\infty}^\infty  {\phi}_\mathrm{xx} (\tau) e^{-j 2 \pi f \tau} \, d \tau = \Phi_\mathrm{xx} (f)
\]

We assert the general form of the normalized standard error:
\[
\varepsilon  = \frac{1}{\sqrt{B_e T}}
\]



Considering now that the data window can be considered as 
\[
x_{_T} = x(t) \times p(t)
\]
where \(p(t) = \mathrm{rect[0,T]}\), then the spectrum of \(p(t)\) is a sinc function of width \(\frac{2}{T}\) as depicted in Figure \ref{fig:FFT_P(f)}.

\begin{figure}[h!]
\centering
\includegraphics[scale=0.6]{figures/FFT_P(f).png}
\caption{The frequency spectrum of the data truncation window.}
\label{fig:FFT_P(f)}
\end{figure}

This suggests a frequency resolution \(B_e \approx \frac{2}{T}\). 

Therefore,
\[
\varepsilon  = \frac{1}{\sqrt{B_e T}} = \frac{1}{\sqrt{2}}
\]

This is an interesting result, as the normalized standard error is fixed, \emph{independent of the data length}, \(T\). \\

To reduce the variance of the FFT PSD estimate, the general approach is to choose a relatively large \(T\), and divide it into segments. The PSD estimates of the segments are then ensemble averaged to provide a lower-variance PSD estimate. This is referred to as the Welch method. There are two drawbacks to this approach:
\begin{itemize}
    \item The resolution of the PSD estimate of each segemnt is \(B_e = \frac{2}{T_{seg}}\). If the data record \(T\) is divided into \(M\) segments, then \(T_{seg} = \frac{T}{M} \) and the effective frequency resolution will be degraded to \(M B_e\).
    \item Choosing a larger \(T\) may increase the likelihood that the data used violate the assumption of stationarity.
\end{itemize}






\end{document}
